% Generated by scripts/build_reports.py; edit reviewed Markdown inputs.
% Compile with XeLaTeX.
% Options for packages loaded elsewhere
\PassOptionsToPackage{unicode}{hyperref}
\PassOptionsToPackage{hyphens}{url}
\PassOptionsToPackage{dvipsnames,svgnames,x11names}{xcolor}
\documentclass[
  11pt,
]{article}
\usepackage{xcolor}
\usepackage[margin=1in]{geometry}
\usepackage{amsmath,amssymb}
\setcounter{secnumdepth}{-\maxdimen} % remove section numbering
\usepackage{iftex}
\ifPDFTeX
  \usepackage[T1]{fontenc}
  \usepackage[utf8]{inputenc}
  \usepackage{textcomp} % provide euro and other symbols
\else % if luatex or xetex
  \usepackage{unicode-math} % this also loads fontspec
  \defaultfontfeatures{Scale=MatchLowercase}
  \defaultfontfeatures[\rmfamily]{Ligatures=TeX,Scale=1}
\fi
\usepackage{lmodern}
\ifPDFTeX\else
  % xetex/luatex font selection
\fi
% Use upquote if available, for straight quotes in verbatim environments
\IfFileExists{upquote.sty}{\usepackage{upquote}}{}
\IfFileExists{microtype.sty}{% use microtype if available
  \usepackage[]{microtype}
  \UseMicrotypeSet[protrusion]{basicmath} % disable protrusion for tt fonts
}{}
\makeatletter
\@ifundefined{KOMAClassName}{% if non-KOMA class
  \IfFileExists{parskip.sty}{%
    \usepackage{parskip}
  }{% else
    \setlength{\parindent}{0pt}
    \setlength{\parskip}{6pt plus 2pt minus 1pt}}
}{% if KOMA class
  \KOMAoptions{parskip=half}}
\makeatother
\usepackage{longtable,booktabs,array}
\usepackage{caption}
\captionsetup[table]{skip=6pt}
\newcounter{none} % for unnumbered tables
\usepackage{calc} % for calculating minipage widths
% Correct order of tables after \paragraph or \subparagraph
\usepackage{etoolbox}
\makeatletter
\patchcmd\longtable{\par}{\if@noskipsec\mbox{}\fi\par}{}{}
\makeatother
% Allow footnotes in longtable head/foot
\IfFileExists{footnotehyper.sty}{\usepackage{footnotehyper}}{\usepackage{footnote}}
\makesavenoteenv{longtable}
\setlength{\emergencystretch}{3em} % prevent overfull lines
\providecommand{\tightlist}{%
  \setlength{\itemsep}{0pt}\setlength{\parskip}{0pt}}
\usepackage{amsmath,amssymb}
\usepackage{xurl}
\usepackage{fancyhdr}
\pagestyle{fancy}
\fancyhf{}
\fancyhead[L]{\small Pure Thought | Revised research specifications}
\fancyfoot[L]{\small 2026-09-17}
\fancyfoot[R]{\small\thepage}
\setlength{\headheight}{15pt}
\setlength{\emergencystretch}{3em}
\allowdisplaybreaks
\usepackage{bookmark}
\IfFileExists{xurl.sty}{\usepackage{xurl}}{} % add URL line breaks if available
\urlstyle{same}
% fallback for those not using the hyperref driver hyperxmp:
\makeatletter
\@ifundefined{xmpquote}{\newcommand{\xmpquote}[1]{#1}}{}
\makeatother
\hypersetup{
  pdftitle={Complete problem summaries and mathematical formulations},
  pdfauthor={Pure Thought research portfolio},
  colorlinks=true,
  linkcolor={blue},
  filecolor={Maroon},
  citecolor={Blue},
  urlcolor={blue},
  pdfcreator={LaTeX via pandoc}}

\title{Complete problem summaries and mathematical formulations}
\author{Pure Thought research portfolio}
\date{Revised 17 September 2026}

\begin{document}
\maketitle
\thispagestyle{empty}
\clearpage

{
\setcounter{tocdepth}{1}
\tableofcontents
\clearpage
}
\section{Complete problem summaries and mathematical formulations}

Updated 2026-09-17. All 30 canonical problems and all 12 supplemental
topics are included in full, with plain-language guides, assumptions,
mathematical targets, outputs and validation gates. The displayed
formulations make the report equations and claim directions explicit.
These are proposed research tasks, not completed research results.

\href{https://github.com/rick-stevens-ai/Pure-Thought/blob/main/reviews/COMPLETE-REVIEW.md}{Review
and critiques} \ensuremath{\cdot}
\href{https://github.com/rick-stevens-ai/Pure-Thought/blob/main/reviews/EVIDENCE-STANDARD.md}{Evidence
standard} \ensuremath{\cdot}
\href{https://github.com/rick-stevens-ai/Pure-Thought/blob/main/reviews/SYNTHESIS-CROSSWALK.md}{Synthesis
crosswalk} \ensuremath{\cdot}
\href{https://github.com/rick-stevens-ai/Pure-Thought/blob/main/latex-reports/Complete-Problem-Summaries.pdf}{PDF
reading copy}

\begin{center}\rule{0.5\linewidth}{0.5pt}\end{center}

\section{01. Extremal holomorphic CFT partition functions: certified
necessary tests}

\textbf{Revised specification --- 2026-09-11.}
\href{https://github.com/rick-stevens-ai/Pure-Thought/blob/main/reviews/01-AdS3-Modular-Bootstrap.md}{Detailed
critique} \ensuremath{\cdot}
\href{https://github.com/rick-stevens-ai/Pure-Thought/blob/main/archive/original-PRDs/01-AdS3-Modular-Bootstrap.md}{Original
description} \ensuremath{\cdot}
\href{https://github.com/rick-stevens-ai/Pure-Thought/blob/main/reviews/EVIDENCE-STANDARD.md}{Evidence
standard}

\subsection{Plain-language guide}

\subsubsection{The problem in everyday terms}

Imagine trying to identify a musical instrument from a list of the notes
it can produce. Some lists are impossible because the notes do not fit
the instrument's mathematical rules. Others pass those checks, but that
still does not tell us how to build the instrument. Here the
``instrument'' is a quantum theory and the ``notes'' are its allowed
states. We ask whether exceptionally sparse lists of states obey the
consistency rules. Such sparse theories are candidates for describing a
simplified universe containing gravity with no additional matter.

\subsubsection{Key terms}

\begin{itemize}
\item
  \textbf{Conformal field theory (CFT)} --- A quantum theory whose laws
  respect transformations that preserve angles, including changes of
  scale; it describes systems with no preferred length scale.
\item
  \textbf{AdS\ensuremath{{}_{3}}} --- A three-dimensional spacetime with
  a particular negative curvature, used as a mathematically tractable
  setting for gravity.
\item
  \textbf{Holomorphic} --- Depending on one complex coordinate rather
  than both it and its complex conjugate; this is a special restriction
  on the theories studied here.
\item
  \textbf{Central charge c} --- A number characterizing aspects of a
  CFT's symmetry and number of degrees of freedom. Here c=24k, with k a
  positive integer.
\item
  \textbf{Partition function} --- A generating function that packages
  the theory's state energies or scaling weights and the number of
  states at each level.
\item
  \textbf{Primary and descendant} --- A primary is a basic type of
  excitation; descendants are related excitations generated from it by
  symmetry.
\item
  \textbf{Extremal spectrum} --- A spectrum with the largest prescribed
  initial gap before additional primary excitations, under this
  particular ansatz.
\item
  \textbf{Modular invariance} --- The requirement that different
  equivalent descriptions of a torus give the same partition function.
\end{itemize}

\subsubsection{Why this matters}

Quantum gravity is hard partly because we do not know which
mathematically plausible theories are internally consistent. These tests
can eliminate candidates without experiments or a complete construction.
They also connect gravity to the arithmetic of modular functions:
counting quantum states becomes a problem about exact number sequences.

\subsubsection{What progress would mean}

A useful first result is a trustworthy table of candidates and any exact
reason for rejecting one. Passing the table's tests would identify a
candidate worth further study; constructing the full theory would remain
a much stronger achievement.

\subsection{Core question}

For c=24k, which extremal holomorphic partition-function candidates pass
exact genus-one spectral tests, and can any candidate be excluded by an
additional, explicitly stated consistency condition?

\subsection{Scope and assumptions}

Start with bosonic holomorphic theories, a unique vacuum, c=24k, and
k=1,\ldots,4. This is a deliberately restricted version of the original
gravity question. Use h for holomorphic weight; do not identify it with
the full scaling dimension \ensuremath{\Delta}=h+\ensuremath{\bar{h}}.
The extremal ansatz has no nonvacuum Virasoro primary with
h\ensuremath{\le}k. Treat the nonchiral modular bootstrap as a separate
extension with independent (h,\ensuremath{\bar{h}}) data.

Keep the motivating existence question, but make the first deliverable a
necessary-condition test. A modular function with integer coefficients
is not a constructed CFT or a demonstrated gravity dual.

\subsection{Mathematical target}

With q=exp(2\ensuremath{\pi}i\ensuremath{\tau}), c\textgreater1, use

\[\chi _0=q^{-c/24}\prod_{n=2}^{\infty }(1-q^n)^{-1},\qquad \chi _h=q^{h-(c-1)/24}/\eta (\tau )\quad(h>0).\]

Construct the unique degree-k polynomial in J=j\ensuremath{-}744 whose
expansion agrees with \ensuremath{\chi}\ensuremath{{}_{0}} through
q\ensuremath{{}^{0}}. Write \(Z_k=\chi_0+\sum_{h\ge k+1}d_h\chi_h\) and
extract primary multiplicities exactly through a declared cutoff H,
initially H=100. Prove modular invariance from the polynomial
representation, not from evaluations at a few \ensuremath{\tau}. A
negative or noninteger d\ensuremath{{}_{h}} excludes the stated
candidate; positivity through H is only a finite test.

\subsection{Required outputs and proof obligations}

Export the polynomial coefficients, exact q-series, primary
multiplicities, conventions, cutoff, and a checker that independently
expands \ensuremath{\eta} and J. An exclusion must name the violated
necessary condition. An additional genus-two or OPE test must state its
own assumptions and how its constraints follow from CFT consistency.

\subsection{Validation and rejection controls}

Recover J=q\ensuremath{{}^{-1}}+196884q+… at k=1 and
d\ensuremath{{}_{2}}=196883 after subtracting the vacuum descendant.
Test the checker with a deliberately altered coefficient. Check
agreement of two series-construction routes. Do not report all-order
coefficient positivity without a tail theorem.

\subsection{Milestones and research extension}

Baseline: exact k=1,\ldots,4 tables with explicitly finite coverage.
Strong: one additional sewing/OPE condition with independently checked
consequences. Research extension: a new exclusion theorem or a full CFT
construction; neither is a required outcome or a promised deadline.

Milestones are acceptance gates, not calendar promises. Before a run,
freeze the concrete instance/domain, tolerance, compute budget, and
permitted dependencies. A bounded search that fails to find a result
must preserve an unresolved outcome.

\subsection{Starting instruction}

Work on the baseline above. First make every input and convention
explicit, then validate the reference case. Build the checker alongside
the calculation. Label numerical evidence, finite-model certificates,
and continuum theorems separately; return the strongest justified result
and identify the exact unmet proof obligation.

\subsection{Primary references and their role}

\begin{itemize}
\tightlist
\item
  \href{https://arxiv.org/abs/0706.3359}{Witten, Three-Dimensional
  Gravity Revisited} --- Motivation for the holomorphic extremal ansatz,
  not a proof that every candidate exists.
\end{itemize}

References support definitions or limitations; the scope choices and
proposed milestones are editorial recommendations, not claims that these
results have already been obtained.

\subsection{Displayed formulation and claim boundary}

\href{https://github.com/rick-stevens-ai/Pure-Thought/blob/main/PRDs/01-AdS3-Modular-Bootstrap.md}{Canonical
specification}.

For the holomorphic sector with \(c=24k\), \(q=e^{2\pi i\tau}\) and
\(J=j-744\),

\[\chi_0=q^{-k}\prod_{n=2}^{\infty}(1-q^n)^{-1},\qquad
\chi_h=\frac{q^{h-k+1/24}}{\eta(\tau)}\quad(h>0).\]

\[Z_k=P_k(J)=\chi_0+\sum_{h\ge k+1}d_h\chi_h,\qquad
P_k(J)-\chi_0=O(q).\]

Determine \(P_k\) exactly and test \(d_h\in\mathbb Z_{\ge0}\) through a
declared \(H\). At \(k=1\), \(d_2=196883\). Passing this finite
necessary test does not construct the CFT or establish OPE consistency.

\begin{center}\rule{0.5\linewidth}{0.5pt}\end{center}

\section{02. Conditional positivity bounds for a specified gravitational
amplitude}

\textbf{Revised specification --- 2026-09-11.}
\href{https://github.com/rick-stevens-ai/Pure-Thought/blob/main/reviews/02-Gravitational-Positivity-Bounds.md}{Detailed
critique} \ensuremath{\cdot}
\href{https://github.com/rick-stevens-ai/Pure-Thought/blob/main/archive/original-PRDs/02-Gravitational-Positivity-Bounds.md}{Original
description} \ensuremath{\cdot}
\href{https://github.com/rick-stevens-ai/Pure-Thought/blob/main/reviews/EVIDENCE-STANDARD.md}{Evidence
standard}

\subsection{Plain-language guide}

\subsubsection{The problem in everyday terms}

A low-energy theory is like a map drawn at a coarse scale: it describes
what we can resolve while summarizing finer details in a few adjustable
numbers. Not every choice of those numbers can come from a consistent
microscopic world. This problem asks whether basic rules about
probabilities, cause and effect, and scattering restrict one such number
when gravity is present. Gravity's long reach makes the usual argument
harder, so part of the task is making the assumptions precise.

\subsubsection{Key terms}

\begin{itemize}
\item
  \textbf{Effective field theory (EFT)} --- An approximate description
  valid over a stated range of energies, with unresolved physics
  represented by additional terms.
\item
  \textbf{Wilson coefficient} --- The number multiplying one of those
  additional terms; it measures the strength of a correction to the
  simpler theory.
\item
  \textbf{Scalar particle} --- A particle with zero intrinsic spin,
  chosen here to simplify the scattering calculation.
\item
  \textbf{Scattering amplitude} --- A quantum quantity used to calculate
  probabilities for incoming particles to emerge in particular outgoing
  states.
\item
  \textbf{Positivity bound} --- An inequality on theory parameters
  derived from a quantity that must be nonnegative under specified
  assumptions.
\item
  \textbf{Infrared (IR) and ultraviolet (UV)} --- The low-energy or
  long-distance regime, and the high-energy or short-distance regime,
  respectively.
\item
  \textbf{Pole and subtraction} --- A pole is a singular term in a
  mathematical expression; a subtraction removes a specified
  contribution so the remaining quantity can be analyzed.
\item
  \textbf{Dispersion relation} --- A mathematical relation connecting an
  amplitude's behavior in different energy regions.
\item
  \textbf{Regge assumption} --- An assumption about how scattering
  behaves at very high energy with momentum transfer held fixed.
\end{itemize}

\subsubsection{Why this matters}

These bounds can tell us that an apparently sensible low-energy model
cannot fit particular assumptions about its underlying physics. That is
valuable when direct access to microscopic gravity is unavailable. A
reliable bound also prevents researchers from interpreting an artifact
of a singular calculation as a physical restriction.

\subsubsection{What progress would mean}

The first success would be one defensible inequality with its
assumptions attached. It would narrow the parameter space of candidate
theories, without claiming that every surviving theory has a complete
microscopic realization.

\subsection{Core question}

Under a fully stated infrared prescription and high-energy assumption,
what certified inequality constrains one selected low-energy coefficient
of scalar scattering coupled to gravity?

\subsection{Scope and assumptions}

Use four-dimensional scattering of one identical real massive scalar
coupled to Einstein gravity. This replaces the redundant schematic
pure-gravity R\ensuremath{{}^{2}} coefficient list with an
amplitude-defined observable. Set m=1 as the mass unit; specify
\ensuremath{\kappa}, the EFT subtraction scale, perturbative order, and
whether the target is a tree coefficient or a renormalized coefficient.
Graviton helicity amplitudes are a later extension.

The first task is to establish a valid sum rule. Do not run an
optimization until the pole subtraction, massless cuts, crossed channel,
and large-contour contribution have been derived.

\subsection{Mathematical target}

Let \ensuremath{\nu}=s\ensuremath{-}2m\ensuremath{{}^{2}}+t/2, so
s\ensuremath{\leftrightarrow}u sends
\ensuremath{\nu}\ensuremath{\mapsto}\ensuremath{-}\ensuremath{\nu}.
Define B(\ensuremath{\nu},t) by subtracting the explicitly calculated
light poles and the chosen low-energy cut contributions from
A(\ensuremath{\nu},t). Define
b\ensuremath{{}_{2}}(t)=(1/2)\ensuremath{\partial}\ensuremath{{}^{2}}\_\ensuremath{\nu}B(0,t).

Derive b\ensuremath{{}_{2}}(t)=I\_UV(t)+R\_\ensuremath{\infty}(t) with
all kernels and subtractions written out. Choose either a finite-t
smeared positive sum rule or a justified regulated forward limit.
Nonnegativity must be established for the resulting partial-wave kernel;
Im A(s,t)\ensuremath{\ge}0 at arbitrary t\textless0 is not an admissible
substitute. State a bound on R\_\ensuremath{\infty} or retain it as an
explicit unknown.

\subsection{Required outputs and proof obligations}

Return a derivation dossier, normalized coefficient definition,
assumptions, and either a conditional inequality
b\ensuremath{{}_{2}}\ensuremath{\ge}L or an explanation of the missing
hypothesis. If optimizing, export a dual functional with certified
kernel positivity and all truncation errors. A surviving coefficient is
``not excluded by these tests,'' not ``UV complete.''

\subsection{Validation and rejection controls}

Check s\ensuremath{\leftrightarrow}u symmetry, dimensional consistency,
and pole residues. Reproduce the nongravitational forward positivity
argument in the \ensuremath{\kappa}\ensuremath{\to}0 limit. Use a fully
specified example amplitude to test the subtraction algebra. Perturb a
certificate beyond its margin and verify rejection.

\subsection{Milestones and research extension}

Baseline: one correct sum rule and its gravity-free limit. Strong: a
numerically useful certified conditional bound with a finite positive
margin. Research extension: improve the bound or relax a high-energy
assumption; label any sharpness claim with both matching constructions
and upper/lower evidence.

Milestones are acceptance gates, not calendar promises. Before a run,
freeze the concrete instance/domain, tolerance, compute budget, and
permitted dependencies. A bounded search that fails to find a result
must preserve an unresolved outcome.

\subsection{Starting instruction}

Work on the baseline above. First make every input and convention
explicit, then validate the reference case. Build the checker alongside
the calculation. Label numerical evidence, finite-model certificates,
and continuum theorems separately; return the strongest justified result
and identify the exact unmet proof obligation.

\subsection{Primary references and their role}

\begin{itemize}
\item
  \href{https://arxiv.org/abs/2007.15009}{Tokuda, Aoki and Hirano,
  Gravitational positivity bounds} --- Regge assumptions and finite
  gravitational corrections.
\item
  \href{https://arxiv.org/abs/2007.12667}{Alberte et al., Positivity
  Bounds and the Massless Spin-2 Pole} --- Why naive pole-subtracted
  forward positivity can fail.
\end{itemize}

References support definitions or limitations; the scope choices and
proposed milestones are editorial recommendations, not claims that these
results have already been obtained.

\subsection{Displayed formulation and claim boundary}

\href{https://github.com/rick-stevens-ai/Pure-Thought/blob/main/PRDs/02-Gravitational-Positivity-Bounds.md}{Canonical
specification}.

\[s+t+u=4m^2,\qquad \nu=s-2m^2+\frac t2,\qquad
b_2(t)=\frac12\partial_\nu^2 B(0,t).\]

\[B=A-A_{\mathrm{light\ poles}}-A_{\mathrm{chosen\ low\ cuts}},\qquad
b_2(t)=I_{\mathrm{UV}}(t)+R_\infty(t).\]

The last equality is the sum rule to derive, with its kernel,
subtractions and contour assumptions. Only a justified smearing or
regulated forward prescription can turn it into a positivity statement.
If that prescription proves \(I_{\mathrm{UV}}\ge0\) and
\(R_\infty\ge-R_*\), then \(b_2\ge-R_*\). An uncontrolled contour term
must remain unknown.

\begin{center}\rule{0.5\linewidth}{0.5pt}\end{center}

\section{03. Celestial amplitudes: distributional crossing and
soft-limit consistency}

\textbf{Revised specification --- 2026-09-11.}
\href{https://github.com/rick-stevens-ai/Pure-Thought/blob/main/reviews/03-Celestial-CFT-Bootstrap.md}{Detailed
critique} \ensuremath{\cdot}
\href{https://github.com/rick-stevens-ai/Pure-Thought/blob/main/archive/original-PRDs/03-Celestial-CFT-Bootstrap.md}{Original
description} \ensuremath{\cdot}
\href{https://github.com/rick-stevens-ai/Pure-Thought/blob/main/reviews/EVIDENCE-STANDARD.md}{Evidence
standard}

\subsection{Plain-language guide}

\subsubsection{The problem in everyday terms}

Instead of describing a collision by particle energies and momenta,
imagine recording the directions in which particles enter or leave on a
sphere surrounding the experiment. Celestial methods reorganize
scattering data into quantities associated with that sphere. The
ambition is a new language for gravity. Before using that language to
discover new theories, this problem checks that it faithfully preserves
familiar facts about a simple graviton collision.

\subsubsection{Key terms}

\begin{itemize}
\item
  \textbf{Celestial sphere} --- The sphere of directions seen from a
  point, used here to label asymptotic particle directions rather than a
  physical shell.
\item
  \textbf{Graviton} --- The hypothetical quantum excitation of the
  gravitational field used in perturbative quantum-gravity calculations.
\item
  \textbf{Helicity} --- The component of a particle's spin along its
  direction of motion.
\item
  \textbf{Mellin transform} --- A mathematical transformation that
  reorganizes dependence on energy into dependence on scaling behavior.
\item
  \textbf{Conformal weight \ensuremath{\Delta}} --- A label describing
  scaling behavior in the celestial description; it can be complex.
\item
  \textbf{Distribution} --- A generalized function, such as a delta
  function, understood through integration against test functions rather
  than ordinary pointwise values.
\item
  \textbf{Crossing} --- The relationship between different scattering
  processes obtained by reinterpreting an incoming particle as an
  outgoing antiparticle, with appropriate continuation.
\item
  \textbf{Soft limit} --- The limit in which one particle's energy
  approaches zero; a soft theorem controls the resulting leading
  behavior.
\item
  \textbf{Regulator} --- A temporary modification that makes a singular
  expression well defined, accompanied by a rule for removing or
  interpreting it.
\end{itemize}

\subsubsection{Why this matters}

A successful reformulation could expose symmetries or constraints that
are difficult to see in ordinary scattering variables. That would give
theorists another way to investigate gravity. But a change of
mathematical language is useful only if singularities, conservation laws
and limiting cases survive the translation correctly.

\subsubsection{What progress would mean}

A verified example would supply a dependable foundation for later
celestial calculations. It would establish that specific identities
translate correctly, rather than establish a complete two-dimensional
theory of the universe.

\subsection{Core question}

Can a regulated celestial transform of a specified four-graviton tree
amplitude reproduce Lorentz covariance, crossing, and the leading
conformally soft residue in one consistent convention?

\subsection{Scope and assumptions}

Fix four-dimensional Einstein gravity at tree level, a nonzero
four-graviton helicity amplitude, incoming/outgoing signs
\ensuremath{\epsilon}\ensuremath{{}_{i}}, and
p\ensuremath{{}_{i}}=\ensuremath{\epsilon}\ensuremath{{}_{i}}\ensuremath{\omega}\ensuremath{{}_{i}}q(z\ensuremath{{}_{i}},\ensuremath{\bar{z}}\ensuremath{{}_{i}}).
Declare metric, helicity, normalization, and analytic-continuation
conventions. The target is an identity pipeline for amplitudes, not a
classification of celestial CFTs.

Use
\ensuremath{\Delta}\ensuremath{{}_{i}}=1+i\ensuremath{\lambda}\ensuremath{{}_{i}}
on the principal series where appropriate and
h\ensuremath{{}_{i}}=(\ensuremath{\Delta}\ensuremath{{}_{i}}+J\ensuremath{{}_{i}})/2,
\ensuremath{\bar{h}}\ensuremath{{}_{i}}=(\ensuremath{\Delta}\ensuremath{{}_{i}}\ensuremath{-}J\ensuremath{{}_{i}})/2.
Identify explicitly which continuations away from that contour are used
for soft residues.

\subsection{Mathematical target}

Define the transformed distribution by

\[\widetilde{\mathcal A}=\prod_i\int_0^\infty d\omega _i\,\omega _i^{\Delta _i-1}\;\delta ^{(4)}\!\left(\sum_i \epsilon _i\omega _iq_i\right)\mathcal M_4.\]

State a regulator or test-function space making the integral meaningful,
and specify regulator removal or meromorphic continuation. From a
leading momentum-space term proportional to
\ensuremath{\omega}\_s\ensuremath{{}^{-1}}, derive the Mellin pole at
\ensuremath{\Delta}\_s=1 in this convention. Subleading residues require
their own derivation. Crossing must transform helicities, signs, and
support, not just permute ordinary Euclidean functions.

\subsection{Required outputs and proof obligations}

Produce exact momentum-space expressions, regulated transform formulas,
support conditions, and symbolic covariance/crossing/soft identities. If
a transform cannot be defined under the chosen assumptions, return a
precise obstruction at that step. A positive bootstrap is an extension
only after deriving an applicable inner product and positivity
statement.

\subsection{Validation and rejection controls}

Check the original amplitude's factorization first. Test Mellin formulas
against analytic one-variable examples in their convergence strips.
Verify soft residues before and after transformation under stated
limit-interchange hypotheses. Use smeared distributions for numerical
checks; never evaluate a delta distribution as an ordinary function.

\subsection{Milestones and research extension}

Baseline: one four-point transform with explicit support and checked
leading soft residue. Strong: a second helicity/process example and a
verified crossing relation. Research extension: derive a justified
positive constraint in a restricted subsector before attempting an SDP
island.

Milestones are acceptance gates, not calendar promises. Before a run,
freeze the concrete instance/domain, tolerance, compute budget, and
permitted dependencies. A bounded search that fails to find a result
must preserve an unresolved outcome.

\subsection{Starting instruction}

Work on the baseline above. First make every input and convention
explicit, then validate the reference case. Build the checker alongside
the calculation. Label numerical evidence, finite-model certificates,
and continuum theorems separately; return the strongest justified result
and identify the exact unmet proof obligation.

\subsection{Primary references and their role}

\begin{itemize}
\tightlist
\item
  \href{https://www.pure.ed.ac.uk/ws/files/121747399/Adamo_2019_Class._Quantum_Grav._36_205018_1_.pdf}{Adamo
  et al., Celestial amplitudes and conformal soft theorems} ---
  Conformally soft graviton operators at \ensuremath{\Delta}=1 and
  \ensuremath{\Delta}=0.
\end{itemize}

References support definitions or limitations; the scope choices and
proposed milestones are editorial recommendations, not claims that these
results have already been obtained.

\subsection{Displayed formulation and claim boundary}

\href{https://github.com/rick-stevens-ai/Pure-Thought/blob/main/PRDs/03-Celestial-CFT-Bootstrap.md}{Canonical
specification}.

\[p_i=\epsilon_i\omega_i q_i,\quad
\Delta_i=1+i\lambda_i,\quad
h_i=\frac{\Delta_i+J_i}{2},\quad
\bar h_i=\frac{\Delta_i-J_i}{2}.\]

\[\widetilde{\mathcal A}=
\left(\prod_{i=1}^4\int_0^\infty d\omega_i\,\omega_i^{\Delta_i-1}\right)
\delta^{(4)}\!\left(\sum_i\epsilon_i\omega_iq_i\right)\mathcal M_4.\]

The expression is a regulated distribution. A leading soft term
proportional to \(\omega_s^{-1}\) yields the local integral
\(\int_0^\Lambda\omega_s^{\Delta_s-2}d\omega_s=\Lambda^{\Delta_s-1}/(\Delta_s-1)\)
before continuation, exposing the pole at \(\Delta_s=1\). Support,
helicity, crossing and regulator removal remain proof obligations;
ordinary Euclidean OPE positivity is not assumed.

\begin{center}\rule{0.5\linewidth}{0.5pt}\end{center}

\section{04. Lightcone bootstrap with controlled large-spin remainders}

\textbf{Revised specification --- 2026-09-11.}
\href{https://github.com/rick-stevens-ai/Pure-Thought/blob/main/reviews/04-Modular-Lightcone-Bootstrap.md}{Detailed
critique} \ensuremath{\cdot}
\href{https://github.com/rick-stevens-ai/Pure-Thought/blob/main/archive/original-PRDs/04-Modular-Lightcone-Bootstrap.md}{Original
description} \ensuremath{\cdot}
\href{https://github.com/rick-stevens-ai/Pure-Thought/blob/main/reviews/EVIDENCE-STANDARD.md}{Evidence
standard}

\subsection{Plain-language guide}

\subsubsection{The problem in everyday terms}

A quantum correlation can be calculated by grouping the participating
excitations in different ways, and all valid groupings must agree. That
agreement can force some excitations to exist and restrict their
properties. This problem studies the pattern at large spin, where useful
approximations become available. The central challenge is to say not
only what the leading approximation predicts, but also how large its
error could be.

\subsubsection{Key terms}

\begin{itemize}
\item
  \textbf{Correlation function} --- A quantity describing how field
  measurements at different spacetime points are related.
\item
  \textbf{Four-point function} --- A correlation involving four field
  insertions; it can be reorganized into different pairings.
\item
  \textbf{Bootstrap} --- A method that constrains a theory by requiring
  its own consistency rules to agree, rather than starting from a
  complete microscopic model.
\item
  \textbf{Scaling dimension \ensuremath{\Delta}} --- A measure of how a
  field or excitation changes when all lengths are rescaled.
\item
  \textbf{Spin \ensuremath{\ell}} --- A label describing how an
  excitation transforms under rotations.
\item
  \textbf{Lightcone limit} --- A limit in which two spacetime points
  approach a separation that light could traverse; it organizes which
  contributions dominate the correlation function.
\item
  \textbf{Twist} --- The combination
  \ensuremath{\Delta}\ensuremath{-}\ensuremath{\ell}, useful for
  ordering contributions near a lightlike separation.
\item
  \textbf{Operator product expansion (OPE)} --- A rule expressing the
  effect of nearby field insertions as a sum of other operators with
  specified coefficients.
\item
  \textbf{Double-twist family} --- A sequence of large-spin operators
  whose dimensions approach a characteristic pattern related to two
  external operators.
\item
  \textbf{Remainder bound} --- A guaranteed limit on the terms omitted
  from an approximation.
\end{itemize}

\subsubsection{Why this matters}

Consistency can reveal features shared by many strongly interacting
theories, even when their detailed dynamics cannot be solved. Large-spin
patterns are also useful in theories proposed to describe gravity
through a lower-dimensional boundary. Controlled errors turn suggestive
asymptotic patterns into statements that can actually constrain a finite
calculation.

\subsubsection{What progress would mean}

The baseline would reproduce an established large-spin pattern with
unambiguous conventions. A stronger result would show how far that
pattern can deviate at specified spins, making later comparisons and
exclusions more trustworthy.

\subsection{Core question}

For one scalar four-point function in a unitary CFT with d\textgreater2,
what controlled bounds on large-spin double-twist data follow from
specified low-twist exchanges?

\subsection{Scope and assumptions}

Start in d=3 with an identical scalar \ensuremath{\phi} of fixed
dimension \ensuremath{\Delta}\ensuremath{\phi}. Specify the
normalization of C\_T, the stress-tensor OPE coefficient fixed by its
Ward identity, and an explicit twist gap above the low-twist exchanges
retained. Use a large-C\_T expansion only when a factorization
assumption and its error order are stated.

This challenge studies analytic large-spin behavior. Challenge 07
instead studies the spinning stress-tensor correlator. Torus modular
invariance is restricted to a separate d=2 variant.

\subsection{Mathematical target}

For a fixed scalar block convention, define
F\_\{\ensuremath{\Delta},\ensuremath{\ell}\}=v\textsuperscript{\{\ensuremath{\Delta}\ensuremath{\phi}\}g\_\{\ensuremath{\Delta},\ensuremath{\ell}\}(u,v)\ensuremath{-}u}\{\ensuremath{\Delta}\ensuremath{\phi}\}g\_\{\ensuremath{\Delta},\ensuremath{\ell}\}(v,u),
and use F\_id+\ensuremath{\Sigma}
\ensuremath{\lambda}\ensuremath{{}^{2}}F=0. Twist is
\ensuremath{\tau}=\ensuremath{\Delta}\ensuremath{-}\ensuremath{\ell}.
The double-twist families have
\ensuremath{\Delta}\_\{n,\ensuremath{\ell}\}=2\ensuremath{\Delta}\ensuremath{\phi}+2n+\ensuremath{\ell}+\ensuremath{\gamma}\_\{n,\ensuremath{\ell}\}
asymptotically.

Choose n=0 initially and derive a retained contribution
\ensuremath{\hat{\gamma}}\emph{\{0,\ensuremath{\ell}\}, together with a
bound
|\ensuremath{\gamma}}\{0,\ensuremath{\ell}\}\ensuremath{-}\ensuremath{\hat{\gamma}}\_\{0,\ensuremath{\ell}\}|\ensuremath{\le}R(\ensuremath{\ell})
for
\ensuremath{\ell}\ensuremath{\ge}\ensuremath{\ell}\ensuremath{{}_{0}}
under explicit spectral/Regge assumptions. If only an asymptotic series
is obtained, label it as such and do not manufacture a finite-spin
bound.

\subsection{Required outputs and proof obligations}

Export block normalizations, low-twist input, the derivation or
inversion formula, coefficients, and a decomposition of R into omitted
exchanges, expansion terms, and numerical errors. State clearly whether
the result is an identity, a conditional bound, or an asymptotic
estimate.

\subsection{Validation and rejection controls}

Use generalized-free-field crossing as an algebraic benchmark, while
noting that it is not by itself the target local CFT with a finite
stress tensor. Recover double-twist accumulation in a known analytic
example. Verify every uniform estimate on the specified cross-ratio
domain and all spins covered by the claim.

\subsection{Milestones and research extension}

Baseline: reproduce an established leading large-spin result with
conventions checked. Strong: obtain one explicit remainder estimate, or
identify the extra data needed to obtain it. Research extension: a
tighter conditional bound on low-twist exchanges or bulk couplings with
a stated holographic dictionary.

Milestones are acceptance gates, not calendar promises. Before a run,
freeze the concrete instance/domain, tolerance, compute budget, and
permitted dependencies. A bounded search that fails to find a result
must preserve an unresolved outcome.

\subsection{Starting instruction}

Work on the baseline above. First make every input and convention
explicit, then validate the reference case. Build the checker alongside
the calculation. Label numerical evidence, finite-model certificates,
and continuum theorems separately; return the strongest justified result
and identify the exact unmet proof obligation.

\subsection{Primary references and their role}

\begin{itemize}
\tightlist
\item
  \href{https://arxiv.org/abs/1212.3616}{Fitzpatrick et al., The
  Analytic Bootstrap and AdS Superhorizon Locality} --- Large-spin
  double-twist structure and analytic bootstrap.
\end{itemize}

References support definitions or limitations; the scope choices and
proposed milestones are editorial recommendations, not claims that these
results have already been obtained.

\subsection{Displayed formulation and claim boundary}

\href{https://github.com/rick-stevens-ai/Pure-Thought/blob/main/PRDs/04-Modular-Lightcone-Bootstrap.md}{Canonical
specification}.

\[F_{\Delta,\ell}=v^{\Delta_\phi}g_{\Delta,\ell}(u,v)
-u^{\Delta_\phi}g_{\Delta,\ell}(v,u),\qquad
F_{\mathrm{id}}+\sum_{\mathcal O}\lambda_{\phi\phi\mathcal O}^{2}F_{\Delta,\ell}=0.\]

\[\Delta_{n,\ell}=2\Delta_\phi+2n+\ell+\gamma_{n,\ell},\qquad
|\gamma_{0,\ell}-\widehat\gamma_{0,\ell}|\le R(\ell)\quad(\ell\ge\ell_0).\]

The inequality is a target requiring explicit spectral and Regge
assumptions and a proved remainder. An asymptotic series alone does not
imply this finite-spin inequality. Higher-dimensional crossing is not
torus modular invariance.

\begin{center}\rule{0.5\linewidth}{0.5pt}\end{center}

\section{05. Positive-geometry tests for a fixed gravity integrand}

\textbf{Revised specification --- 2026-09-11.}
\href{https://github.com/rick-stevens-ai/Pure-Thought/blob/main/reviews/05-Positive-Geometry-Gravity.md}{Detailed
critique} \ensuremath{\cdot}
\href{https://github.com/rick-stevens-ai/Pure-Thought/blob/main/archive/original-PRDs/05-Positive-Geometry-Gravity.md}{Original
description} \ensuremath{\cdot}
\href{https://github.com/rick-stevens-ai/Pure-Thought/blob/main/reviews/EVIDENCE-STANDARD.md}{Evidence
standard}

\subsection{Plain-language guide}

\subsubsection{The problem in everyday terms}

Particle-collision formulas often contain many terms whose cancellations
hide a simpler structure. In some theories, a geometric shape organizes
the answer: its boundaries encode the ways a collision can split into
simpler events. This problem asks whether a tightly specified class of
shapes can do the same for one gravity calculation. It treats geometry
as a possible explanation for the structure of a formula, not as a
literal picture of spacetime.

\subsubsection{Key terms}

\begin{itemize}
\item
  \textbf{Integrand} --- The mathematical expression to be integrated to
  obtain a quantity such as a scattering amplitude.
\item
  \textbf{Loop} --- A feature of a quantum correction involving internal
  momenta that must be integrated over; one loop is the first such level
  here.
\item
  \textbf{Supergravity} --- A theory combining gravity with
  supersymmetry, a proposed relation between types of particles; N=8
  specifies the chosen four-dimensional model.
\item
  \textbf{Positive geometry} --- A region with mathematical boundary
  properties that determine a special differential form.
\item
  \textbf{Canonical form} --- That distinguished differential form,
  whose singularities and boundary residues obey precise recursive
  rules.
\item
  \textbf{Residue} --- The coefficient describing the leading
  contribution of a simple pole; here it encodes boundary information.
\item
  \textbf{Factorization} --- The way a scattering process separates into
  simpler processes near certain physical singularities.
\item
  \textbf{Ansatz} --- A deliberately restricted class of candidate
  expressions or geometric constructions searched for a solution.
\item
  \textbf{Pole at infinity} --- A singularity revealed when a coordinate
  is taken to infinity, which must be included in a global comparison.
\end{itemize}

\subsubsection{Why this matters}

A geometric organization could explain cancellations and suggest more
efficient calculations or new identities in quantum gravity. Even a
carefully limited impossibility theorem is informative: it tells us
which geometric proposals are too restrictive and where a different idea
is needed.

\subsubsection{What progress would mean}

Success would mean either a verified shape-to-form correspondence or a
proof that the declared candidate class cannot work. Failure to find a
shape would not rule out every possible geometric description of
gravity.

\subsection{Core question}

For a specified low-loop supergravity integrand, does a precisely
bounded class of candidate positive geometries reproduce its
differential form, including every boundary and pole at infinity?

\subsection{Scope and assumptions}

Start with the four-point one-loop N=8 supergravity integrand, in a
declared four-dimensional representation and helicity sector. Record any
regulator and the limits of four-dimensional cut reconstruction. Derive
or reproduce the target before searching. Tree-level and scalar-box
canonical forms are calibration examples.

Define the search class before claiming an obstruction: ambient
coordinates, dimension, allowed rational maps, inequalities, permitted
external prefactors, and bounds on degree/number of boundaries. ``No
geometry exists'' is not a finite computational question without such
restrictions.

\subsection{Mathematical target}

Compare the top-form
\ensuremath{\Omega}\_target=I(\ensuremath{\ell};k)d\ensuremath{{}^{4}}\ensuremath{\ell}
to f(k)\ensuremath{\pi}\_*\ensuremath{\Omega}\_G for a declared map
\ensuremath{\pi} and external factor f.~State whether equality is
literal rational-form equality or equality modulo a specified class of
terms integrating to zero. Work on a defined real slice with
orientations.

For a standard positive geometry, \ensuremath{\Omega}\_G has logarithmic
boundary poles and recursively normalized residues. Check all finite
boundaries, their intersections, and infinity. Distinguish a canonical
differential form from the symbol of an integrated polylogarithmic
function.

\subsection{Required outputs and proof obligations}

Return either an explicit G, map, canonical form, and exact equality
check, or a proof excluding the declared ansatz class. A failed search
returns ``unresolved within budget.'' A dlog decomposition alone is a
candidate diagnostic, not a geometry certificate.

\subsection{Validation and rejection controls}

Verify simplex/box normalization independently. Check factorization and
generalized cuts of the target. Rational reconstruction must be followed
by an exact polynomial identity after denominators are cleared, with
singular loci accounted for. Check orientation signs and residues at
infinity explicitly.

\subsection{Milestones and research extension}

Baseline: exact target and pole/residue audit. Strong: one verified
representation or one restricted no-go theorem. Research extension:
broaden the geometry class or multiplicity; explain how the result
changes understanding of gravity amplitudes rather than treating a
computational failure as a universal obstruction.

Milestones are acceptance gates, not calendar promises. Before a run,
freeze the concrete instance/domain, tolerance, compute budget, and
permitted dependencies. A bounded search that fails to find a result
must preserve an unresolved outcome.

\subsection{Starting instruction}

Work on the baseline above. First make every input and convention
explicit, then validate the reference case. Build the checker alongside
the calculation. Label numerical evidence, finite-model certificates,
and continuum theorems separately; return the strongest justified result
and identify the exact unmet proof obligation.

\subsection{Primary references and their role}

\begin{itemize}
\tightlist
\item
  \href{https://arxiv.org/abs/1703.04541}{Arkani-Hamed, Bai and Lam,
  Positive Geometries and Canonical Forms} --- Definition of logarithmic
  canonical forms and boundary recursion.
\end{itemize}

References support definitions or limitations; the scope choices and
proposed milestones are editorial recommendations, not claims that these
results have already been obtained.

\subsection{Displayed formulation and claim boundary}

\href{https://github.com/rick-stevens-ai/Pure-Thought/blob/main/PRDs/05-Positive-Geometry-Gravity.md}{Canonical
specification}.

\[\Omega_{\mathrm{target}}=I(\ell;k)\,d^4\ell,
\qquad \Omega_{\mathrm{target}}\stackrel{?}{=}f(k)\,\pi_*\Omega_G.\]

For a standard positive geometry, each oriented boundary must satisfy
the canonical residue relation
\(\operatorname{Res}_{\partial G}\Omega_G=\Omega_{\partial G}\). Fix the
geometry class, maps, prefactors, representation and equality convention
first. Verify finite boundaries, intersections and infinity; an
obstruction is limited to that class. An integrated symbol is not the
integrand's canonical form.

\begin{center}\rule{0.5\linewidth}{0.5pt}\end{center}

\section{06. S-matrix bootstrap with an explicit gravitational infrared
prescription}

\textbf{Revised specification --- 2026-09-11.}
\href{https://github.com/rick-stevens-ai/Pure-Thought/blob/main/reviews/06-Nonperturbative-S-matrix-Bootstrap.md}{Detailed
critique} \ensuremath{\cdot}
\href{https://github.com/rick-stevens-ai/Pure-Thought/blob/main/archive/original-PRDs/06-Nonperturbative-S-matrix-Bootstrap.md}{Original
description} \ensuremath{\cdot}
\href{https://github.com/rick-stevens-ai/Pure-Thought/blob/main/reviews/EVIDENCE-STANDARD.md}{Evidence
standard}

\subsection{Plain-language guide}

\subsubsection{The problem in everyday terms}

Suppose we do not know a theory's microscopic machinery, but do know
that collision probabilities must be sensible and different descriptions
of the same process must agree. Can those rules alone rule out a
proposed interaction strength? This problem searches the space of
scattering possibilities rather than choosing one detailed microscopic
theory. Because gravity acts over long distances and allows very
low-energy radiation, defining the scattering question correctly is
itself part of the work.

\subsubsection{Key terms}

\begin{itemize}
\item
  \textbf{S-matrix} --- The object relating incoming quantum states to
  outgoing states in a scattering experiment.
\item
  \textbf{Partial wave} --- A contribution to scattering organized by
  angular momentum, which makes some probability constraints easier to
  state.
\item
  \textbf{Unitarity} --- Conservation of total probability; an allowed
  scattering process cannot produce a total probability greater than
  one.
\item
  \textbf{Elastic scattering} --- Scattering that stays within the
  chosen particle channel without producing additional particles or
  excitations.
\item
  \textbf{Crossing symmetry} --- Relations between amplitudes for
  processes with differently assigned incoming and outgoing particles.
\item
  \textbf{Infrared prescription} --- A stated way to handle low-energy
  radiation and long-distance singularities.
\item
  \textbf{Outer relaxation} --- A larger, easier-to-test set containing
  all valid possibilities under the assumptions. Excluding a point from
  it can exclude the original possibility.
\item
  \textbf{Inner ansatz} --- A selected subset of candidate amplitudes;
  failing to find a solution there need not exclude other amplitudes.
\item
  \textbf{Dual certificate} --- An independently checkable mathematical
  witness establishing a bound or inconsistency.
\end{itemize}

\subsubsection{Why this matters}

The method could identify restrictions shared by many possible
microscopic theories. That is useful when no single theory is
established. Its reliability depends on not mistaking a limitation of
the search formula, grid or cutoff for a limitation imposed by physics.

\subsubsection{What progress would mean}

A certified exclusion would eliminate a parameter choice under the
stated scattering assumptions. The gravity-free reference case first
checks the machinery; adding a defensible treatment of gravity would be
the important next step.

\subsection{Core question}

Can one certify an exclusion for a low-energy scalar-scattering
parameter using a crossing-symmetric partial-wave relaxation whose
infrared and truncation assumptions are explicit?

\subsection{Scope and assumptions}

Begin with identical massive scalars in four dimensions and no gravity
to validate the machinery. The gravitational stage must first define an
infrared-finite observable or a justified regulator/pole treatment,
including massless radiation. Keep one scalar mass and one amplitude
coefficient fixed; optimize a second coefficient in a declared
normalization.

This is complementary to Challenge 02: the focus is amplitude
feasibility and partial-wave constraints rather than deriving a
low-energy positivity sum rule. A nongravitational baseline does not
count as solving the gravitational stage.

\subsection{Mathematical target}

Choose \(A=16\pi\sum_{\ell\ge0}(2\ell+1)a_\ell(s)P_\ell(\cos\theta)\),
s+t+u=4m\ensuremath{{}^{2}},
\ensuremath{\rho}=\ensuremath{\surd}(1\ensuremath{-}4m\ensuremath{{}^{2}}/s),
and S\_\ensuremath{\ell}=1+2i\ensuremath{\rho}a\_\ensuremath{\ell}. Then
|S\_\ensuremath{\ell}|\ensuremath{\le}1 is equivalent to Im
a\_\ensuremath{\ell}\ensuremath{\ge}\ensuremath{\rho}|a\_\ensuremath{\ell}|\ensuremath{{}^{2}}.
Equality needs an actually elastic sector.

State how crossing is imposed and distinguish an outer relaxation of
necessary constraints from an inner finite amplitude ansatz. A dual
exclusion from an outer relaxation transfers to the full problem only
when every physical amplitude maps into that relaxation. Bound energy,
angular-momentum, and basis tails; checking a finite mesh is
insufficient.

\subsection{Required outputs and proof obligations}

Export the amplitude conventions, IR prescription, finite optimization
problem, mapping from the continuum problem, and dual certificate with
its positivity/error margins. Feasible ansatz points are examples
satisfying the verified constraints only. If the IR construction or tail
control fails, state which conclusion remains unavailable.

\subsection{Validation and rejection controls}

Check free scattering and a specified analytic nongravitational
benchmark. Verify the unitarity disk algebra independently. Test between
grid points and beyond cutoffs using proved bounds. Demonstrate that
known admissible examples are not falsely excluded. Increasing precision
is a diagnostic; rigorous enclosures are needed for certification.

\subsection{Milestones and research extension}

Baseline: certified finite nongravitational relaxation. Strong: justify
an IR treatment and certify one conditional gravitational exclusion.
Research extension: improved continuum bounds or a constructive
amplitude with a precisely delimited consistency claim.

Milestones are acceptance gates, not calendar promises. Before a run,
freeze the concrete instance/domain, tolerance, compute budget, and
permitted dependencies. A bounded search that fails to find a result
must preserve an unresolved outcome.

\subsection{Starting instruction}

Work on the baseline above. First make every input and convention
explicit, then validate the reference case. Build the checker alongside
the calculation. Label numerical evidence, finite-model certificates,
and continuum theorems separately; return the strongest justified result
and identify the exact unmet proof obligation.

\subsection{Primary references and their role}

\begin{itemize}
\tightlist
\item
  \href{https://arxiv.org/abs/2007.12667}{Alberte et al., Positivity
  Bounds and the Massless Spin-2 Pole} --- Infrared and analyticity
  issues that must be resolved before importing standard bounds.
\end{itemize}

References support definitions or limitations; the scope choices and
proposed milestones are editorial recommendations, not claims that these
results have already been obtained.

\subsection{Displayed formulation and claim boundary}

\href{https://github.com/rick-stevens-ai/Pure-Thought/blob/main/PRDs/06-Nonperturbative-S-matrix-Bootstrap.md}{Canonical
specification}.

\[A=16\pi\sum_{\ell\ge0}(2\ell+1)a_\ell(s)P_\ell(\cos\theta),\quad
\rho=\sqrt{1-4m^2/s},\quad S_\ell=1+2i\rho a_\ell.\]

\[|S_\ell|\le1\quad\Longleftrightarrow\quad
\operatorname{Im}a_\ell\ge\rho|a_\ell|^2\qquad(s>4m^2).\]

Equality requires an elastic sector. A certified outer relaxation can
exclude a physical amplitude only after the inclusion of every
admissible amplitude and all tail bounds are proved. Failure of an inner
ansatz excludes only that ansatz. The gravitational stage has a separate
infrared prerequisite.

\begin{center}\rule{0.5\linewidth}{0.5pt}\end{center}

\section{07. Stress-tensor bootstrap: a sector-specific gap bound in
three dimensions}

\textbf{Revised specification --- 2026-09-11.}
\href{https://github.com/rick-stevens-ai/Pure-Thought/blob/main/reviews/07-Extremal-CFTs-Stress-Tensor.md}{Detailed
critique} \ensuremath{\cdot}
\href{https://github.com/rick-stevens-ai/Pure-Thought/blob/main/archive/original-PRDs/07-Extremal-CFTs-Stress-Tensor.md}{Original
description} \ensuremath{\cdot}
\href{https://github.com/rick-stevens-ai/Pure-Thought/blob/main/reviews/EVIDENCE-STANDARD.md}{Evidence
standard}

\subsection{Plain-language guide}

\subsubsection{The problem in everyday terms}

Energy and momentum must be accounted for consistently in a quantum
theory. The stress tensor is the mathematical object that tracks them.
By examining how four stress-tensor insertions relate to one another, we
can test whether a proposed theory is missing an excitation it ought to
contain. The question here is how large a particular gap in the spectrum
can be before the consistency rules contradict each other.

\subsubsection{Key terms}

\begin{itemize}
\item
  \textbf{Stress tensor T} --- The field that describes the local flow
  and distribution of energy and momentum.
\item
  \textbf{Conformal field theory} --- A quantum theory with symmetry
  under angle-preserving transformations, including rescaling lengths.
\item
  \textbf{Central-charge parameter C\_T} --- The normalization of the
  stress tensor's two-point correlation, used to characterize the theory
  in a fixed convention.
\item
  \textbf{TTT data} --- Coefficients describing a three-stress-tensor
  correlation; they contain information beyond C\_T.
\item
  \textbf{Scalar sector} --- The set of exchanged operators with zero
  spin, chosen as the focus of the gap question.
\item
  \textbf{Parity-even} --- Unchanged in sign under the specified
  spatial-reflection transformation.
\item
  \textbf{Spectral gap} --- Here, the threshold in scaling dimension
  before the first nonidentity operator in the selected sector appears,
  not necessarily a material's energy gap.
\item
  \textbf{Tensor structure} --- One of the distinct ways the directional
  indices of spinning fields can enter a correlation function.
\item
  \textbf{Positive semidefinite matrix} --- A matrix whose quadratic
  form is nonnegative for every vector; it expresses the needed
  positivity for multiple coupled tensor structures.
\end{itemize}

\subsubsection{Why this matters}

The stress tensor is present in a broad class of local theories, so
restrictions derived from it can have wide reach. Such restrictions help
map which strongly interacting quantum theories are possible. They can
also constrain proposed boundary descriptions of gravity, where the
stress tensor has a special role.

\subsubsection{What progress would mean}

A verified gap bound would show that a theory with the declared inputs
must contain a certain kind of excitation below a threshold. It would
not determine the entire spectrum or prove that a gravity theory is
unique.

\subsection{Core question}

At fixed stress-tensor two- and three-point data, can crossing exclude a
proposed gap in the parity-even scalar sector of T\ensuremath{\times}T
in a unitary three-dimensional CFT?

\subsection{Scope and assumptions}

Fix d=3, parity invariance, a unique conserved stress tensor, a
normalization for C\_T, and the remaining independent TTT parameter. The
initial observable is the first parity-even nonidentity scalar appearing
with nonzero coupling in T\ensuremath{\times}T. It is not a bound on
every operator in the theory. State separate gaps for any other
restricted spin/parity sectors.

Use the full conserved spinning correlator. A scalar block toy
implementation is permitted for learning but is not the target
stress-tensor bootstrap.

\subsection{Mathematical target}

Write the crossing equations in a nonredundant tensor basis as

\[V_{id}+V_T(C_T,t_4)+\sum_{\mathcal O}\lambda_{\mathcal O}^{\mathsf T}V_{\Delta ,\ell ,\pi}\lambda_{\mathcal O}=0.\]

The OPE vectors account for multiple tensor structures. In d=3, a
conserved symmetric-traceless spin-\ensuremath{\ell} primary has
\ensuremath{\Delta}=\ensuremath{\ell}+1, not \ensuremath{\Delta}=3 for
all spins. For a proposed scalar gap g, seek \ensuremath{\alpha} with
\ensuremath{\alpha}(V\_id+V\_T)>0 and
\ensuremath{\alpha}(V\_\{\ensuremath{\Delta},\ensuremath{\ell},\ensuremath{\pi}\})
positive semidefinite on every allowed exchange. This yields a
contradiction and therefore excludes g under the specified assumptions.

\subsection{Required outputs and proof obligations}

Export tensor/block conventions, conservation and Ward identities, fixed
input data, derivative basis, spin treatment, and functional
coefficients. Attach rigorous block-approximation and high-spin control.
If only the finite approximant is checked, label the result accordingly.

\subsection{Validation and rejection controls}

Recover a selected published stress-tensor benchmark using exactly its
conventions and assumptions. Check the allowed region rather than the
excluded region. Verify matrix positivity, not just individual scalar
components. Free scalar/fermion data provide normalization tests;
numerical agreement alone is not a proof of the continuum bound.

\subsection{Milestones and research extension}

Baseline: reproduce one finite benchmark and publish the assumptions.
Strong: turn one exclusion into a certified result with omitted-sector
control. Research extension: improve a bound, add another sector, or
obtain a conditional extremal spectrum. Saturation does not construct a
CFT.

Milestones are acceptance gates, not calendar promises. Before a run,
freeze the concrete instance/domain, tolerance, compute budget, and
permitted dependencies. A bounded search that fails to find a result
must preserve an unresolved outcome.

\subsection{Starting instruction}

Work on the baseline above. First make every input and convention
explicit, then validate the reference case. Build the checker alongside
the calculation. Label numerical evidence, finite-model certificates,
and continuum theorems separately; return the strongest justified result
and identify the exact unmet proof obligation.

\subsection{Primary references and their role}

\begin{itemize}
\tightlist
\item
  \href{https://arxiv.org/abs/1708.05718}{Dymarsky et al., The 3d
  Stress-Tensor Bootstrap} --- Conserved tensor crossing and
  sector-specific numerical gap bounds.
\end{itemize}

References support definitions or limitations; the scope choices and
proposed milestones are editorial recommendations, not claims that these
results have already been obtained.

\subsection{Displayed formulation and claim boundary}

\href{https://github.com/rick-stevens-ai/Pure-Thought/blob/main/PRDs/07-Extremal-CFTs-Stress-Tensor.md}{Canonical
specification}.

\[V_{\mathrm{id}}+V_T(C_T,t_4)+
\sum_{\mathcal O}\lambda_{\mathcal O}^{\mathsf T}V_{\Delta,\ell,\pi}
\lambda_{\mathcal O}=0.\]

\[\alpha(V_{\mathrm{id}}+V_T)>0,\qquad
\alpha(V_{\Delta,\ell,\pi})\succeq0
\quad\text{for every allowed exchange}.\]

These inequalities contradict crossing and exclude the proposed
sector-specific gap. Tensor structures, conserved spinning blocks, Ward
identities and continuous exchange coverage are required. In three
dimensions a conserved symmetric-traceless spin-\(\ell\) primary has
\(\Delta=\ell+1\).

\begin{center}\rule{0.5\linewidth}{0.5pt}\end{center}

\section{08. Exact modular-data and symmetry-anomaly consistency tests}

\textbf{Revised specification --- 2026-09-11.}
\href{https://github.com/rick-stevens-ai/Pure-Thought/blob/main/reviews/08-Swampland-Modularity-Symmetries.md}{Detailed
critique} \ensuremath{\cdot}
\href{https://github.com/rick-stevens-ai/Pure-Thought/blob/main/archive/original-PRDs/08-Swampland-Modularity-Symmetries.md}{Original
description} \ensuremath{\cdot}
\href{https://github.com/rick-stevens-ai/Pure-Thought/blob/main/reviews/EVIDENCE-STANDARD.md}{Evidence
standard}

\subsection{Plain-language guide}

\subsubsection{The problem in everyday terms}

A proposed quantum theory comes with mathematical rules for combining
excitations and transforming its descriptions. These rules should fit
together like the pieces of a puzzle. This problem builds exact checks
for that fit and then asks a second question: can a chosen global
symmetry be made into a gauge symmetry? A theory can be consistent even
when this second operation is obstructed, so the two verdicts must
remain separate.

\subsubsection{Key terms}

\begin{itemize}
\item
  \textbf{Rational CFT} --- A conformal field theory with finitely many
  basic representation types relative to a specified chiral algebra.
\item
  \textbf{Character} --- A generating function counting states within
  one representation of that algebra.
\item
  \textbf{Modular S and T matrices} --- Matrices describing how the
  characters change under two basic transformations of the torus.
\item
  \textbf{Fusion rules} --- Rules specifying which excitation types can
  occur when two types are combined.
\item
  \textbf{Global symmetry} --- A transformation applied consistently
  throughout a system that leaves its laws unchanged.
\item
  \textbf{Gauging} --- Promoting an appropriate symmetry to a local
  redundancy and adding the corresponding gauge structure.
\item
  \textbf{Anomaly} --- An obstruction to implementing a classical
  symmetry or a proposed gauging consistently in the quantum setting; an
  't Hooft anomaly can exist in a valid theory.
\item
  \textbf{Cocycle} --- Algebraic data satisfying a compatibility
  identity, used here to encode an anomaly.
\item
  \textbf{Bulk and boundary} --- The interior gravitational theory and
  its lower-dimensional boundary description; symmetries can have
  different interpretations in the two.
\end{itemize}

\subsubsection{Why this matters}

Exact consistency checks can detect false candidate data before
researchers invest in a full construction. They also clarify which
symmetries can be gauged and which require additional structure, such as
anomaly cancellation or inflow from another system. This is useful both
for studying quantum field theories and for interpreting holographic
gravity proposals.

\subsubsection{What progress would mean}

The first contribution would be a dependable checker with known passing
and failing examples. A failure would identify a particular algebraic
inconsistency or gauging obstruction, rather than issuing an unsupported
verdict about all of quantum gravity.

\subsection{Core question}

For supplied rational-CFT modular data and a specified finite internal
symmetry, which algebraic consistency tests can be certified, and which
failures obstruct a proposed gauging?

\subsection{Scope and assumptions}

Use a bosonic rational CFT with a declared chiral algebra and finitely
many characters. Input exact algebraic S,T matrices, vacuum label, and a
nonnegative integer multiplicity matrix M. Start with the Ising model.
An optional finite 0-form group G requires separate
defect/twisted-sector data and an anomaly cocycle; these cannot be
inferred from the untwisted partition function alone.

The result concerns the supplied CFT data and gauging problem. A general
bulk quantum-gravity compatibility verdict is outside the baseline.

\subsection{Mathematical target}

Check S\ensuremath{\dagger}S=I, symmetry of S where assumed,
S\ensuremath{{}^{2}}=C, (ST)\ensuremath{{}^{3}}=C in the declared T
convention, M\ensuremath{{}_{00}}=1, MS=SM, MT=TM, and
M\ensuremath{{}_{ij}}\ensuremath{\in}Z\ensuremath{\ge}0. Check Verlinde
coefficients \ensuremath{\Sigma}\_a S\_ia S\_ja S*\_ka/S\_0a are
nonnegative integers.

For a finite symmetry cocycle
\ensuremath{\omega}\ensuremath{\in}Z\ensuremath{{}^{3}}(G,U(1)), check
the cocycle identity exactly. To certify anomaly-free gauging in the
stated setting, supply an explicit 2-cochain \ensuremath{\beta} with
\ensuremath{\delta}\ensuremath{\beta}=\ensuremath{\omega}, or prove its
nonexistence in a complete cohomology computation. Distinguish this
obstruction from inconsistency of an anomalous but valid boundary
theory.

\subsection{Required outputs and proof obligations}

Return pass/fail/unresolved for each named axiom, with exact arithmetic
traces. Include the number field and embeddings. Even passing every
modular-data test gives necessary algebraic consistency only; full CFT
sewing and realizability are separate obligations.

\subsection{Validation and rejection controls}

Verify Ising S entries including 1/\ensuremath{\surd}2 in algebraic
arithmetic. Reject altered modular relations, noninteger multiplicities,
and a failed cocycle identity. Accept the existence of a boundary global
symmetry rather than flagging it as automatically inconsistent.
Cross-check a small cyclic-group cohomology example.

\subsection{Milestones and research extension}

Baseline: exact modular-data checker on Ising and one second supplied
model. Strong: one finite-group twisted-sector/gauging analysis.
Research extension: classify a explicitly finite family of data or
derive a conditional bulk consequence using a stated holographic
theorem.

Milestones are acceptance gates, not calendar promises. Before a run,
freeze the concrete instance/domain, tolerance, compute budget, and
permitted dependencies. A bounded search that fails to find a result
must preserve an unresolved outcome.

\subsection{Starting instruction}

Work on the baseline above. First make every input and convention
explicit, then validate the reference case. Build the checker alongside
the calculation. Label numerical evidence, finite-model certificates,
and continuum theorems separately; return the strongest justified result
and identify the exact unmet proof obligation.

\subsection{Primary references and their role}

\begin{itemize}
\tightlist
\item
  \href{https://arxiv.org/abs/1810.05338}{Harlow and Ooguri, Symmetries
  in quantum field theory and quantum gravity} --- Bulk gauge versus
  boundary global symmetry, including higher-form extensions.
\end{itemize}

References support definitions or limitations; the scope choices and
proposed milestones are editorial recommendations, not claims that these
results have already been obtained.

\subsection{Displayed formulation and claim boundary}

\href{https://github.com/rick-stevens-ai/Pure-Thought/blob/main/PRDs/08-Swampland-Modularity-Symmetries.md}{Canonical
specification}.

\[S^\dagger S=I,\quad S^2=C,\quad(ST)^3=C,\quad
MS=SM,\quad MT=TM,\quad M_{00}=1,\quad M_{ij}\in\mathbb Z_{\ge0}.\]

\[N_{ij}^{\ k}=\sum_a\frac{S_{ia}S_{ja}S_{ka}^*}{S_{0a}}
\in\mathbb Z_{\ge0},\qquad \omega\in Z^3(G,U(1)),\quad\delta\beta=\omega.\]

Fix the \(T\) convention and any assumed symmetry of \(S\). An explicit
cochain \(\beta\) trivializes the specified anomaly; a complete
cohomology calculation may prove no such cochain exists. This gauging
obstruction does not make an anomalous boundary theory inconsistent, and
untwisted modular data alone do not supply defect data.

\begin{center}\rule{0.5\linewidth}{0.5pt}\end{center}

\section{09. Certified Chern phases in a bounded tight-binding family}

\textbf{Revised specification --- 2026-09-11.}
\href{https://github.com/rick-stevens-ai/Pure-Thought/blob/main/reviews/09-Topological-Band-Theory.md}{Detailed
critique} \ensuremath{\cdot}
\href{https://github.com/rick-stevens-ai/Pure-Thought/blob/main/archive/original-PRDs/09-Topological-Band-Theory.md}{Original
description} \ensuremath{\cdot}
\href{https://github.com/rick-stevens-ai/Pure-Thought/blob/main/reviews/EVIDENCE-STANDARD.md}{Evidence
standard}

\subsection{Plain-language guide}

\subsubsection{The problem in everyday terms}

Two solids can both resist ordinary electrical conduction in their
interiors yet differ in how their quantum states fit together. In a
topological phase, that global organization cannot be smoothly changed
without passing through a gap closing or leaving the assumptions that
protect it. Here we build simple lattice models and prove which
parameter ranges belong to which phase. The task is closer to drawing an
exact map of a model than predicting a particular material.

\subsubsection{Key terms}

\begin{itemize}
\item
  \textbf{Hamiltonian H} --- The mathematical operator specifying a
  quantum system's energy and dynamics.
\item
  \textbf{Tight-binding model} --- A lattice model in which particles
  occupy localized orbitals and move between them through specified
  hopping terms.
\item
  \textbf{Orbital} --- One of the local quantum states included at a
  lattice site or within a unit cell.
\item
  \textbf{Band} --- A range of allowed energies labeled by crystal
  momentum in a periodic system.
\item
  \textbf{Band gap} --- An energy separation between the selected
  occupied bands and the next bands.
\item
  \textbf{Brillouin zone} --- The space of distinct crystal momenta for
  a periodic lattice.
\item
  \textbf{Chern number} --- An integer measuring a global twist in the
  selected quantum states across a two-dimensional Brillouin zone.
\item
  \textbf{Occupied projector} --- An operator selecting the band states
  being treated as occupied.
\item
  \textbf{Phase diagram} --- A map showing which kind of behavior occurs
  for each model-parameter range.
\end{itemize}

\subsubsection{Why this matters}

Topological models explain why some properties depend on global
structure rather than microscopic details. A certified phase map
provides reliable examples for theory and simulation, and it helps
identify which ingredients a proposed design actually needs. Proving a
gap everywhere is important because a small missed gap closing can
invalidate a computed topological label.

\subsubsection{What progress would mean}

A useful result would give explicit models, guaranteed parameter ranges
and checked Chern numbers. These could guide later physical designs,
while leaving material stability, fabrication and experimental behavior
to additional work.

\subsection{Core question}

Within a declared finite-range Hamiltonian family, which parameter
regions have a proven band gap and occupied-band Chern number, and what
minimality statements follow within that family?

\subsection{Scope and assumptions}

Begin with the two-band square-lattice model
\(H_m(k)=\sin(k_x)\sigma_x+\sin(k_y)\sigma_y+[m+\cos(k_x)+\cos(k_y)]\sigma_z\),
k\ensuremath{\in}[\ensuremath{-}\ensuremath{\pi},\ensuremath{\pi}]\ensuremath{{}^{2}}.
Use one occupied band and a fixed orientation. Analyze rational
parameter intervals separated from m=\ensuremath{-}2,0,2.

For subsequent searches fix dimension, orbitals, filling, hopping
support, coefficient domain, and symmetry representation. A rational
grid is finite; a continuous coefficient box requires a different
completeness argument. ``All Hamiltonians for a space group'' is not an
enumerable specification.

\subsection{Mathematical target}

For a general family \(H(k)=\sum_R t_Re^{ik\cdot R}\), impose
t\_\ensuremath{-}R=t\_R\ensuremath{\dagger} and the stated symmetry
equations. Prove
E\_\{r+1\}(k)\ensuremath{-}E\_r(k)\ensuremath{\ge}\ensuremath{\delta}>0
for all k and all parameters claimed. With P the occupied projector,
choose

\[C=\frac{1}{2\pi i}\int_{BZ}\operatorname{Tr}P[\partial _xP,\partial _yP]\,d^{2}k.\]

Certify C through an analytic degree/homotopy argument or validated
integration with an error enclosure containing a unique integer,
including discretization and projector errors.

\subsection{Required outputs and proof obligations}

Return model coefficients, gap lower bounds, invariant and sign
convention, certified parameter coverage, and uncovered boxes. A
minimality claim must specify whether it minimizes orbitals, hopping
range, or coefficients and must exclude every smaller object in that
declared class.

\subsection{Validation and rejection controls}

Recover gap closings at \ensuremath{-}2,0,2 and trivial/nontrivial
intervals on either side. Reverse orientation and verify C changes sign.
A fine-grid lattice Chern calculation is a useful cross-check, but must
not substitute for the continuous gap proof. Check a
time-reversal-invariant occupied bundle has zero total Chern number.

\subsection{Milestones and research extension}

Baseline: complete phase diagram for H\_m away from transition values.
Strong: a certified atlas for one bounded symmetry-compatible family.
Research extension: new range/band tradeoffs or minimality theorems,
explicitly relative to the permitted family.

Milestones are acceptance gates, not calendar promises. Before a run,
freeze the concrete instance/domain, tolerance, compute budget, and
permitted dependencies. A bounded search that fails to find a result
must preserve an unresolved outcome.

\subsection{Starting instruction}

Work on the baseline above. First make every input and convention
explicit, then validate the reference case. Build the checker alongside
the calculation. Label numerical evidence, finite-model certificates,
and continuum theorems separately; return the strongest justified result
and identify the exact unmet proof obligation.

\subsection{Primary references and their role}

\begin{itemize}
\tightlist
\item
  \href{https://arxiv.org/abs/1311.4956}{Chen et al., The impossibility
  of exactly flat non-trivial Chern bands in strictly local periodic
  tight binding models} --- A concrete example of why range/locality
  assumptions change realizability claims.
\end{itemize}

References support definitions or limitations; the scope choices and
proposed milestones are editorial recommendations, not claims that these
results have already been obtained.

\subsection{Displayed formulation and claim boundary}

\href{https://github.com/rick-stevens-ai/Pure-Thought/blob/main/PRDs/09-Topological-Band-Theory.md}{Canonical
specification}.

\[H_m(k)=\sin k_x\,\sigma_x+\sin k_y\,\sigma_y+
(m+\cos k_x+\cos k_y)\sigma_z.\]

\[E_{r+1}(k)-E_r(k)\ge\delta>0,\qquad
C=\frac{1}{2\pi i}\int_{\mathrm{BZ}}
\operatorname{Tr}\!\left(P[\partial_{k_x}P,\partial_{k_y}P]\right)\,d^2k.\]

Prove the gap over the full claimed parameter and momentum domain. Fix
orientation and the occupied projector. An analytic argument or
validated integral enclosing one integer certifies \(C\); a momentum
mesh alone does not.

\begin{center}\rule{0.5\linewidth}{0.5pt}\end{center}

\section{10. Flatness--locality--geometry tradeoffs for Chern bands}

\textbf{Revised specification --- 2026-09-11.}
\href{https://github.com/rick-stevens-ai/Pure-Thought/blob/main/reviews/10-Flat-Chern-Bands.md}{Detailed
critique} \ensuremath{\cdot}
\href{https://github.com/rick-stevens-ai/Pure-Thought/blob/main/archive/original-PRDs/10-Flat-Chern-Bands.md}{Original
description} \ensuremath{\cdot}
\href{https://github.com/rick-stevens-ai/Pure-Thought/blob/main/reviews/EVIDENCE-STANDARD.md}{Evidence
standard}

\subsection{Plain-language guide}

\subsubsection{The problem in everyday terms}

In an ordinary energy band, particles can have quite different energies
depending on how they move. A nearly flat band reduces those energy
differences, potentially making interactions between particles more
influential. We want that band to also have a nontrivial topological
structure. The problem asks how well these goals can be combined when
particles are allowed to hop only a limited distance and the model has a
limited number of orbitals.

\subsubsection{Key terms}

\begin{itemize}
\item
  \textbf{Flat band} --- A band whose energy is constant across momentum
  space; a nearly flat band has a small but nonzero variation.
\item
  \textbf{Bandwidth W} --- The difference between the largest and
  smallest energy in the selected band.
\item
  \textbf{Gap \ensuremath{\Delta}} --- The minimum energy separation
  between that band and other bands.
\item
  \textbf{Flatness ratio W/\ensuremath{\Delta}} --- A way to compare
  bandwidth to isolation from other bands; smaller values mean a
  flatter, well-separated band.
\item
  \textbf{Hopping range} --- The greatest lattice distance over which
  the model allows direct hopping.
\item
  \textbf{Chern band} --- A band with a nonzero Chern number, indicating
  nontrivial global topology.
\item
  \textbf{Berry curvature} --- A local measure of how quantum states
  twist as momentum changes; its integral determines the Chern number in
  the stated convention.
\item
  \textbf{Quantum metric} --- A measure of how distinguishable nearby
  momentum-dependent quantum states are.
\item
  \textbf{Spectral flattening} --- A mathematical change to band
  energies that preserves the chosen states but generally changes the
  spatial range of hopping.
\end{itemize}

\subsubsection{Why this matters}

Flat topological bands are useful settings for studying collective
quantum behavior driven by interactions. But there are tradeoffs:
restricting the model's spatial range changes which flatness properties
are possible. Quantifying those tradeoffs can show whether a proposed
improvement is attainable or requires a more complicated model.

\subsubsection{What progress would mean}

Progress would provide certified examples and limits within a defined
design family. Favorable single-particle geometry would make a model
interesting to investigate further; it would not by itself prove an
interacting fractional topological phase.

\subsection{Core question}

For fixed orbital count, hopping range, and gap normalization, how small
can a Chern band's bandwidth and geometric nonuniformity be, with
certified bounds?

\subsection{Scope and assumptions}

Choose a two-dimensional isolated rank-one band, target C=1, and a
compact coefficient family with range R and fixed orbital positions.
Initially use a specified two-band family extending Challenge 09. Impose
a positive gap floor and a norm bound to prevent trivial energy
rescaling.

Maintain two distinct tracks: finite-range approximate flatness, and
exact spectral flattening with generally non-finite-range hopping. Do
not require exact nonzero-Chern flatness with strictly finite-range
hopping under the no-go theorem's hypotheses.

\subsection{Mathematical target}

Define W=max\_k E\_n\ensuremath{-}min\_k E\_n and
\ensuremath{\Delta}=min\_k
min\_\{j\ensuremath{\ne}n\}|E\_j\ensuremath{-}E\_n|. Optimize
W/\ensuremath{\Delta}. For the fixed Brillouin-zone coordinates define
g\_ij=Re\ensuremath{\langle}\ensuremath{\partial}\_iu|(1\ensuremath{-}P)|\ensuremath{\partial}\_ju\ensuremath{\rangle}
and F with the same Chern convention as Challenge 09. Then
C=(1/2\ensuremath{\pi})\ensuremath{\int}F and
\ensuremath{\bar{F}}=2\ensuremath{\pi}C/Area(BZ).

Report separately
Var(F)=Area\ensuremath{{}^{-1}}\ensuremath{\int}(F\ensuremath{-}\ensuremath{\bar{F}})\ensuremath{{}^{2}}
and D\_g=Area\ensuremath{{}^{-1}}\ensuremath{\int}(tr
g\ensuremath{-}|F|). These are diagnostics under a specified Euclidean
coordinate metric, not a universal fractional-Chern-insulator stability
score.

\subsection{Required outputs and proof obligations}

Export coefficient/domain constraints, W upper and \ensuremath{\Delta}
lower bounds, certified C, geometric enclosures, and optimization gaps.
An achieved objective gives an upper bound on the minimum; a global
lower bound needs a covering/relaxation proof. For infinite-range
flattening, give decay and truncation estimates.

\subsection{Validation and rejection controls}

Verify normalization \ensuremath{\int}F=2\ensuremath{\pi}C and the
pointwise metric inequality in the chosen convention. Check that
truncating a spectrally flattened model generally reintroduces
bandwidth. Include a trivial exactly flat band as a negative topology
test. Compare candidate objectives using identical units, embedding, and
norm constraints.

\subsection{Milestones and research extension}

Baseline: one certified near-flat Chern model and a correct locality
audit. Strong: global objective bounds on one compact family or a
rigorous truncation tradeoff. Research extension: optimal geometry
within a specified class; many-body FCI stability requires an added
interacting Hamiltonian and filling.

Milestones are acceptance gates, not calendar promises. Before a run,
freeze the concrete instance/domain, tolerance, compute budget, and
permitted dependencies. A bounded search that fails to find a result
must preserve an unresolved outcome.

\subsection{Starting instruction}

Work on the baseline above. First make every input and convention
explicit, then validate the reference case. Build the checker alongside
the calculation. Label numerical evidence, finite-model certificates,
and continuum theorems separately; return the strongest justified result
and identify the exact unmet proof obligation.

\subsection{Primary references and their role}

\begin{itemize}
\tightlist
\item
  \href{https://arxiv.org/abs/1311.4956}{Chen et al., Flat Chern-band
  locality obstruction} --- No-go theorem for simultaneous exact
  flatness, nonzero Chern number, and strictly local hopping.
\end{itemize}

References support definitions or limitations; the scope choices and
proposed milestones are editorial recommendations, not claims that these
results have already been obtained.

\subsection{Displayed formulation and claim boundary}

\href{https://github.com/rick-stevens-ai/Pure-Thought/blob/main/PRDs/10-Flat-Chern-Bands.md}{Canonical
specification}.

\[W=\max_k E_n(k)-\min_k E_n(k),\quad
\Delta=\min_k\min_{j\ne n}|E_j(k)-E_n(k)|,\quad\text{minimize }W/\Delta.\]

\[g_{ij}=\operatorname{Re}\langle\partial_i u|(1-P)|\partial_j u\rangle,
\quad \bar F=\frac{2\pi C}{A},\quad A=\operatorname{Area}(\mathrm{BZ}),\]

\[\operatorname{Var}(F)=\frac1A\int_{\mathrm{BZ}}(F-\bar F)^2d^2k,
\qquad D_g=\frac1A\int_{\mathrm{BZ}}(\operatorname{tr}g-|F|)d^2k.\]

Fix a gap floor, norm bound, orbital positions and coordinate metric.
Use the Chern convention of 09. Strictly finite-range approximate
flatness and generally long-range exact flattening are separate tasks;
these diagnostics alone do not prove an interacting fractional phase.

\begin{center}\rule{0.5\linewidth}{0.5pt}\end{center}

\section{11. Certified topological bands in an idealized Maxwell medium}

\textbf{Revised specification --- 2026-09-11.}
\href{https://github.com/rick-stevens-ai/Pure-Thought/blob/main/reviews/11-Photonic-Topological-Crystals.md}{Detailed
critique} \ensuremath{\cdot}
\href{https://github.com/rick-stevens-ai/Pure-Thought/blob/main/archive/original-PRDs/11-Photonic-Topological-Crystals.md}{Original
description} \ensuremath{\cdot}
\href{https://github.com/rick-stevens-ai/Pure-Thought/blob/main/reviews/EVIDENCE-STANDARD.md}{Evidence
standard}

\subsection{Plain-language guide}

\subsubsection{The problem in everyday terms}

A repeating optical structure changes which light waves can travel
through it. With suitable nonreciprocal response, it may also support
light that travels along an interface in a preferred direction. This
problem asks whether a precisely defined ideal electromagnetic medium
really has the required band structure and topology. It separates a
mathematically guaranteed wave effect from the later task of finding a
material that realizes the assumed response.

\subsubsection{Key terms}

\begin{itemize}
\item
  \textbf{Maxwell equations} --- The equations governing classical
  electric and magnetic fields.
\item
  \textbf{Periodic medium} --- A medium whose properties repeat from
  cell to cell.
\item
  \textbf{Permittivity \ensuremath{\epsilon} and permeability
  \ensuremath{\mu}} --- Quantities describing how a medium responds to
  electric and magnetic fields; they can be direction-dependent tensors.
\item
  \textbf{Nonreciprocity} --- Response that breaks the usual interchange
  symmetry between source and receiver; here it provides the required
  time-reversal-breaking mechanism.
\item
  \textbf{Photonic band} --- An allowed family of electromagnetic wave
  frequencies across crystal momenta.
\item
  \textbf{Transverse constraint} --- The condition selecting physical
  wave fields consistent with the relevant divergence equation.
\item
  \textbf{Chern number} --- An integer characterizing the topology of an
  isolated group of wave bands.
\item
  \textbf{Interface mode} --- A wave concentrated near the boundary
  between two media.
\item
  \textbf{Lossless and nondispersive} --- Idealizations in which energy
  is not absorbed and the constitutive response is taken independent of
  frequency.
\end{itemize}

\subsubsection{Why this matters}

Controlling where and in which direction waves travel is a basic aim of
optical engineering. Topology offers one mechanism for robust interface
behavior under appropriate conditions. A certified Maxwell calculation
would also expose whether an apparent effect is caused by numerical
artifacts or by assumptions a real material cannot satisfy.

\subsubsection{What progress would mean}

The first meaningful achievement would be a validated ideal medium with
a proved gap and topological index. It could inform later waveguide
designs, but performance in a lossy, frequency-dependent device would
require a separate physical model.

\subsection{Core question}

Can a specified lossless periodic electromagnetic model support a
rigorously isolated band group with nonzero Chern number, and can its
interface modes be validated?

\subsection{Scope and assumptions}

Choose a two-dimensional periodic, lossless, frequency-independent
constitutive model with Hermitian positive \ensuremath{\epsilon}(r),
\ensuremath{\mu}(r), specified boundary conditions, and a mechanism
breaking physical time reversal. Treat the constitutive tensors as
mathematical inputs. Real-frequency dispersion, loss, and fabrication
feasibility are extensions requiring additional physical data.

Start by validating a reciprocal scalar dielectric benchmark with zero
total Chern number for a time-reversal-invariant band group; then
introduce an explicitly nonreciprocal tensor.

\subsection{Mathematical target}

In a formulation supporting the chosen model, solve

\[(\nabla +ik)\times \mu ^{-1}(r)(\nabla +ik)\times u=\omega ^{2}\epsilon (r)u,\]

in units with c=1, subject to the appropriate transverse constraint. Use
the \ensuremath{\epsilon}-weighted inner product or an equivalent
Hermitian formulation. Specify an isolated positive-frequency band
group; photons do not require an electron-like filling assumption. Prove
a frequency gap over the entire Brillouin zone and certify the projector
invariant in the correct metric.

\subsection{Required outputs and proof obligations}

Return ideal geometry and constitutive tensors, solver discretization,
eigenvalue enclosures, spurious-mode controls, gap and topology
certificates. Interface calculations must state both media and
termination. Report robustness only for an explicit perturbation class
and bound that preserves the required gap/symmetry.

\subsection{Validation and rejection controls}

Check homogeneous-medium dispersion, mesh/basis convergence,
Hermiticity, and transverse modes. Separate Galerkin error from k-space
interpolation error. Verify reciprocal controls cannot acquire a net
Chern number from numerical artifacts. Compare an interface spectrum
with the bulk-index difference without claiming immunity to arbitrary
disorder.

\subsection{Milestones and research extension}

Baseline: validated reciprocal eigenproblem. Strong: one certified ideal
nonreciprocal topological design. Research extension: a realistic
dispersive constitutive model or gap optimization within a bounded
geometry family. Fabrication drawings alone do not certify material
behavior.

Milestones are acceptance gates, not calendar promises. Before a run,
freeze the concrete instance/domain, tolerance, compute budget, and
permitted dependencies. A bounded search that fails to find a result
must preserve an unresolved outcome.

\subsection{Starting instruction}

Work on the baseline above. First make every input and convention
explicit, then validate the reference case. Build the checker alongside
the calculation. Label numerical evidence, finite-model certificates,
and continuum theorems separately; return the strongest justified result
and identify the exact unmet proof obligation.

\subsection{Primary references and their role}

\begin{itemize}
\tightlist
\item
  \href{https://journals.aps.org/prl/abstract/10.1103/PhysRevLett.100.013904}{Haldane
  and Raghu, Directional optical waveguides with broken time-reversal
  symmetry} --- Nonreciprocal media as the mechanism for photonic
  quantum-Hall-like edge modes.
\end{itemize}

References support definitions or limitations; the scope choices and
proposed milestones are editorial recommendations, not claims that these
results have already been obtained.

\subsection{Displayed formulation and claim boundary}

\href{https://github.com/rick-stevens-ai/Pure-Thought/blob/main/PRDs/11-Photonic-Topological-Crystals.md}{Canonical
specification}.

\[(\nabla+ik)\times\mu^{-1}(r)(\nabla+ik)\times u
=\omega^2\epsilon(r)u,\qquad
\langle u,v\rangle_\epsilon=\int_{\mathrm{cell}}u^*\cdot\epsilon v\,dr.\]

Use positive Hermitian constitutive tensors, the transverse constraint,
units \(c=1\), and an isolated positive-frequency band group. Certify
its gap and projector invariant in the weighted metric. Nonzero total
Chern number requires breaking the relevant physical time-reversal
symmetry; scalar reciprocal benchmarks provide rejection controls.

\begin{center}\rule{0.5\linewidth}{0.5pt}\end{center}

\section{12. Rigidity and boundary modes of periodic Maxwell frames}

\textbf{Revised specification --- 2026-09-11.}
\href{https://github.com/rick-stevens-ai/Pure-Thought/blob/main/reviews/12-Topological-Mechanical-Metamaterials.md}{Detailed
critique} \ensuremath{\cdot}
\href{https://github.com/rick-stevens-ai/Pure-Thought/blob/main/archive/original-PRDs/12-Topological-Mechanical-Metamaterials.md}{Original
description} \ensuremath{\cdot}
\href{https://github.com/rick-stevens-ai/Pure-Thought/blob/main/reviews/EVIDENCE-STANDARD.md}{Evidence
standard}

\subsection{Plain-language guide}

\subsubsection{The problem in everyday terms}

A framework of bars or springs can be rigid in one place and easy to
deform in another. Sometimes a carefully arranged repeating framework
has soft motions concentrated at its boundary. The question is whether
we can predict and prove those motions from the exact geometry, rather
than simply observing them in a simulation. Counting connections helps,
but the direction of each connection matters too.

\subsubsection{Key terms}

\begin{itemize}
\item
  \textbf{Frame} --- A network of sites joined by bars or springs, used
  as an ideal mechanical model.
\item
  \textbf{Degree of freedom} --- An independently variable coordinate of
  the system's motion.
\item
  \textbf{Isostatic or Maxwell frame} --- A frame balanced between
  degrees of freedom and constraints in the relevant counting sense;
  details of zero modes and self stress still matter.
\item
  \textbf{Zero mode} --- An infinitesimal displacement requiring no
  restoring energy in the linearized model.
\item
  \textbf{State of self stress} --- A pattern of internal bond tensions
  that balances forces without external loads.
\item
  \textbf{Compatibility matrix C} --- The linear map from small site
  displacements to bond extensions.
\item
  \textbf{Stiffness matrix K} --- The operator relating displacements to
  restoring forces in the linear approximation.
\item
  \textbf{Winding number} --- An integer counting how a complex quantity
  winds around zero along a closed path; here it helps characterize the
  periodic frame.
\item
  \textbf{Termination} --- The precise way a periodic structure is cut
  to make a boundary.
\end{itemize}

\subsubsection{Why this matters}

Geometry-controlled softness could help explain and design structures
with directional compliance, boundary motion or localized mechanical
response. The scientific value is understanding which behaviors follow
from a robust index and which are accidental features of one geometry.
That distinction matters before translating an ideal spring model into a
physical structure.

\subsubsection{What progress would mean}

A useful result would give a frame, an exact count of its allowed soft
motions and a proof of where they concentrate. Material failure,
nonlinear deformation and manufacturing tolerance would remain
additional engineering questions.

\subsection{Core question}

For a geometrically specified isostatic frame and boundary termination,
can the zero-mode count and localization be proved from its
compatibility matrix?

\subsection{Scope and assumptions}

Use unstressed central-force springs with exact rational or algebraic
site coordinates, positive spring constants, and explicit periodic bond
offsets. Fix the unit-cell gauge, dimension, masses, and boundary
termination. Begin with a small frame and a known periodic Maxwell
lattice.

Connectivity alone is insufficient: bond directions determine the
compatibility matrix. Distinguish rigid-body motions, floppy modes, and
states of self stress. Claims concern linearized ideal mechanics.

\subsection{Mathematical target}

Let C map displacements to bond extensions and Q=C\ensuremath{\dagger}.
For a finite free frame,

\[n_0- n_{ss}=dN- N_b,\qquad K=C^{\dagger }\operatorname{diag}(k_b)C,\]

with dynamical matrix M\ensuremath{{}^{-1/2}}KM\ensuremath{{}^{-1/2}}.
For a square periodic C(k), define winding
(2\ensuremath{\pi}i)\ensuremath{{}^{-1}}\ensuremath{\oint}d log det C(k)
only along loops where the determinant is nonzero. Account for acoustic
translation zeros and the local boundary count when translating bulk
winding into an edge index.

\subsection{Required outputs and proof obligations}

Export positions, bonds, C, exact rank/kernel calculations, self-stress
basis, winding certificate, and a boundary-specific mode count. For
localization, provide a transfer-matrix/root-modulus bound separated
from the unit circle. Specify perturbations preserving the nonvanishing
condition and positivity of spring constants.

\subsection{Validation and rejection controls}

Check one elementary underconstrained and one overconstrained frame
against the index equation. Verify K is positive semidefinite and
translations are zero modes where appropriate. Compare two terminations
of the same bulk. Reject winding calculations crossing a zero of det C.

\subsection{Milestones and research extension}

Baseline: exact Maxwell--Calladine counts on finite frames. Strong: one
periodic frame with certified winding and termination-dependent edge
localization. Research extension: a bounded family with robust mode
counts. Elastic beams, prestress, nonlinear stability, and manufacturing
tolerance require separate models.

Milestones are acceptance gates, not calendar promises. Before a run,
freeze the concrete instance/domain, tolerance, compute budget, and
permitted dependencies. A bounded search that fails to find a result
must preserve an unresolved outcome.

\subsection{Starting instruction}

Work on the baseline above. First make every input and convention
explicit, then validate the reference case. Build the checker alongside
the calculation. Label numerical evidence, finite-model certificates,
and continuum theorems separately; return the strongest justified result
and identify the exact unmet proof obligation.

\subsection{Primary references and their role}

\begin{itemize}
\tightlist
\item
  \href{https://arxiv.org/abs/1308.0554}{Kane and Lubensky, Topological
  Boundary Modes in Isostatic Lattices} --- Compatibility matrices,
  topological indices, and boundary modes.
\end{itemize}

References support definitions or limitations; the scope choices and
proposed milestones are editorial recommendations, not claims that these
results have already been obtained.

\subsection{Displayed formulation and claim boundary}

\href{https://github.com/rick-stevens-ai/Pure-Thought/blob/main/PRDs/12-Topological-Mechanical-Metamaterials.md}{Canonical
specification}.

\[Q=C^\dagger,\qquad n_0-n_{ss}=dN-N_b,\qquad
K=C^\dagger\operatorname{diag}(k_b)C,\quad D=M^{-1/2}KM^{-1/2}.\]

\[\nu_\gamma=\frac1{2\pi i}\oint_\gamma d\log\det C(k),
\qquad \det C(k)\ne0\text{ on }\gamma.\]

The first count includes rigid-body motions and self stress. Supply
actual bond directions, unit-cell gauge and termination. Acoustic zeros
and the local boundary count must be included before identifying edge
modes from winding.

\begin{center}\rule{0.5\linewidth}{0.5pt}\end{center}

\section{13. Higher-order topology with explicit boundary and symmetry
assumptions}

\textbf{Revised specification --- 2026-09-11.}
\href{https://github.com/rick-stevens-ai/Pure-Thought/blob/main/reviews/13-Higher-Order-Topological-Insulators.md}{Detailed
critique} \ensuremath{\cdot}
\href{https://github.com/rick-stevens-ai/Pure-Thought/blob/main/archive/original-PRDs/13-Higher-Order-Topological-Insulators.md}{Original
description} \ensuremath{\cdot}
\href{https://github.com/rick-stevens-ai/Pure-Thought/blob/main/reviews/EVIDENCE-STANDARD.md}{Evidence
standard}

\subsection{Plain-language guide}

\subsubsection{The problem in everyday terms}

A topological material can have interesting boundary behavior even when
its broad edges do not conduct. In certain models, charge accumulates in
a quantized way near corners instead. This problem asks how to establish
that behavior carefully: what property belongs to the bulk, what depends
on the way the sample is cut, and what extra symmetry is needed to place
a corner state at exactly zero energy?

\subsubsection{Key terms}

\begin{itemize}
\item
  \textbf{Higher-order topology} --- Topology associated with
  lower-dimensional boundaries, such as corners of a two-dimensional
  sample or hinges of a three-dimensional sample.
\item
  \textbf{Quadrupole} --- A pattern of electric charge organization
  beyond total charge and ordinary dipole polarization; the selected
  model supports a quantized version under particular assumptions.
\item
  \textbf{Bulk} --- The interior of a system, far from boundaries.
\item
  \textbf{Wilson loop} --- A matrix describing how a chosen set of
  quantum states changes when transported around a closed path in
  momentum space.
\item
  \textbf{Wannier sector} --- A selected group of Wilson-loop
  eigenstates, separated from the others by a gap in that loop's
  spectrum.
\item
  \textbf{Nested polarization} --- A polarization calculated within such
  a Wannier sector, using a second transport calculation.
\item
  \textbf{Corner excess charge} --- Charge near a corner measured
  relative to a specified background, after accounting for other
  contributions.
\item
  \textbf{Chiral symmetry} --- An additional spectral symmetry that can
  relate positive and negative energies; it can matter for pinning
  special states to zero.
\item
  \textbf{Termination} --- The actual boundary cut and local boundary
  structure of the sample.
\end{itemize}

\subsubsection{Why this matters}

This problem broadens our understanding of how a material's interior
organizes behavior at its boundaries. It also corrects a common
interpretive shortcut: a quantized corner charge and a zero-energy
corner state are not the same observation. Reliable distinctions are
needed before such features can be used as design principles.

\subsubsection{What progress would mean}

Progress would produce a model with verified gaps, a well-defined bulk
invariant and a carefully measured boundary consequence. It would
establish a controlled example, rather than guarantee protected corner
states for every sample shape or boundary.

\subsection{Core question}

For a fixed quadrupole-insulator family, can one certify the bulk and
Wannier gaps, a quantized nested polarization, and the corresponding
boundary charge under a specified termination?

\subsection{Scope and assumptions}

Start with a four-band Benalcazar--Bernevig--Hughes-type square-lattice
model at half filling. Supply all hopping signs, orbital positions,
mirror representations, and a rectangular termination. State whether
chiral symmetry is imposed. Restrict parameters to a compact region away
from bulk and Wannier-gap closings.

Separate fractional corner charge from an exactly zero-energy corner
eigenstate. The latter requires additional spectral symmetry and
boundary assumptions.

\subsection{Mathematical target}

Let P(k) be the rank-two occupied projector. Build the path-ordered
Wilson loop W\_x(k\_y) from occupied-state overlaps. If its spectrum has
a certified separation into Wannier sectors, form a smooth sector
projector and its Berry holonomy along y. Define the nested polarization
from that holonomy modulo one.

Do not replace a nested Wilson loop by the integral of the first
Wilson-loop eigenphases. Define corner excess charge relative to an
explicit ionic/reference background and a stated spatial window, with
finite-size and edge contributions controlled.

\subsection{Required outputs and proof obligations}

Return H, symmetry operators and identities, bulk and Wannier-gap
enclosures, nested invariant, boundary charge definition, and
finite-size bounds. Any robustness statement must preserve the
protecting symmetry and relevant gaps. State whether conclusions are
intrinsic bulk statements or termination-dependent observables.

\subsection{Validation and rejection controls}

Recover topological and trivial limits of the selected model. Verify
zero total Chern number where required while the nested invariant is
nontrivial. Break a protecting symmetry as a control; the checker must
stop claiming quantization. Compare two terminations and explain any
boundary difference.

\subsection{Milestones and research extension}

Baseline: reproduce one quadrupole benchmark with valid Wilson-loop
construction. Strong: certify its invariant and boundary charge on a
parameter box. Research extension: classify a bounded family or one
additional symmetry setting; no universal HOTI classification is
implied.

Milestones are acceptance gates, not calendar promises. Before a run,
freeze the concrete instance/domain, tolerance, compute budget, and
permitted dependencies. A bounded search that fails to find a result
must preserve an unresolved outcome.

\subsection{Starting instruction}

Work on the baseline above. First make every input and convention
explicit, then validate the reference case. Build the checker alongside
the calculation. Label numerical evidence, finite-model certificates,
and continuum theorems separately; return the strongest justified result
and identify the exact unmet proof obligation.

\subsection{Primary references and their role}

\begin{itemize}
\tightlist
\item
  \href{https://arxiv.org/abs/1611.07987}{Benalcazar, Bernevig and
  Hughes, Quantized Electric Multipole Insulators} --- Nested Wilson
  loops and assumptions for quantized multipole response.
\end{itemize}

References support definitions or limitations; the scope choices and
proposed milestones are editorial recommendations, not claims that these
results have already been obtained.

\subsection{Displayed formulation and claim boundary}

\href{https://github.com/rick-stevens-ai/Pure-Thought/blob/main/PRDs/13-Higher-Order-Topological-Insulators.md}{Canonical
specification}.

\[W_x(k_y)=\mathcal P\exp\!\left(i\oint\mathcal A_x(k_x,k_y)\,dk_x\right),
\qquad \mathcal A_{x,ab}=i\langle u_a|\partial_{k_x}u_b\rangle.\]

\[\widetilde W_y=\mathcal P\exp\!\left(i\oint\widetilde{\mathcal A}_y\,dk_y\right),
\qquad p_y^{\mathrm{sector}}=\frac1{2\pi}\arg\det\widetilde W_y\pmod1.\]

Use consistent holonomy orientation in an overlap implementation. The
second connection belongs to a smoothly transported, spectrally isolated
Wannier sector, not to the first loop's eigenphases alone. Prove both
bulk and Wannier gaps and the relevant quantizing symmetry. Define
corner charge with a reference background and termination; zero-energy
corner states require additional assumptions.

\begin{center}\rule{0.5\linewidth}{0.5pt}\end{center}

\section{14. Certified Weyl nodes and slice topology in a lattice model}

\textbf{Revised specification --- 2026-09-11.}
\href{https://github.com/rick-stevens-ai/Pure-Thought/blob/main/reviews/14-Topological-Semimetals-Weyl-Dirac.md}{Detailed
critique} \ensuremath{\cdot}
\href{https://github.com/rick-stevens-ai/Pure-Thought/blob/main/archive/original-PRDs/14-Topological-Semimetals-Weyl-Dirac.md}{Original
description} \ensuremath{\cdot}
\href{https://github.com/rick-stevens-ai/Pure-Thought/blob/main/reviews/EVIDENCE-STANDARD.md}{Evidence
standard}

\subsection{Plain-language guide}

\subsubsection{The problem in everyday terms}

Energy bands in a three-dimensional crystal can meet at isolated points
instead of remaining separated everywhere. Near certain crossings, the
equations resemble those of a Weyl particle. Each crossing carries a
handedness-like topological label. This problem asks whether we can find
every such crossing in one lattice model, prove their labels, and
connect them to the behavior of two-dimensional momentum slices.

\subsubsection{Key terms}

\begin{itemize}
\item
  \textbf{Semimetal} --- A system in which the occupied and unoccupied
  electronic states are not separated by a full insulating gap; the
  selected model has isolated band crossings.
\item
  \textbf{Weyl node} --- A simple crossing of two bands with a
  characteristic linear structure and nonzero topological charge.
\item
  \textbf{Chirality} --- The signed handedness of a simple Weyl
  crossing, defined here from its velocity matrix.
\item
  \textbf{Velocity matrix} --- The derivatives describing how the
  Hamiltonian changes in each momentum direction near a node.
\item
  \textbf{Berry flux} --- The integrated Berry curvature through a
  closed surface around a crossing.
\item
  \textbf{Momentum slice} --- A two-dimensional section of the
  three-dimensional Brillouin zone, obtained by holding one momentum
  coordinate fixed.
\item
  \textbf{Fermi arc} --- An open surface-state contour whose relation to
  bulk nodes depends on the chosen surface and energy.
\item
  \textbf{Root isolation} --- A mathematical procedure proving that a
  small region contains a solution, often uniquely.
\item
  \textbf{Completeness} --- A proof that no additional nodes have been
  missed anywhere in the claimed domain.
\end{itemize}

\subsubsection{Why this matters}

Weyl nodes connect a local band crossing to a global topological
constraint. They help explain unusual electronic structures and surface
states, and provide controlled settings for studying the relation
between particle-like excitations and lattice physics. Missing even a
pair of nodes can change the interpretation of a phase diagram.

\subsubsection{What progress would mean}

A certified catalog would establish all crossings and slice invariants
of the selected model. It would give a reliable starting point for
surface and response calculations, without pretending that bulk nodes
uniquely determine every surface arc.

\subsection{Core question}

Can all nodes in a specified two-band three-dimensional Hamiltonian be
isolated, assigned chirality, and related to the Chern numbers of gapped
momentum slices?

\subsection{Scope and assumptions}

Begin with
\(H(k)=\sin k_x\,\sigma_x+\sin k_y\,\sigma_y+(2-\cos k_x-\cos k_y-\cos k_z)\sigma_z\),
k\ensuremath{\in}T\ensuremath{{}^{3}}, lower band occupied away from
nodes. Fix orientation and Pauli-matrix conventions. This exact
benchmark has nodes at (0,0,\ensuremath{\pm}\ensuremath{\pi}/2).

Subsequent searches must fix orbitals, hopping support, parameter box,
and symmetry representation. Dirac points and nodal lines are separate
extensions because their protection criteria differ.

\subsection{Mathematical target}

For H=d\ensuremath{{}_{0}}I+d\ensuremath{\cdot}\ensuremath{\sigma},
nodes solve d(k)=0. A simple node has invertible velocity matrix
\(V_{ij}=\partial d_i/\partial k_j\). Define Weyl chirality
\ensuremath{\chi}=sgn det V, and explicitly relate this to the
lower-band Chern flux under the chosen convention. Prove the complement
of all isolating neighborhoods has \textbar d(k)\textbar\textgreater0.

For fixed k\_z away from nodes, certify the two-dimensional occupied
Chern number. Its jumps track node charge. A slab calculation requires a
chosen surface normal and termination; arc details are not determined
solely by bulk node positions.

\subsection{Required outputs and proof obligations}

Return all root enclosures, uniqueness/completeness proof, Jacobian
signs, slice Chern numbers, and optional slab spectrum with error
control. Nodes need not have closed-form coordinates: certified
isolating boxes are acceptable. Robustness is scoped to perturbations
that prevent annihilation with opposite total charge.

\subsection{Validation and rejection controls}

Verify the two benchmark roots analytically and opposite determinant
signs. Check total charge on the periodic Brillouin zone vanishes.
Confirm slice invariants only on gapped slices. Test a perturbation
moving nodes and a parameter change permitting pair annihilation.

\subsection{Milestones and research extension}

Baseline: complete benchmark node and slice audit. Strong: certified
parameter-dependent node motion or one bounded model family. Research
extension: a symmetry-protected Dirac model with explicit little-group
constraints, or nodal-link invariants with their own definitions.

Milestones are acceptance gates, not calendar promises. Before a run,
freeze the concrete instance/domain, tolerance, compute budget, and
permitted dependencies. A bounded search that fails to find a result
must preserve an unresolved outcome.

\subsection{Starting instruction}

Work on the baseline above. First make every input and convention
explicit, then validate the reference case. Build the checker alongside
the calculation. Label numerical evidence, finite-model certificates,
and continuum theorems separately; return the strongest justified result
and identify the exact unmet proof obligation.

\subsection{Primary references and their role}

\begin{itemize}
\tightlist
\item
  \href{https://arxiv.org/abs/1007.0016}{Wan et al., Electronic
  Structure of Pyrochlore Iridates (published as Topological semimetal
  and Fermi-arc surface states)} --- Reference for Weyl nodes and
  surface arcs; benchmark and proof obligations here are explicitly
  specified.
\end{itemize}

References support definitions or limitations; the scope choices and
proposed milestones are editorial recommendations, not claims that these
results have already been obtained.

\subsection{Displayed formulation and claim boundary}

\href{https://github.com/rick-stevens-ai/Pure-Thought/blob/main/PRDs/14-Topological-Semimetals-Weyl-Dirac.md}{Canonical
specification}.

\[d(k)=(\sin k_x,\sin k_y,2-\cos k_x-\cos k_y-\cos k_z),
\quad d(k_*)=0,\]

\[k_*=(0,0,\pm\pi/2),\qquad
V_{ij}=\partial_{k_j}d_i(k_*),\qquad \chi=\operatorname{sgn}\det V.\]

The two determinants are \(+1\) and \(-1\) in the stated coordinate
order. Relate this convention explicitly to the occupied-band flux.
Certify absence of other nodes on the complement and the topology of
gapped slices; surface arcs additionally depend on the chosen boundary.

\begin{center}\rule{0.5\linewidth}{0.5pt}\end{center}

\section{15. Symmetry-indicator algebra and the limits of band
diagnosis}

\textbf{Revised specification --- 2026-09-11.}
\href{https://github.com/rick-stevens-ai/Pure-Thought/blob/main/reviews/15-Topological-Quantum-Chemistry.md}{Detailed
critique} \ensuremath{\cdot}
\href{https://github.com/rick-stevens-ai/Pure-Thought/blob/main/archive/original-PRDs/15-Topological-Quantum-Chemistry.md}{Original
description} \ensuremath{\cdot}
\href{https://github.com/rick-stevens-ai/Pure-Thought/blob/main/reviews/EVIDENCE-STANDARD.md}{Evidence
standard}

\subsection{Plain-language guide}

\subsubsection{The problem in everyday terms}

Symmetry can act like a fingerprint for quantum bands: it provides a
compact set of labels that can sometimes reveal nontrivial topology. But
different objects can share a fingerprint. This problem asks exactly
what can be inferred from symmetry labels, how to compute that inference
reliably, and when the correct answer is that more information about the
quantum states is needed.

\subsubsection{Key terms}

\begin{itemize}
\item
  \textbf{Crystal symmetry} --- A translation, rotation, reflection or
  related operation preserving the lattice structure and the stated
  physical setting.
\item
  \textbf{Representation} --- Matrices describing how quantum states
  transform under a symmetry.
\item
  \textbf{Irreducible representation (irrep)} --- A basic representation
  that cannot be decomposed into smaller independent representations of
  the same group.
\item
  \textbf{Compatibility relations} --- Rules requiring symmetry labels
  at different momenta to fit into continuous bands.
\item
  \textbf{Atomic band representation} --- Band data generated from
  suitably localized, symmetry-compatible orbitals.
\item
  \textbf{Wannier function} --- A localized real-space representation of
  band states; the existence of an appropriate symmetry-respecting basis
  is central to atomicity.
\item
  \textbf{Symmetry indicator} --- An algebraic label detecting certain
  obstructions to atomic band data, but not every possible topological
  distinction.
\item
  \textbf{Smith normal form} --- An exact integer-matrix decomposition
  used to calculate the quotient of symmetry-data lattices.
\item
  \textbf{Fragile topology} --- An obstruction that can disappear when
  suitable trivial bands are added, under the chosen definition.
\end{itemize}

\subsubsection{Why this matters}

Symmetry labels are often much simpler to work with than the full
wavefunctions across momentum space. Knowing what those labels do and do
not prove can make topological diagnosis more efficient and prevent
false declarations that a system is ordinary. The distinction is
essential for turning a classification table into a trustworthy
scientific tool.

\subsubsection{What progress would mean}

A useful result would be an exact indicator calculation for one symmetry
setting, together with examples showing its reach and blind spots.
Stronger conclusions about a particular Hamiltonian would use additional
information, such as Berry phases or a Wannier construction.

\subsection{Core question}

For one fixed symmetry setting, can compatibility relations and atomic
band representations be generated exactly, and what topology can their
quotient actually detect?

\subsection{Scope and assumptions}

Choose one small symmetry setting, initially a spinless two-dimensional
inversion-symmetric lattice without time reversal. Specify the
space-group generators, orbital/site representations, and high-symmetry
momenta. Derive all input representations from group operations, with
external tables used only as checks.

The result is an algebraic classification of symmetry data. A zero
indicator is not a certificate of a trivial occupied bundle.

\subsection{Mathematical target}

Let B\ensuremath{\subset}Z\textasciicircum{}r be the lattice of integer
irrep vectors satisfying compatibility constraints, and let
A\ensuremath{\subset}B be the integer lattice generated by atomic
band-representation vectors. Compute X=B/A by Smith normal form,
including basis-change matrices.

For a physical nonnegative irrep vector b, distinguish integer-lattice
membership from nonnegative decomposition b=A\_mat n,
n\ensuremath{\ge}0. Nonnegative decomposition establishes compatibility
with atomic symmetry data, not by itself a Wannier construction for a
supplied Hamiltonian. Conversely, an actual symmetry-preserving
localized Wannier basis is constructive atomic evidence.

\subsection{Required outputs and proof obligations}

Export group conventions, irreps, compatibility matrix, atomic
generators, Smith decomposition, indicator mapping, and witnesses for
any integer/nonnegative decomposition. Report ``indicator trivial;
topology unresolved'' when symmetry data are insufficient. Use
Berry/Wilson-loop information when making a stronger claim about a
specific H.

\subsection{Validation and rejection controls}

Verify all representation identities and dimension counts. Reconstruct
the Smith identity U A\_mat V=D in the chosen B basis, with U,V
unimodular. Test known atomic data, incompatible data, and models
sharing symmetry labels but differing in an additional invariant. Keep
stable and fragile notions distinct.

\subsection{Milestones and research extension}

Baseline: exact quotient and decompositions for one setting. Strong:
connect the diagnosis to explicit Hamiltonians and identify its kernel.
Research extension: another symmetry group or a rigorous fragile
obstruction with a definition of allowed added bands.

Milestones are acceptance gates, not calendar promises. Before a run,
freeze the concrete instance/domain, tolerance, compute budget, and
permitted dependencies. A bounded search that fails to find a result
must preserve an unresolved outcome.

\subsection{Starting instruction}

Work on the baseline above. First make every input and convention
explicit, then validate the reference case. Build the checker alongside
the calculation. Label numerical evidence, finite-model certificates,
and continuum theorems separately; return the strongest justified result
and identify the exact unmet proof obligation.

\subsection{Primary references and their role}

\begin{itemize}
\tightlist
\item
  \href{https://arxiv.org/abs/1703.00911}{Po, Vishwanath and Watanabe,
  Symmetry-based Indicators of Band Topology} --- Quotient of compatible
  symmetry data by atomic data; symmetry-based diagnosis.
\end{itemize}

References support definitions or limitations; the scope choices and
proposed milestones are editorial recommendations, not claims that these
results have already been obtained.

\subsection{Displayed formulation and claim boundary}

\href{https://github.com/rick-stevens-ai/Pure-Thought/blob/main/PRDs/15-Topological-Quantum-Chemistry.md}{Canonical
specification}.

\[A\subset B\subset\mathbb Z^r,\qquad X=B/A,\qquad
UA_{\mathrm{mat}}V=D,\quad U,V\text{ unimodular}.\]

\[b=A_{\mathrm{mat}}n,\quad n\in\mathbb Z^s
\qquad\text{versus}\qquad b=A_{\mathrm{mat}}n,\quad n\in\mathbb Z_{\ge0}^s.\]

Use a declared basis for the compatibility lattice \(B\) in the Smith
decomposition. Integer membership and nonnegative atomic symmetry-data
decomposition are distinct tests. Neither a zero indicator nor matching
atomic labels alone constructs localized Wannier functions for a given
Hamiltonian.

\begin{center}\rule{0.5\linewidth}{0.5pt}\end{center}

\section{16. Certified energy brackets from reduced-density-matrix
relaxations}

\textbf{Revised specification --- 2026-09-11.}
\href{https://github.com/rick-stevens-ai/Pure-Thought/blob/main/reviews/16-N-Representability-2RDM.md}{Detailed
critique} \ensuremath{\cdot}
\href{https://github.com/rick-stevens-ai/Pure-Thought/blob/main/archive/original-PRDs/16-N-Representability-2RDM.md}{Original
description} \ensuremath{\cdot}
\href{https://github.com/rick-stevens-ai/Pure-Thought/blob/main/reviews/EVIDENCE-STANDARD.md}{Evidence
standard}

\subsection{Plain-language guide}

\subsubsection{The problem in everyday terms}

Describing every possible configuration of many electrons becomes
overwhelmingly expensive as the system grows. Yet the energy of a system
with one- and two-particle interactions depends only on certain
one-electron and two-electron summaries. The challenge is that not every
plausible summary can come from a real many-electron state. We use
necessary consistency tests to calculate a guaranteed lower energy limit
and an actual trial state to obtain an upper limit.

\subsubsection{Key terms}

\begin{itemize}
\item
  \textbf{Fermion} --- A particle, such as an electron, whose quantum
  states obey antisymmetry under exchange.
\item
  \textbf{Reduced density matrix (RDM)} --- A summary of selected
  particles' quantum statistics obtained from a larger system's state.
\item
  \textbf{2-RDM} --- The two-particle reduced density matrix, sufficient
  with the appropriate one-particle data to compute the chosen two-body
  Hamiltonian's energy.
\item
  \textbf{N-representability} --- The requirement that the reduced data
  arise from some valid state of N fermions.
\item
  \textbf{P, Q and G conditions} --- Positivity tests involving particle
  pairs, hole pairs and particle--hole combinations; they are necessary
  representability conditions.
\item
  \textbf{Hole} --- An unoccupied orbital state, treated as an absence
  of a particle relative to the finite orbital description.
\item
  \textbf{Semidefinite program (SDP)} --- An optimization problem with
  linear conditions and positive-semidefinite matrix constraints.
\item
  \textbf{Ground-state energy} --- The lowest allowed energy of the
  chosen quantum system.
\item
  \textbf{Energy bracket} --- A pair of verified lower and upper numbers
  containing the true ground-state energy.
\end{itemize}

\subsubsection{Why this matters}

Reliable energy calculations underpin the theoretical study of bonding
and interacting electrons. A bracket tells us how much uncertainty
remains, rather than merely reporting a convincing decimal number. The
approach may extract useful information without storing a complete
many-electron wavefunction, although general representability remains
computationally difficult.

\subsubsection{What progress would mean}

Progress would mean tighter verified brackets at an affordable cost for
a specified finite model. That would show how much the extra constraints
help; extending the conclusion to an actual molecule would also require
controlling orbital-basis and Hamiltonian errors.

\subsection{Core question}

For a specified finite fermionic Hamiltonian, how tight an independently
certified ground-state energy interval can P,Q,G and selected
higher-order RDM constraints produce?

\subsection{Scope and assumptions}

Fix N\ensuremath{\ge}2 fermions in M spin orbitals with an exact
rational/algebraic, number-conserving one- and two-body Hamiltonian.
Start with a small Hubbard-type system where exact diagonalization in
the same N-particle sector is feasible. Molecular integrals are a later
input with an explicit basis and numerical error enclosure.

Use ensemble N-representability. The task is to certify bounds within a
finite orbital model, not solve general electronic structure in
polynomial time.

\subsection{Mathematical target}

Set \(\gamma_{ij}=\langle a_i^\dagger a_j\rangle\) and
\(\Gamma_{ij,kl}=\langle a_i^\dagger a_j^\dagger a_l a_k\rangle\).
Require Hermiticity, fermionic antisymmetry, Tr \ensuremath{\gamma}=N,
\(\sum_{ij}\Gamma_{ij,ij}=N(N-1)\), and
\(\sum_j\Gamma_{ij,kj}=(N-1)\gamma_{ik}\). Define Q and G using two-hole
and particle-hole operator moments and derive their affine forms from
anticommutation identities.

Minimizing the linear energy over a necessary-condition outer relaxation
gives E\_relax\ensuremath{\le}E\ensuremath{{}_{0}}. A verified dual
feasible point gives L\ensuremath{\le}E\_relax; an explicit normalized
N-fermion state gives U\ensuremath{\ge}E\ensuremath{{}_{0}}. Report
L\ensuremath{\le}E\ensuremath{{}_{0}}\ensuremath{\le}U, not an
unsupported accuracy estimate from \ensuremath{\Gamma} alone.

\subsection{Required outputs and proof obligations}

Export the Hamiltonian and indexing convention, moment constraints,
certified dual objective L, state witness for U, and gap
U\ensuremath{-}L. Separate solver error, relaxation error, and
finite-basis/model error. Report any property bounds via separately
optimized observables; a relaxed \ensuremath{\Gamma} need not be
physical.

\subsection{Validation and rejection controls}

Compare with exact diagonalization in identical conventions. Check
contractions using explicit Slater determinants and correlated states.
Adding valid necessary constraints must not lower the exact relaxed
minimum. Verify dual PSD matrices and equality residual corrections
rigorously, rather than interpreting the raw solver objective as a
bound.

\subsection{Milestones and research extension}

Baseline: valid L,U brackets for small fixed models. Strong: quantify
improvement from G or T constraints across a specified instance set.
Research extension: a new valid constraint with demonstrated tightening,
or a theorem on a tractable family; novelty needs comparison with
existing hierarchies.

Milestones are acceptance gates, not calendar promises. Before a run,
freeze the concrete instance/domain, tolerance, compute budget, and
permitted dependencies. A bounded search that fails to find a result
must preserve an unresolved outcome.

\subsection{Starting instruction}

Work on the baseline above. First make every input and convention
explicit, then validate the reference case. Build the checker alongside
the calculation. Label numerical evidence, finite-model certificates,
and continuum theorems separately; return the strongest justified result
and identify the exact unmet proof obligation.

\subsection{Primary references and their role}

\begin{itemize}
\tightlist
\item
  \href{https://arxiv.org/abs/quant-ph/0609125}{Liu, Christandl and
  Verstraete, N-representability is QMA-complete} --- Complexity barrier
  to general exact representability.
\end{itemize}

References support definitions or limitations; the scope choices and
proposed milestones are editorial recommendations, not claims that these
results have already been obtained.

\subsection{Displayed formulation and claim boundary}

\href{https://github.com/rick-stevens-ai/Pure-Thought/blob/main/PRDs/16-N-Representability-2RDM.md}{Canonical
specification}.

\[\gamma_{ij}=\langle a_i^\dagger a_j\rangle,\qquad
\Gamma_{ij,kl}=\langle a_i^\dagger a_j^\dagger a_l a_k\rangle,\]

\[\operatorname{Tr}\gamma=N,\quad
\sum_{ij}\Gamma_{ij,ij}=N(N-1),\quad
\sum_j\Gamma_{ij,kj}=(N-1)\gamma_{ik}.\]

\[L_{\mathrm{dual}}\le E_{\mathrm{relax}}\le E_0\le
\langle\Psi|H|\Psi\rangle=U,\qquad \langle\Psi|\Psi\rangle=1.\]

Derive two-particle, two-hole and particle-hole moment matrices from
anticommutation in this convention. Necessary PSD constraints give an
outer relaxation. Certify dual feasibility for \(L\) and a physical
\(N\)-fermion state for \(U\); a relaxed RDM need not represent any
state.

\begin{center}\rule{0.5\linewidth}{0.5pt}\end{center}

\section{17. Complete enumeration of a precisely defined molecular-graph
class}

\textbf{Revised specification --- 2026-09-11.}
\href{https://github.com/rick-stevens-ai/Pure-Thought/blob/main/reviews/17-Isomer-Enumeration-Molecular-Graphs.md}{Detailed
critique} \ensuremath{\cdot}
\href{https://github.com/rick-stevens-ai/Pure-Thought/blob/main/archive/original-PRDs/17-Isomer-Enumeration-Molecular-Graphs.md}{Original
description} \ensuremath{\cdot}
\href{https://github.com/rick-stevens-ai/Pure-Thought/blob/main/reviews/EVIDENCE-STANDARD.md}{Evidence
standard}

\subsection{Plain-language guide}

\subsubsection{The problem in everyday terms}

A molecular formula tells us how many atoms of each kind are present,
but not how they are connected. The same counts can describe different
structures: butane and isobutane are a familiar example. This problem
asks for every distinct connection pattern allowed by an explicit set of
bonding rules, with a proof that none was counted twice and none was
missed.

\subsubsection{Key terms}

\begin{itemize}
\item
  \textbf{Molecular formula} --- A count of atoms by element, such as
  C\ensuremath{{}_{4}}H\ensuremath{{}_{10}}.
\item
  \textbf{Constitutional isomer} --- A structure with the same molecular
  formula as another but a different pattern of atom connections.
\item
  \textbf{Molecular graph} --- A mathematical representation with atoms
  as vertices and bonds as edges carrying bond-order labels.
\item
  \textbf{Valence grammar} --- The exact bonding rules permitted in this
  enumeration; it deliberately defines a restricted chemical model.
\item
  \textbf{Bond order} --- The bond label distinguishing single, double
  and triple bonds in this graph model.
\item
  \textbf{Graph isomorphism} --- A relabeling of atoms that preserves
  element types and bonds, showing two drawings represent the same
  graph.
\item
  \textbf{Canonical labeling} --- A consistent representative label for
  a graph, used to identify duplicates.
\item
  \textbf{Stereoisomer} --- A distinction in spatial arrangement not
  captured merely by the basic connection graph.
\item
  \textbf{Completeness certificate} --- Evidence that the enumeration
  covers every object in the specified graph class.
\end{itemize}

\subsubsection{Why this matters}

Complete enumeration turns an intuitive picture of ``possible
structures'' into a checkable mathematical set. Such sets are useful for
testing structure-generation methods and exploring how bonding
constraints shape chemical possibilities. The guarantee is especially
valuable when a search is used to support a claim that no alternative
structure exists.

\subsubsection{What progress would mean}

Success would supply exact graph counts and explicit structures under
the chosen rules. A graph's presence in the list would not imply that
the corresponding molecule is stable, easily synthesized or realizable
in a particular three-dimensional shape.

\subsection{Core question}

For a fixed formula and explicit valence grammar, can every connected
constitutional graph be generated exactly once, with a checkable
completeness argument?

\subsection{Scope and assumptions}

Initially allow neutral closed-shell C,H,O graphs with bond orders
1,2,3, no self-bonds, and valences 4,1,2. Hydrogens may be implicit if
reconstructed uniquely from residual valence. Specify whether different
Kekulé graphs are distinct; the baseline uses literal bond-order graphs
and no aromatic equivalence.

Constitutional graph validity is not chemical stability or
synthesizability. Stereoisomers, conformers, ions, radicals, isotopes,
and 3D geometries require added definitions and are not baseline
outputs.

\subsection{Mathematical target}

For labeled heavy atoms set symmetric b\_ij\ensuremath{\in}\{0,1,2,3\},
b\_ii=0, h\_i=v\_i\ensuremath{-}\ensuremath{\Sigma}\_j
b\_ij\ensuremath{\ge}0, and \ensuremath{\Sigma}\_i h\_i=n\_H. Require
connectivity and the requested counts of element labels. Quotient by
permutations preserving element labels and bond orders.

Prove completeness via canonical augmentation with a proved parent rule,
or exhaustive labeled enumeration plus verified orbit reduction. A SAT
route needs the exact encoding and orbit-blocking/symmetry-breaking
justification; the final UNSAT proof must cover all remaining admissible
graphs.

\subsection{Required outputs and proof obligations}

Export canonical graphs, reconstructed hydrogens, counts, labeling
conventions, and generation/coverage proof. Provide a validator
independent of the optimized generator. If a resource limit stops
enumeration, output a partial catalog and coverage status rather than
``all isomers.''

\subsection{Validation and rejection controls}

Recover alkane constitutional counts for
C\ensuremath{{}_{4}}H\ensuremath{{}_{10}} and
C\ensuremath{{}_{5}}H\ensuremath{{}_{12}}, namely 2 and 3. Cross-check
all tiny instances with brute force. Random relabelings must preserve
canonical codes. Verify each graph's formula and weighted valence;
reject disconnected mixtures.

\subsection{Milestones and research extension}

Baseline: complete small-formula catalog with independent counts.
Strong: complete C\ensuremath{{}_{6}}H\ensuremath{{}_{12}}O under the
stated grammar. Research extension: stereochemical orbit enumeration
with meso cases and explicit stereocenter/double-bond conventions; keep
conformational continua separate.

Milestones are acceptance gates, not calendar promises. Before a run,
freeze the concrete instance/domain, tolerance, compute budget, and
permitted dependencies. A bounded search that fails to find a result
must preserve an unresolved outcome.

\subsection{Starting instruction}

Work on the baseline above. First make every input and convention
explicit, then validate the reference case. Build the checker alongside
the calculation. Label numerical evidence, finite-model certificates,
and continuum theorems separately; return the strongest justified result
and identify the exact unmet proof obligation.

\subsection{Primary references and their role}

\begin{itemize}
\tightlist
\item
  \href{https://users.cecs.anu.edu.au/~bdm/papers/orderly.pdf}{McKay,
  Isomorph-free exhaustive generation} --- Canonical generation and
  completeness methodology.
\end{itemize}

References support definitions or limitations; the scope choices and
proposed milestones are editorial recommendations, not claims that these
results have already been obtained.

\subsection{Displayed formulation and claim boundary}

\href{https://github.com/rick-stevens-ai/Pure-Thought/blob/main/PRDs/17-Isomer-Enumeration-Molecular-Graphs.md}{Canonical
specification}.

Let \(\mathcal G\) contain exactly the graphs satisfying the fixed atom,
bond, valence and connectivity grammar. Then

\[\text{catalog}=\mathcal G/\!\cong.\]

Isomorphisms preserve the declared atom and bond labels. Every output
must satisfy the grammar, no two representatives may be equivalent, and
every allowed input must map to a representative. Canonical labeling
proves deduplication only; completeness needs a coverage proof for
generation and pruning. Stereochemistry requires a separately declared
equivalence model.

\begin{center}\rule{0.5\linewidth}{0.5pt}\end{center}

\section{18. Certified optimal-transport bounds for strictly correlated
electrons}

\textbf{Revised specification --- 2026-09-11.}
\href{https://github.com/rick-stevens-ai/Pure-Thought/blob/main/reviews/18-Optimal-Transport-Chemistry.md}{Detailed
critique} \ensuremath{\cdot}
\href{https://github.com/rick-stevens-ai/Pure-Thought/blob/main/archive/original-PRDs/18-Optimal-Transport-Chemistry.md}{Original
description} \ensuremath{\cdot}
\href{https://github.com/rick-stevens-ai/Pure-Thought/blob/main/reviews/EVIDENCE-STANDARD.md}{Evidence
standard}

\subsection{Plain-language guide}

\subsubsection{The problem in everyday terms}

Electrons repel one another, but a prescribed electron density also
constrains where they spend their time. Imagine arranging their joint
positions so they avoid one another as effectively as possible while
preserving that average density. This problem makes that arrangement
question precise and seeks reliable upper and lower bounds on its
repulsion energy. It represents a limiting strong-correlation model, not
the full energy of ordinary electrons.

\subsubsection{Key terms}

\begin{itemize}
\item
  \textbf{Electron density \ensuremath{\rho}} --- The average spatial
  distribution of electrons, integrating to their total number N.
\item
  \textbf{Coulomb cost} --- The repulsion contribution between charged
  particles, proportional here to the inverse of their separation.
\item
  \textbf{Optimal transport} --- Optimization over ways of distributing
  or coupling quantities while satisfying prescribed distribution
  constraints.
\item
  \textbf{Marginal distribution} --- The distribution of one variable
  after the other variables in a joint distribution are ignored.
\item
  \textbf{Coupling \ensuremath{\Gamma}} --- A joint distribution of all
  electron positions consistent with the prescribed individual
  marginals.
\item
  \textbf{Multi-marginal transport} --- Transport optimization with more
  than two prescribed marginals.
\item
  \textbf{Strictly correlated electrons (SCE)} --- A limiting model in
  which electron positions are organized to minimize interaction energy
  at fixed density.
\item
  \textbf{Dual potential} --- An auxiliary function providing a
  certified lower bound when it satisfies the required inequality
  everywhere.
\item
  \textbf{Discretization error} --- The difference introduced by
  replacing a continuous spatial problem with a finite representation.
\end{itemize}

\subsubsection{Why this matters}

Strongly correlated electrons are difficult to approximate because their
motions are strongly linked. This limiting problem gives a
mathematically structured reference for studying that linkage and for
developing density-functional ideas. Verified bounds help distinguish a
true property of the continuum model from an artifact of a spatial grid
or optimization tolerance.

\subsubsection{What progress would mean}

A useful result would enclose the SCE interaction energy for a declared
density. It could inform further electronic-structure theory, while
leaving kinetic energy, finite-correlation effects and actual molecular
reaction paths to additional modeling.

\subsection{Core question}

For a supplied normalized electron density, can a discretized Coulomb
multi-marginal transport problem yield verified bounds, with continuum
error separated from optimization error?

\subsection{Scope and assumptions}

Take an electron density \ensuremath{\rho}\ensuremath{\ge}0 with
\ensuremath{\int}\ensuremath{\rho}=N and marginal
\ensuremath{\mu}=\ensuremath{\rho}/N. Begin with N=2 and a compactly
supported mathematical density or a finite rational discrete model.
Specify treatment of the diagonal Coulomb singularity. A
continuous-density result additionally needs domain-tail and
discretization estimates.

The main target is the strictly correlated interaction functional.
Wasserstein comparison of densities is a separate optional metric task;
it is not a minimum-energy chemical reaction path.

\subsection{Mathematical target}

Define

\[V_{SCE}[\rho ]=\inf_{\Gamma \in\Pi(\mu ,\ldots,\mu )}\int\sum_{i<j}|r_i-r_j|^{-1}\,d\Gamma .\]

\ensuremath{\Gamma} is a probability measure and may be symmetrized; it
is not an antisymmetric probability distribution. The symmetric dual
uses u with \ensuremath{\Sigma}\_i
u(r\_i)\ensuremath{\le}\ensuremath{\Sigma}\_\{i<j\}|r\_i\ensuremath{-}r\_j|\ensuremath{{}^{-1}}
and objective \ensuremath{\int}u\ensuremath{\rho}. A feasible coupling
gives an upper bound and a dual-feasible potential gives a lower bound
for the same problem. Define
W\_\ensuremath{\infty}=V\_SCE\ensuremath{-}U\_H if subtracting Hartree
energy; do not call V\_SCE exact exchange.

\subsection{Required outputs and proof obligations}

Return marginal conventions, cost and singularity policy, primal plan,
dual potentials, rigorously checked constraints, and
L\ensuremath{\le}V\ensuremath{\le}U. State whether V is a finite
discrete optimum or the continuum functional. Sinkhorn candidates must
be converted to feasible bounds for the unregularized problem, with
entropic bias controlled.

\subsection{Validation and rejection controls}

Verify a small rational-cost transport LP exactly and a one-dimensional
quadratic-cost quantile example as a separate solver test. For Coulomb
densities, use a case with a known co-motion construction under its
actual hypotheses. Reject comparing normalized densities with unequal
total mass unless normalization or an unbalanced-OT model is explicitly
chosen.

\subsection{Milestones and research extension}

Baseline: certified two-marginal discrete bounds. Strong: one continuum
enclosure with singularity/tail control. Research extension:
N\textgreater2 with symmetry reduction and validated bounds, or a
demonstrably improved relaxation. Do not promise deterministic Monge
maps for arbitrary multimarginal inputs.

Milestones are acceptance gates, not calendar promises. Before a run,
freeze the concrete instance/domain, tolerance, compute budget, and
permitted dependencies. A bounded search that fails to find a result
must preserve an unresolved outcome.

\subsection{Starting instruction}

Work on the baseline above. First make every input and convention
explicit, then validate the reference case. Build the checker alongside
the calculation. Label numerical evidence, finite-model certificates,
and continuum theorems separately; return the strongest justified result
and identify the exact unmet proof obligation.

\subsection{Primary references and their role}

\begin{itemize}
\tightlist
\item
  \href{https://arxiv.org/abs/1104.0603}{Cotar, Friesecke and
  Klüppelberg, Density functional theory and optimal transportation with
  Coulomb cost} --- Strong-correlation/semiclassical connection; theorem
  hypotheses must be checked for each density.
\end{itemize}

References support definitions or limitations; the scope choices and
proposed milestones are editorial recommendations, not claims that these
results have already been obtained.

\subsection{Displayed formulation and claim boundary}

\href{https://github.com/rick-stevens-ai/Pure-Thought/blob/main/PRDs/18-Optimal-Transport-Chemistry.md}{Canonical
specification}.

\[\int\rho=N,\qquad \mu=\rho/N,\qquad
c(r_1,\ldots,r_N)=\sum_{i<j}|r_i-r_j|^{-1},\]

\[V_{\mathrm{SCE}}[\rho]=\inf_{\Gamma\in\Pi(\mu,\ldots,\mu)}\int c\,d\Gamma,
\qquad \sum_i u(r_i)\le c(r_1,\ldots,r_N),\]

\[\underbrace{\int u\rho}_{L}\ \le\ V_{\mathrm{SCE}}[\rho]\ \le\
\underbrace{\int c\,d\Gamma_{\mathrm{feasible}}}_{U}.\]

The coupling is a probability measure; it may be symmetrized and is not
antisymmetric. This is a Coulomb interaction functional, not exact
exchange. For a discrete problem, the same inequalities apply only to
the discrete optimum until continuum error is bounded.

\begin{center}\rule{0.5\linewidth}{0.5pt}\end{center}

\section{19. Certified equilibria and conditional persistence in
mass-action networks}

\textbf{Revised specification --- 2026-09-11.}
\href{https://github.com/rick-stevens-ai/Pure-Thought/blob/main/reviews/19-Chemical-Reaction-Networks.md}{Detailed
critique} \ensuremath{\cdot}
\href{https://github.com/rick-stevens-ai/Pure-Thought/blob/main/archive/original-PRDs/19-Chemical-Reaction-Networks.md}{Original
description} \ensuremath{\cdot}
\href{https://github.com/rick-stevens-ai/Pure-Thought/blob/main/reviews/EVIDENCE-STANDARD.md}{Evidence
standard}

\subsection{Plain-language guide}

\subsubsection{The problem in everyday terms}

A list of chemical reactions can settle into different long-term
concentration patterns. Some networks have one resting state; others can
behave like a switch, with two stable states separated by an unstable
one. This problem asks whether we can find every positive equilibrium of
a small, explicitly specified network and prove which are locally
stable. It also distinguishes never reaching zero at a finite time from
genuinely avoiding extinction in the long run.

\subsubsection{Key terms}

\begin{itemize}
\item
  \textbf{Reaction network} --- A set of chemical species and rules
  converting specified combinations of reactants into products.
\item
  \textbf{Mass-action kinetics} --- Rate laws in which each reaction
  rate is proportional to powers of its reactant concentrations,
  multiplied by a rate constant.
\item
  \textbf{Stoichiometry matrix S} --- A table describing how each
  reaction changes the amount of each species.
\item
  \textbf{Conservation law} --- A combination of species amounts that
  remains constant under the network's reactions.
\item
  \textbf{Compatibility class} --- The concentration states accessible
  under the same conserved quantities and stoichiometric constraints.
\item
  \textbf{Equilibrium or steady state} --- A concentration pattern where
  the net rate of change is zero.
\item
  \textbf{Local stability} --- The tendency of states starting
  sufficiently near an equilibrium to remain nearby, with asymptotic
  stability also implying return toward it.
\item
  \textbf{Multistationarity and multistability} --- Having multiple
  equilibria, and having multiple stable equilibria, respectively.
\item
  \textbf{Persistence} --- A long-time condition preventing each
  initially present species from approaching zero along the trajectory.
\end{itemize}

\subsubsection{Why this matters}

Stable switches and coexistence are important organizing ideas in
chemistry and biological regulation. An exact small-network analysis can
establish whether a proposed mechanism even supports those behaviors. It
also makes numerical simulation more interpretable: a trajectory can
miss another equilibrium or appear stable simply because it has not been
followed long enough.

\subsubsection{What progress would mean}

Progress would give complete equilibrium information for the stated
network, rates and conserved totals. A theorem covering many rate
choices would be stronger; identifying a switch in a model would not by
itself establish its role in a living system.

\subsection{Core question}

For a finite mass-action network with specified rates and conservation
totals, can all positive equilibria in one compatibility class be
isolated and their local stability certified?

\subsection{Scope and assumptions}

Input integer reactant/product complexes, positive rational rate
constants, and positive rational initial totals. Use small networks
first, initially at most four species and eight reactions. Distinguish
fixed-parameter root analysis from theorems quantified over all rates.

Keep persistence as a separate theorem-driven module. Catalytic RAF
structure belongs primarily to Challenge 29 and does not by itself
determine mass-action stability.

\subsection{Mathematical target}

Let S be the stoichiometric matrix, f(x)=Sv(x),
v\_r=\ensuremath{\kappa}\_r x\^{}\{y\_r\}, and let rows of L span ker
S\ensuremath{{}^{T}}. Solve f(x)=0, Lx=Lx\ensuremath{{}_{0}},
x\textgreater0. Remove dependent equations before root isolation.
Deficiency is \ensuremath{\delta}=number of
complexes\ensuremath{-}number of linkage classes\ensuremath{-}rank S.

For an isolated equilibrium, test the Jacobian on im S, not on the full
space with conservation-induced zero eigenvalues. Multistationarity
means at least two positive equilibria in the same class; multistability
additionally needs stability. Persistence means
liminf\_\{t\ensuremath{\to}\ensuremath{\infty}\}x\_i(t)>0 and requires
global hypotheses and boundedness checks.

\subsection{Required outputs and proof obligations}

Return exact network data, deficiency/reversibility facts, conservation
laws, isolating boxes with existence/uniqueness proofs, and coverage of
the admissible domain. A positive-dimensional equilibrium set is a
separate valid outcome. Persistence output must name a sufficient
theorem and check every hypothesis, or return unresolved.

\subsection{Validation and rejection controls}

Use A\ensuremath{\rightleftharpoons}B as a deficiency-zero benchmark.
Use the one-species open network
\ensuremath{\varnothing}\ensuremath{\rightleftharpoons}X and
2X\ensuremath{\rightleftharpoons}3X with rates giving
f(x)=6\ensuremath{-}11x+6x\ensuremath{{}^{2}}\ensuremath{-}x\ensuremath{{}^{3}};
isolate x=1,2,3 and distinguish two stable roots from three equilibria.
Include a boundary-approaching trajectory example to distinguish
positivity from persistence.

\subsection{Milestones and research extension}

Baseline: exact structural facts and complete small fixed-parameter root
enumeration. Strong: certified local stability and one applicable
persistence theorem. Research extension: parameter-region proofs or a
new structural criterion, with quantifiers stated explicitly.

Milestones are acceptance gates, not calendar promises. Before a run,
freeze the concrete instance/domain, tolerance, compute budget, and
permitted dependencies. A bounded search that fails to find a result
must preserve an unresolved outcome.

\subsection{Starting instruction}

Work on the baseline above. First make every input and convention
explicit, then validate the reference case. Build the checker alongside
the calculation. Label numerical evidence, finite-model certificates,
and continuum theorems separately; return the strongest justified result
and identify the exact unmet proof obligation.

\subsection{Primary references and their role}

\begin{itemize}
\tightlist
\item
  \href{https://arxiv.org/abs/q-bio/0608019}{Angeli, De Leenheer and
  Sontag, A Petri Net approach to persistence} --- Checkable sufficient
  conditions for persistence, with hypotheses.
\end{itemize}

References support definitions or limitations; the scope choices and
proposed milestones are editorial recommendations, not claims that these
results have already been obtained.

\subsection{Displayed formulation and claim boundary}

\href{https://github.com/rick-stevens-ai/Pure-Thought/blob/main/PRDs/19-Chemical-Reaction-Networks.md}{Canonical
specification}.

\[\dot x=S v(x),\qquad v_j(x)=k_j\prod_i x_i^{\alpha_{ij}},\qquad
x\in(x_0+\operatorname{im}S)\cap\mathbb R_{\ge0}^n.\]

\[f(x)=6-11x+6x^2-x^3=-(x-1)(x-2)(x-3),
\qquad (f'(1),f'(2),f'(3))=(-2,1,-2).\]

The scalar benchmark has two locally stable positive equilibria and one
unstable equilibrium. General equilibrium enumeration needs complete
coverage within the compatibility class, boundary treatment, and a
reduced Jacobian. Positivity invariance, persistence and multistability
are different claims.

\begin{center}\rule{0.5\linewidth}{0.5pt}\end{center}

\section{20. Path-integral thermodynamics with a transparent error
budget}

\textbf{Revised specification --- 2026-09-11.}
\href{https://github.com/rick-stevens-ai/Pure-Thought/blob/main/reviews/20-Ab-Initio-Path-Integrals.md}{Detailed
critique} \ensuremath{\cdot}
\href{https://github.com/rick-stevens-ai/Pure-Thought/blob/main/archive/original-PRDs/20-Ab-Initio-Path-Integrals.md}{Original
description} \ensuremath{\cdot}
\href{https://github.com/rick-stevens-ai/Pure-Thought/blob/main/reviews/EVIDENCE-STANDARD.md}{Evidence
standard}

\subsection{Plain-language guide}

\subsubsection{The problem in everyday terms}

At small scales, particles cannot always be treated as tiny classical
balls sitting at definite positions. Their quantum spreading affects
thermal averages, especially for light particles. Path-integral methods
replace part of this quantum calculation with a sampling problem
involving an artificial ring of linked copies of a particle. This
challenge checks that the sampling produces the right equilibrium
quantities and that each source of error is reported separately.

\subsubsection{Key terms}

\begin{itemize}
\item
  \textbf{Thermodynamics} --- The study of quantities such as energy and
  temperature describing equilibrium behavior.
\item
  \textbf{Path integral} --- A formulation of quantum mechanics
  involving a sum over possible paths; here it is used in imaginary time
  for equilibrium properties.
\item
  \textbf{Imaginary time} --- A mathematical reformulation of time used
  to connect quantum calculations with thermal statistics, not a
  physical clock reading.
\item
  \textbf{Ring polymer and beads} --- The artificial closed chain and
  its particle copies used to approximate the imaginary-time path
  integral.
\item
  \textbf{PIMD} --- Path-integral molecular dynamics, a method using
  auxiliary dynamics to sample the ring-polymer equilibrium
  distribution.
\item
  \textbf{Estimator} --- A formula calculated from sampled
  configurations to estimate a desired physical average.
\item
  \textbf{Discretization bias} --- Systematic error caused here by using
  a finite number of beads or a finite integration step.
\item
  \textbf{Sampling uncertainty} --- Uncertainty from using only a finite
  sample of configurations.
\item
  \textbf{Born--Oppenheimer surface} --- An approximate potential energy
  for nuclear motion obtained by treating electronic behavior
  separately.
\end{itemize}

\subsubsection{Why this matters}

Quantum thermal effects matter for understanding how light particles
behave in molecules and materials. But an accurate-looking simulation
can combine several hidden errors: too few beads, incomplete sampling,
time-step effects and an approximate potential. A transparent error
budget makes the calculation useful as a scientific measurement of the
model rather than just a plausible plot.

\subsubsection{What progress would mean}

The initial result would reproduce known oscillator averages and then
handle one controlled anharmonic case. It would validate equilibrium
predictions; it would not automatically give exact real-time reaction
rates or remove approximations in an electronic-structure model.

\subsection{Core question}

For a specified one-dimensional quantum Hamiltonian, can path-integral
sampling reproduce equilibrium energy and position moments with
separately assessed discretization and statistical errors?

\subsection{Scope and assumptions}

Start with H=p\ensuremath{{}^{2}}/(2m)+V(q), m\textgreater0,
\ensuremath{\beta}>0, and a confining harmonic potential; extend to an
explicitly defined anharmonic potential. Nuclei are distinguishable
Boltzmann particles. Report only equilibrium observables initially.

An ab initio extension must specify the Born--Oppenheimer surface,
electronic method, basis, and force errors. It is not exact quantum
chemistry merely because no empirical force field is used. Real-time
rates and exchange statistics require distinct methods.

\subsection{Mathematical target}

Set \ensuremath{\beta}\_P=\ensuremath{\beta}/P and
\ensuremath{\omega}\_P=P/(\ensuremath{\beta}\ensuremath{\hbar}). Sample
the coordinate density proportional to

\[\exp\left[-\beta _P\sum_{s=1}^P\left\{\tfrac12m\omega _P^{2}(q_s-q_{s+1})^{2}+V(q_s)\right\}\right],\quad q_{P+1}=q_1.\]

Coordinate observables use P\ensuremath{{}^{-1}}\ensuremath{\Sigma}\_s
A(q\_s), generally not A(q\_centroid). For one dimension the primitive
energy estimator is
P/(2\ensuremath{\beta})\ensuremath{-}mP/(2\ensuremath{\beta}\ensuremath{{}^{2}}\ensuremath{\hbar}\ensuremath{{}^{2}})\ensuremath{\Sigma}\_s(q\_s\ensuremath{-}q\_\{s+1\})\ensuremath{{}^{2}}+P\ensuremath{{}^{-1}}\ensuremath{\Sigma}\_sV(q\_s).
Derive the chosen estimator from the partition function to fix all
temperature factors.

\subsection{Required outputs and proof obligations}

Return algorithm, seeds, temperatures, bead counts, time-step settings,
estimators, and an error ledger separating sampling, integration,
finite-P, and potential/model errors. Use ``statistical interval'' for
sampling uncertainty. Reserve ``rigorous enclosure'' for cases with
proved sampling/discretization/tail bounds.

\subsection{Validation and rejection controls}

Recover
E=(\ensuremath{\hbar}\ensuremath{\omega}/2)coth(\ensuremath{\beta}\ensuremath{\hbar}\ensuremath{\omega}/2)
and
\ensuremath{\langle}q\ensuremath{{}^{2}}\ensuremath{\rangle}=\ensuremath{\hbar}coth(\ensuremath{\beta}\ensuremath{\hbar}\ensuremath{\omega}/2)/(2m\ensuremath{\omega})
for the oscillator over a declared \ensuremath{\beta} range. Compare
multiple P and time steps, independent chains, and autocorrelation
diagnostics. Validate anharmonic results against a separately converged
or certified spectral method. A convergence fit alone cannot certify its
asymptote.

\subsection{Milestones and research extension}

Baseline: oscillator thermodynamics within predeclared statistical and
bias tolerances. Strong: one anharmonic error-controlled result.
Research extension: a small ab initio system with a full error budget;
rigorous certification requires more than SCF convergence and stable
trajectories.

Milestones are acceptance gates, not calendar promises. Before a run,
freeze the concrete instance/domain, tolerance, compute budget, and
permitted dependencies. A bounded search that fails to find a result
must preserve an unresolved outcome.

\subsection{Starting instruction}

Work on the baseline above. First make every input and convention
explicit, then validate the reference case. Build the checker alongside
the calculation. Label numerical evidence, finite-model certificates,
and continuum theorems separately; return the strongest justified result
and identify the exact unmet proof obligation.

\subsection{Primary references and their role}

\begin{itemize}
\tightlist
\item
  \href{https://arxiv.org/abs/1009.1045}{Ceriotti et al., Efficient
  stochastic thermostatting of path integral molecular dynamics} ---
  Normal-mode sampling and stochastic thermostat methods.
\end{itemize}

References support definitions or limitations; the scope choices and
proposed milestones are editorial recommendations, not claims that these
results have already been obtained.

\subsection{Displayed formulation and claim boundary}

\href{https://github.com/rick-stevens-ai/Pure-Thought/blob/main/PRDs/20-Ab-Initio-Path-Integrals.md}{Canonical
specification}.

\[\beta_P=\beta/P,\quad\omega_P=P/(\beta\hbar),\quad q_{P+1}=q_1,\]

\[\pi_P(q)\propto\exp\!\left[-\beta_P\sum_{s=1}^P
\left\{\frac{m\omega_P^2}{2}(q_s-q_{s+1})^2+V(q_s)\right\}\right],\]

\[E_{\mathrm{prim}}=\frac{P}{2\beta}
-\frac{mP}{2\beta^2\hbar^2}\sum_s(q_s-q_{s+1})^2
+\frac1P\sum_sV(q_s).\]

Coordinate observables use bead averages. Derive the estimator from the
same partition-function convention. Separate sampling, integration,
finite-\(P\) and model errors; convergence fits and empirical
autocorrelation diagnostics do not by themselves give rigorous
enclosures.

\begin{center}\rule{0.5\linewidth}{0.5pt}\end{center}

\section{21. Quantum LDPC instances with independently certified code
parameters}

\textbf{Revised specification --- 2026-09-11.}
\href{https://github.com/rick-stevens-ai/Pure-Thought/blob/main/reviews/21-Quantum-LDPC-Codes.md}{Detailed
critique} \ensuremath{\cdot}
\href{https://github.com/rick-stevens-ai/Pure-Thought/blob/main/archive/original-PRDs/21-Quantum-LDPC-Codes.md}{Original
description} \ensuremath{\cdot}
\href{https://github.com/rick-stevens-ai/Pure-Thought/blob/main/reviews/EVIDENCE-STANDARD.md}{Evidence
standard}

\subsection{Plain-language guide}

\subsubsection{The problem in everyday terms}

Quantum information is fragile, but simply copying an unknown quantum
state is not allowed. An error-correcting code instead stores
information jointly across many physical qubits so that selected
measurements reveal errors without directly reading the encoded
information. This problem builds sparse-check codes and verifies how
much information they store, how many errors their structure can
withstand, and how a chosen recovery algorithm performs.

\subsubsection{Key terms}

\begin{itemize}
\item
  \textbf{Qubit} --- A quantum information unit whose state can involve
  superpositions of two basis states.
\item
  \textbf{Physical and logical qubits} --- The individual qubits used by
  the device, and the protected information units encoded across them.
\item
  \textbf{Quantum error-correcting code} --- A structured encoding that
  permits detection and correction of specified errors.
\item
  \textbf{LDPC} --- Low-density parity-check: a code family with checks
  involving bounded numbers of qubits and each qubit participating in
  bounded numbers of checks.
\item
  \textbf{CSS code} --- A code built from compatible binary checks for
  two types of quantum error, named after Calderbank, Shor and Steane.
\item
  \textbf{Stabilizer} --- An operator that leaves every valid encoded
  state unchanged; measuring checks can reveal violations.
\item
  \textbf{Syndrome} --- The collection of check outcomes indicating what
  kind of error may have occurred.
\item
  \textbf{Code distance d} --- The smallest number of qubits on which an
  undetectable nontrivial logical operation can act.
\item
  \textbf{Decoder} --- An algorithm that uses a syndrome to choose a
  recovery operation.
\item
  \textbf{Hypergraph product} --- A construction combining two classical
  parity-check matrices into a quantum code.
\end{itemize}

\subsubsection{Why this matters}

Useful quantum computation requires reliable logical information despite
noisy components. Sparse checks offer a route to distributing error
correction efficiently, but code parameters and decoder behavior must be
demonstrated rather than assumed. Exact small examples are valuable for
detecting mistakes that large simulations can conceal.

\subsubsection{What progress would mean}

Success would provide verified code parameters and decoder results under
a stated error model. It would help compare constructions fairly, while
leaving hardware connectivity, noisy measurements and the full cost of
fault-tolerant computation as separate requirements.

\subsection{Core question}

For a specified CSS construction, can code dimension and distance bounds
be verified independently, and can decoder performance be measured under
one explicit noise model?

\subsection{Scope and assumptions}

Start with hypergraph-product codes from supplied binary parity-check
matrices H\ensuremath{{}_{1}} and H\ensuremath{{}_{2}}. Include the
product of two full-rank 3\ensuremath{\times}7 Hamming parity-check
matrices as a small benchmark. Separate finite instances, asymptotic
families, and noisy syndrome extraction.

Use independent X errors with perfect syndrome measurements for the
first decoder experiment. A family is LDPC only when both row and column
weights stay bounded as n grows.

\subsection{Mathematical target}

For H\_i of size m\_i\ensuremath{\times}n\_i define

\[H_X=[H_1\otimes I_{n_2}\mid I_{m_1}\otimes H_2^{\mathsf T}],\quad H_Z=[I_{n_1}\otimes H_2\mid H_1^{\mathsf T}\otimes I_{m_2}].\]

Check \(H_XH_Z^{\mathsf T}=0\) over F\ensuremath{{}_{2}},
n=n\ensuremath{{}_{1}}n\ensuremath{{}_{2}}+m\ensuremath{{}_{1}}m\ensuremath{{}_{2}},
\(k=n-\operatorname{rank}H_X-\operatorname{rank}H_Z\). Define
\(d_X=\min\{|x|:x\in\ker H_Z\setminus\operatorname{row}H_X\}\) and d\_Z
analogously, \(d=\min(d_X,d_Z)\). A found logical operator upper-bounds
distance; ruling out all smaller logical operators lower-bounds it.

\subsection{Required outputs and proof obligations}

Export sparse matrices, GF(2) row reductions, logical bases, explicit
upper witnesses and lower-bound proofs. SAT certificates need a verified
encoding of nontrivial homology and weight. Decoder experiments report
physical noise, recovery rule, logical-failure test, sample counts,
confidence intervals, and runtime.

\subsection{Validation and rejection controls}

For the two 3\ensuremath{\times}7 inputs verify n=58 and k=16; certify
the claimed distance independently rather than assuming it. Inject all
small-weight errors on tiny codes. A matching syndrome is not success
unless error plus recovery is a stabilizer. Compare optimized checks
with an independent GF(2) implementation.

\subsection{Milestones and research extension}

Baseline: exact small product code and certified distance bracket.
Strong: several instances plus honest finite-size decoder curves.
Research extension: a new family tradeoff or proved decoder guarantee.
Reproducing asymptotically good codes is not itself a new existence
result.

Milestones are acceptance gates, not calendar promises. Before a run,
freeze the concrete instance/domain, tolerance, compute budget, and
permitted dependencies. A bounded search that fails to find a result
must preserve an unresolved outcome.

\subsection{Starting instruction}

Work on the baseline above. First make every input and convention
explicit, then validate the reference case. Build the checker alongside
the calculation. Label numerical evidence, finite-model certificates,
and continuum theorems separately; return the strongest justified result
and identify the exact unmet proof obligation.

\subsection{Primary references and their role}

\begin{itemize}
\tightlist
\item
  \href{https://arxiv.org/abs/2111.03654}{Panteleev and Kalachev,
  Asymptotically Good Quantum and Locally Testable Classical LDPC Codes}
  --- Existing asymptotically good quantum LDPC constructions; novelty
  baseline.
\end{itemize}

References support definitions or limitations; the scope choices and
proposed milestones are editorial recommendations, not claims that these
results have already been obtained.

\subsection{Displayed formulation and claim boundary}

\href{https://github.com/rick-stevens-ai/Pure-Thought/blob/main/PRDs/21-Quantum-LDPC-Codes.md}{Canonical
specification}.

\[H_X=[H_1\otimes I_{n_2}\mid I_{m_1}\otimes H_2^{\mathsf T}],
\qquad H_Z=[I_{n_1}\otimes H_2\mid H_1^{\mathsf T}\otimes I_{m_2}],\]

\[H_XH_Z^{\mathsf T}=0,\quad n=n_1n_2+m_1m_2,\quad
k=n-\operatorname{rank}_{\mathbb F_2}H_X-\operatorname{rank}_{\mathbb F_2}H_Z,\]

\[d_X=\min_{x\in\ker H_Z\setminus\operatorname{row}H_X}|x|,
\qquad d_Z=\min_{z\in\ker H_X\setminus\operatorname{row}H_Z}|z|,
\qquad d=\min(d_X,d_Z).\]

Two full-rank \(3\times7\) Hamming checks give \(n=58\), \(k=16\); the
included independent validator also checks \(d=3\). A logical witness
gives an upper distance bound; exclusion of all smaller logicals gives
the lower bound. Finite code parameters do not establish an asymptotic
LDPC family or decoding threshold.

\begin{center}\rule{0.5\linewidth}{0.5pt}\end{center}

\section{22. Bell inequalities with certified local and quantum value
brackets}

\textbf{Revised specification --- 2026-09-11.}
\href{https://github.com/rick-stevens-ai/Pure-Thought/blob/main/reviews/22-Bell-Inequalities-Nonlocality.md}{Detailed
critique} \ensuremath{\cdot}
\href{https://github.com/rick-stevens-ai/Pure-Thought/blob/main/archive/original-PRDs/22-Bell-Inequalities-Nonlocality.md}{Original
description} \ensuremath{\cdot}
\href{https://github.com/rick-stevens-ai/Pure-Thought/blob/main/reviews/EVIDENCE-STANDARD.md}{Evidence
standard}

\subsection{Plain-language guide}

\subsubsection{The problem in everyday terms}

Two distant experimenters choose measurements and compare their
outcomes. If the outcomes were produced only by shared classical
instructions, their correlations would obey particular inequalities.
Quantum systems can exceed some of those limits. This problem calculates
both the classical ceiling and a reliable quantum range for a specified
test, so that the size of the difference is established mathematically.

\subsubsection{Key terms}

\begin{itemize}
\item
  \textbf{Bell scenario} --- A specification of the parties, their
  measurement choices and the possible outcomes.
\item
  \textbf{Local hidden-variable model} --- A model in which outcomes
  depend on local choices and shared classical information, with
  measurement-choice independence among the assumptions.
\item
  \textbf{Bell inequality} --- A linear restriction satisfied by the
  correlations of such local models.
\item
  \textbf{CHSH} --- The simplest standard two-party test with two binary
  measurement choices per party.
\item
  \textbf{Quantum strategy} --- A specified shared quantum state and
  local measurements producing correlations.
\item
  \textbf{POVM} --- A set of positive operators summing to the identity
  that describes a general quantum measurement.
\item
  \textbf{Tsirelson bound} --- A quantum upper limit on a Bell
  expression, with 2\ensuremath{\surd}2 the familiar CHSH value in its
  standard normalization.
\item
  \textbf{NPA hierarchy} --- A sequence of matrix-positivity tests
  providing outer bounds on quantum correlations.
\item
  \textbf{Commuting-operator model} --- A framework in which different
  parties' measurements commute; its general correlation set differs
  from the finite-dimensional tensor-product setting.
\item
  \textbf{Facet} --- A full boundary face of the appropriate dimension
  of the classical correlation polytope.
\end{itemize}

\subsubsection{Why this matters}

Bell inequalities make the difference between classical local
explanations and quantum correlations precise. They also underlie ways
to certify quantum behavior from observed input-output relationships.
Reliable quantum bounds are essential: a numerical optimization value
without a valid upper certificate or achieving strategy may misstate
what quantum theory permits.

\subsubsection{What progress would mean}

A complete small test would include an exact classical maximum, an
explicit quantum strategy and a matching or enclosing upper bound.
Applications to randomness or cryptography would need additional
statistical and security assumptions beyond this mathematical benchmark.

\subsection{Core question}

For a fixed rational Bell functional in a small bipartite scenario, can
its exact local value and a certified quantum lower/upper bracket be
produced?

\subsection{Scope and assumptions}

Start with two parties, two binary settings each, and CHSH. Fix
normalization and no-signaling affine coordinates. The target quantum
value is the supremum over finite-dimensional tensor-product strategies;
identify its closure when discussing limits. Finite NPA levels provide
outer bounds through the commuting-operator framework.

Facet discovery and randomness certification are extensions with
separate definitions. Do not equate a finite SDP solution with a
realizable strategy.

\subsection{Mathematical target}

For \ensuremath{\beta}(P)=\ensuremath{\Sigma}\_abxy
c\_abxyP(ab\textbar xy), compute \ensuremath{\beta}\_L by evaluating
every deterministic local strategy. A verified finite-dimensional state
and POVMs give L\_Q\ensuremath{\le}\ensuremath{\beta}\_Q. A certified
NPA dual or noncommutative SOS gives
\ensuremath{\beta}\_Q\ensuremath{\le}U\_Q. Report \ensuremath{\beta}\_L
and {[}L\_Q,U\_Q{]}.

The NPA limit characterizes commuting-operator correlations, which need
not equal the closure of finite-dimensional tensor-product correlations.
Finite-level tightness needs an explicit achieving strategy or a valid
extraction theorem.

\subsection{Required outputs and proof obligations}

Export coefficients, local vertices, state/measurements, moment-word
list, algebraic relations, dual matrices/SOS, and exact or interval
objective bounds. To certify a facet in affine dimension D, verify the
inequality for every vertex and show saturating vertices have affine
dimension D\ensuremath{-}1. Declare relabeling symmetries when counting
classes.

\subsection{Validation and rejection controls}

Check all 16 CHSH deterministic strategies and \ensuremath{\beta}\_L=2.
Reproduce 2\ensuremath{\surd}2 with an explicit two-qubit strategy and
an operator upper-bound identity. Reject an invalid dual PSD matrix or a
wrong commutation relation. Distinguish eight CHSH variants from one
class under relabeling; positivity facets must be handled separately.

\subsection{Milestones and research extension}

Baseline: exact CHSH lower and upper certificates. Strong: a preselected
additional rational functional with a certified bracket, whether or not
it closes. Research extension: a new inequality with normalized
comparison, or device-independent entropy bounds with an explicit
adversary and statistical model.

Milestones are acceptance gates, not calendar promises. Before a run,
freeze the concrete instance/domain, tolerance, compute budget, and
permitted dependencies. A bounded search that fails to find a result
must preserve an unresolved outcome.

\subsection{Starting instruction}

Work on the baseline above. First make every input and convention
explicit, then validate the reference case. Build the checker alongside
the calculation. Label numerical evidence, finite-model certificates,
and continuum theorems separately; return the strongest justified result
and identify the exact unmet proof obligation.

\subsection{Primary references and their role}

\begin{itemize}
\item
  \href{https://arxiv.org/abs/0803.4290}{Navascués, Pironio and Acín, A
  convergent hierarchy of semidefinite programs} --- NPA convergence to
  a commuting-measurement representation.
\item
  \href{https://arxiv.org/abs/2001.04383}{Ji et al., MIP*=RE} --- Strict
  separation of the relevant correlation sets.
\end{itemize}

References support definitions or limitations; the scope choices and
proposed milestones are editorial recommendations, not claims that these
results have already been obtained.

\subsection{Displayed formulation and claim boundary}

\href{https://github.com/rick-stevens-ai/Pure-Thought/blob/main/PRDs/22-Bell-Inequalities-Nonlocality.md}{Canonical
specification}.

\[\beta(P)=\sum_{abxy}c_{abxy}P(ab|xy),\quad
\beta_L=\max_{P\text{ deterministic local}}\beta(P),\quad
L_Q\le\beta_Q\le U_Q.\]

A finite-dimensional state and measurements certify \(L_Q\); a verified
NPA dual or operator SOS certifies \(U_Q\). For CHSH, \(\beta_L=2\) and
an achieving two-qubit strategy plus an upper-bound identity give
\(\beta_Q=2\sqrt2\). NPA's commuting-operator limit must not be
identified with the closure of finite-dimensional tensor-product
strategies. A facet in affine dimension \(D\) needs saturating affine
rank \(D-1\).

\begin{center}\rule{0.5\linewidth}{0.5pt}\end{center}

\section{23. Entanglement certification with explicit inconclusive
outcomes}

\textbf{Revised specification --- 2026-09-11.}
\href{https://github.com/rick-stevens-ai/Pure-Thought/blob/main/reviews/23-Entanglement-Measures-Witnesses.md}{Detailed
critique} \ensuremath{\cdot}
\href{https://github.com/rick-stevens-ai/Pure-Thought/blob/main/archive/original-PRDs/23-Entanglement-Measures-Witnesses.md}{Original
description} \ensuremath{\cdot}
\href{https://github.com/rick-stevens-ai/Pure-Thought/blob/main/reviews/EVIDENCE-STANDARD.md}{Evidence
standard}

\subsection{Plain-language guide}

\subsubsection{The problem in everyday terms}

Two quantum systems can be related in ways that no mixture of
independent local states can reproduce. That is entanglement. Given a
mathematical description of a joint state, this problem asks for
evidence that it is entangled, evidence that it is separable, or an
honest statement that the available tests cannot decide. It also asks
how to quantify one particular aspect of entanglement without mixing
different measures together.

\subsubsection{Key terms}

\begin{itemize}
\item
  \textbf{Bipartite system} --- A quantum system divided into two
  identified parts, such as two qubits held by different parties.
\item
  \textbf{Density matrix \ensuremath{\rho}} --- An operator describing a
  quantum state, including statistical mixtures.
\item
  \textbf{Separable state} --- A mixture of product states of the two
  parts; correlations in such a decomposition can arise from shared
  classical information.
\item
  \textbf{Entangled state} --- A state that has no separable
  decomposition across the chosen division.
\item
  \textbf{Entanglement witness} --- An observable nonnegative on every
  separable state but negative on the state being tested.
\item
  \textbf{Partial transpose} --- A matrix operation transposing the
  indices of just one subsystem, useful for testing entanglement.
\item
  \textbf{PPT} --- Positive partial transpose: a necessary separability
  condition that is sufficient only in certain small dimensions.
\item
  \textbf{Negativity} --- A quantity measuring the negative part of the
  partial-transpose spectrum; zero does not generally imply
  separability.
\item
  \textbf{DPS hierarchy} --- Tests based on whether the state can be
  extended to larger symmetric states; failure can certify entanglement.
\item
  \textbf{Concurrence} --- A particular entanglement measure with a
  useful explicit formula for two-qubit states.
\end{itemize}

\subsubsection{Why this matters}

Entanglement is an important resource in many quantum-information
protocols, but claiming it from an inconclusive test can invalidate a
proposed demonstration. Conversely, explicit witnesses show exactly what
establishes its presence. Quantitative bounds are useful because
different states may possess very different strengths or kinds of
quantum correlations.

\subsubsection{What progress would mean}

The deliverable would make its verdict inspectable: a witness, a
decomposition, or unresolved status. That would provide a reliable
component for analyzing quantum states, rather than promise a
universally easy entanglement test.

\subsection{Core question}

For a finite bipartite density matrix, can one produce a checked
entanglement witness or a separable decomposition, and bound a
specifically chosen entanglement measure?

\subsection{Scope and assumptions}

Input subsystem dimensions and an exact rational/algebraic density
matrix, or a rigorously described uncertainty set. Begin with
2\ensuremath{\times}2 systems; then include 3\ensuremath{\times}3
PPT-entangled examples. Validate Hermiticity, positivity and unit trace
before classification.

The primary decision outputs are entangled, separable, or unresolved.
Restrict concurrence and its entanglement-of-formation formula to two
qubits. Treat general distillable entanglement and multipartite
classification as separate research problems.

\subsection{Mathematical target}

Compute partial transpose in a fixed tensor-index convention. Negativity
is \(\mathcal N(\rho)=(\|\rho^{T_B}\|_1-1)/2\). For entanglement give W
satisfying \(\langle a\otimes b|W|a\otimes b\rangle\ge0\) for all
product vectors and Tr(W\ensuremath{\rho})<0. For separability give
\(\rho=\sum_jp_j|a_jb_j\rangle\langle a_jb_j|\) with
p\_j\ensuremath{\ge}0 and \ensuremath{\Sigma}p\_j=1, or another
sufficient theorem in its valid dimensions.

PPT/DPS feasibility is generally only a necessary separability test. DPS
infeasibility with a validated dual can certify entanglement; passing a
finite level generally remains unresolved.

\subsection{Required outputs and proof obligations}

Export input validation, partial-transpose spectrum enclosures, measure
intervals, and a witness/decomposition with an independent checker.
Normalize witnesses before optimizing them. For an uncertain input,
require the detection margin to hold throughout the uncertainty set.

\subsection{Validation and rejection controls}

Use Bell states, product states, and
\ensuremath{\rho}\_p=p|\ensuremath{\psi}\ensuremath{{}^{-}}\ensuremath{\rangle}\ensuremath{\langle}\ensuremath{\psi}\ensuremath{{}^{-}}|+(1\ensuremath{-}p)I/4,
0\ensuremath{\le}p\ensuremath{\le}1. In this declared convention
separability holds for p\ensuremath{\le}1/3. Verify negativity 1/2 and
concurrence 1 for a Bell state. Include a 3\ensuremath{\times}3
PPT-entangled control so zero negativity is not mislabeled separability.

\subsection{Milestones and research extension}

Baseline: checked two-qubit decisions and measure values. Strong: a
validated DPS witness beyond PPT and a separable-decomposition example.
Research extension: robust uncertainty-set certification or normalized
optimal witnesses on a stated family.

Milestones are acceptance gates, not calendar promises. Before a run,
freeze the concrete instance/domain, tolerance, compute budget, and
permitted dependencies. A bounded search that fails to find a result
must preserve an unresolved outcome.

\subsection{Starting instruction}

Work on the baseline above. First make every input and convention
explicit, then validate the reference case. Build the checker alongside
the calculation. Label numerical evidence, finite-model certificates,
and continuum theorems separately; return the strongest justified result
and identify the exact unmet proof obligation.

\subsection{Primary references and their role}

\begin{itemize}
\tightlist
\item
  \href{https://arxiv.org/abs/quant-ph/0308032}{Doherty, Parrilo and
  Spedalieri, A complete family of separability criteria} ---
  Symmetric-extension hierarchy and dual entanglement witnesses.
\end{itemize}

References support definitions or limitations; the scope choices and
proposed milestones are editorial recommendations, not claims that these
results have already been obtained.

\subsection{Displayed formulation and claim boundary}

\href{https://github.com/rick-stevens-ai/Pure-Thought/blob/main/PRDs/23-Entanglement-Measures-Witnesses.md}{Canonical
specification}.

\[\mathcal N(\rho)=\frac{\|\rho^{T_B}\|_1-1}{2},\qquad
\langle a\otimes b|W|a\otimes b\rangle\ge0,\quad \operatorname{Tr}(W\rho)<0.\]

\[\rho=\sum_jp_j|a_jb_j\rangle\langle a_jb_j|,
\qquad p_j\ge0,\quad\sum_jp_j=1.\]

The witness conditions certify entanglement; the normalized
product-state decomposition certifies separability. In general, passing
finite PPT/DPS tests remains unresolved. For the two-qubit family
\(\rho_p=p|\psi^-\rangle\langle\psi^-|+(1-p)I/4\), \(0\le p\le1\),
separability holds exactly for \(p\le1/3\).

\begin{center}\rule{0.5\linewidth}{0.5pt}\end{center}

\section{24. Topological code verification and noise-specific decoder
evaluation}

\textbf{Revised specification --- 2026-09-11.}
\href{https://github.com/rick-stevens-ai/Pure-Thought/blob/main/reviews/24-Topological-Quantum-Error-Correction.md}{Detailed
critique} \ensuremath{\cdot}
\href{https://github.com/rick-stevens-ai/Pure-Thought/blob/main/archive/original-PRDs/24-Topological-Quantum-Error-Correction.md}{Original
description} \ensuremath{\cdot}
\href{https://github.com/rick-stevens-ai/Pure-Thought/blob/main/reviews/EVIDENCE-STANDARD.md}{Evidence
standard}

\subsection{Plain-language guide}

\subsubsection{The problem in everyday terms}

Imagine storing information in loops spread across a lattice so that a
small local disturbance usually cannot change the stored value. A
topological quantum code uses a quantum version of this idea. Local
measurements reveal the endpoints of error patterns, and a decoder tries
to join those endpoints correctly. The subtlety is that two repairs can
have the same endpoints yet differ by a loop that changes the logical
information.

\subsubsection{Key terms}

\begin{itemize}
\item
  \textbf{Toric code} --- A quantum error-correcting code on a lattice
  with periodic boundaries, topologically equivalent to a torus.
\item
  \textbf{Surface code} --- A related geometric code with carefully
  chosen open boundaries.
\item
  \textbf{Logical operator} --- An operation changing the encoded
  information while remaining consistent with the code's checks.
\item
  \textbf{Noncontractible loop} --- A loop winding around the torus that
  cannot be shrunk to a point within it.
\item
  \textbf{Homology} --- A mathematical way of distinguishing cycles
  modulo boundaries; here it identifies whether a closed error pattern
  acts logically.
\item
  \textbf{Syndrome} --- Check outcomes locating the boundary of an error
  chain.
\item
  \textbf{Recovery} --- An operation chosen to cancel the error;
  matching the syndrome is necessary but does not ensure logical
  success.
\item
  \textbf{Minimum-weight matching} --- A graph optimization method used
  in certain decoders to pair defects efficiently.
\item
  \textbf{Logical failure probability} --- The probability that error
  and recovery together change encoded information.
\item
  \textbf{Threshold} --- An asymptotic noise boundary for successful
  suppression of logical errors in a specified code-and-decoder family.
\end{itemize}

\subsubsection{Why this matters}

Geometric codes connect an intuitive picture of local errors to precise
algebraic guarantees. They are important models for protecting quantum
information using structured measurements. Understanding logical
failure, rather than only syndrome cancellation, is essential for
assessing whether increasing the code size actually improves
reliability.

\subsubsection{What progress would mean}

A useful result would prove the code's parameters and measure a fixed
decoder under a declared noise model. An estimated threshold would
remain specific to that model and decoder, not a universal device error
tolerance.

\subsection{Core question}

Can a toric-code family's algebraic parameters be proved, and can a
fixed decoder's logical failure probability be evaluated under a
precisely specified error model?

\subsection{Scope and assumptions}

Use an L\ensuremath{\times}L square torus with qubits on edges and
L\ensuremath{\ge}3. Start with independent Z errors of probability p and
perfect syndrome extraction. Specify the decoder, tie-breaking rule, and
whether recovery weights are uniform or likelihood based.

Keep exact code properties separate from statistical decoder
performance. Planar boundaries, color codes, measurement noise, and
circuit-level noise are distinct extensions.

\subsection{Mathematical target}

Over F\ensuremath{{}_{2}} use
C\ensuremath{{}_{2}}\ensuremath{\to}\ensuremath{\partial}\ensuremath{{}_{2}}
C\ensuremath{{}_{1}}\ensuremath{\to}\ensuremath{\partial}\ensuremath{{}_{1}}
C\ensuremath{{}_{0}} with
\ensuremath{\partial}\ensuremath{{}_{1}}\ensuremath{\partial}\ensuremath{{}_{2}}=0,
H\_X=\ensuremath{\partial}\ensuremath{{}_{1}},
H\_Z=\ensuremath{\partial}\ensuremath{{}_{2}}\ensuremath{{}^{T}}. Then
n=2L\ensuremath{{}^{2}}, k=2, and prove d=L using noncontractible primal
and dual cycles.

For Z error e and recovery r, success means e+r\ensuremath{\in}im
\ensuremath{\partial}\ensuremath{{}_{2}}, not merely
\ensuremath{\partial}\ensuremath{{}_{1}}(e+r)=0. A threshold assertion
concerns the limit of P\_fail(L,p) as
L\ensuremath{\to}\ensuremath{\infty} for this fixed noise/decoder
family; a finite crossing is an estimate.

\subsection{Required outputs and proof obligations}

Export lattice incidence data, ranks and logical bases, distance
argument, decoder specification, and failure counts with confidence
intervals. Exhaustive small-size evaluation can give an exact failure
polynomial in p.~Monte Carlo output must retain trials and failure
counts, including zero-failure upper intervals.

\subsection{Validation and rejection controls}

Check n,k,d on several small tori. Verify logical anticommutation via
primal/dual intersection. Enumerate all low-weight errors on tiny
instances and test the stabilizer-membership success criterion. Validate
a matching solution's optimality separately from logical decoding
success.

\subsection{Milestones and research extension}

Baseline: exact construction and distance proof with a tested decoder.
Strong: reproducible finite-size failure curves and uncertainty-aware
threshold estimates. Research extension: a rigorous threshold bound or
noisy-measurement/circuit model, with no transfer of numerical
thresholds across noise conventions.

Milestones are acceptance gates, not calendar promises. Before a run,
freeze the concrete instance/domain, tolerance, compute budget, and
permitted dependencies. A bounded search that fails to find a result
must preserve an unresolved outcome.

\subsection{Starting instruction}

Work on the baseline above. First make every input and convention
explicit, then validate the reference case. Build the checker alongside
the calculation. Label numerical evidence, finite-model certificates,
and continuum theorems separately; return the strongest justified result
and identify the exact unmet proof obligation.

\subsection{Primary references and their role}

\begin{itemize}
\tightlist
\item
  \href{https://arxiv.org/abs/quant-ph/0110143}{Dennis, Kitaev, Landahl
  and Preskill, Topological quantum memory} --- Homological coding and
  model-dependent error-correction thresholds.
\end{itemize}

References support definitions or limitations; the scope choices and
proposed milestones are editorial recommendations, not claims that these
results have already been obtained.

\subsection{Displayed formulation and claim boundary}

\href{https://github.com/rick-stevens-ai/Pure-Thought/blob/main/PRDs/24-Topological-Quantum-Error-Correction.md}{Canonical
specification}.

\[C_2\xrightarrow{\partial_2}C_1\xrightarrow{\partial_1}C_0,
\quad\partial_1\partial_2=0,\quad H_X=\partial_1,\quad H_Z=\partial_2^{\mathsf T},\]

\[n=2L^2,\quad k=2,\quad d=L\quad(L\ge3),\qquad
e+r\in\operatorname{im}\partial_2\quad\text{for successful }Z\text{-error recovery}.\]

A vanishing residual syndrome only gives \(e+r\in\ker\partial_1\) and
can hide a nontrivial logical cycle. Statistical decoder curves require
the fixed noise model and decoder; finite-size crossings are threshold
estimates, not limit theorems.

\begin{center}\rule{0.5\linewidth}{0.5pt}\end{center}

\section{25. Quantum query complexity with explicit oracle and output
models}

\textbf{Revised specification --- 2026-09-11.}
\href{https://github.com/rick-stevens-ai/Pure-Thought/blob/main/reviews/25-Quantum-Algorithms-Complexity.md}{Detailed
critique} \ensuremath{\cdot}
\href{https://github.com/rick-stevens-ai/Pure-Thought/blob/main/archive/original-PRDs/25-Quantum-Algorithms-Complexity.md}{Original
description} \ensuremath{\cdot}
\href{https://github.com/rick-stevens-ai/Pure-Thought/blob/main/reviews/EVIDENCE-STANDARD.md}{Evidence
standard}

\subsection{Plain-language guide}

\subsubsection{The problem in everyday terms}

A claim that a quantum computer is faster is incomplete until we specify
what information it receives, what answer it must return, and which
operations we count. Here the problem is deliberately simple: find one
marked item in a list when the only way to learn about an item is to
query it. We compare a quantum search procedure with a proof of how many
queries any procedure of the chosen kind must use.

\subsubsection{Key terms}

\begin{itemize}
\item
  \textbf{Oracle} --- A precisely defined black-box operation providing
  access to input information; it is a mathematical access model, not a
  prediction device.
\item
  \textbf{Query} --- One use of the oracle, counted separately from
  other computation.
\item
  \textbf{Promise problem} --- A problem whose inputs are guaranteed to
  satisfy a stated restriction, here exactly one marked item.
\item
  \textbf{Grover search} --- A quantum algorithm that amplifies the
  probability of observing a marked item using repeated oracle and
  interference operations.
\item
  \textbf{Amplitude} --- A quantum coefficient whose squared magnitude
  contributes to an outcome probability.
\item
  \textbf{Query complexity} --- The number of oracle uses required to
  solve a problem with the specified success probability.
\item
  \textbf{Upper and lower complexity bounds} --- An algorithm showing
  that a cost suffices, and a proof that no allowed algorithm can beat a
  stated cost.
\item
  \textbf{BPP and BQP} --- Classes of problems efficiently solvable with
  bounded error by classical randomized and quantum computation,
  respectively.
\item
  \textbf{Oracle separation} --- A proof that two classes differ when
  both have access to a particular oracle; it does not by itself
  separate the ordinary classes.
\item
  \textbf{Gate complexity} --- The number of elementary circuit
  operations, a different resource from oracle queries.
\end{itemize}

\subsubsection{Why this matters}

This is a clean setting in which a quantum advantage can be proved
rather than inferred from a small simulation. It also teaches how to
assess broader speedup claims fairly. Input preparation, precision and
output requirements can dominate a real application even when a query
bound looks attractive.

\subsubsection{What progress would mean}

Success would combine a working small implementation with an all-size
argument matching upper and lower query bounds. That establishes an
advantage in this access model, not a blanket claim that quantum
computers speed up every search task.

\subsection{Core question}

For unstructured search with a specified promise, can an implementation
and an analytic lower bound establish a matched quantum query
complexity?

\subsection{Scope and assumptions}

Use N=2\^{}n items with exactly one marked item and phase-oracle access
O\_f|x\ensuremath{\rangle}=(\ensuremath{-}1)\textasciicircum{}f(x)|x\ensuremath{\rangle}.
State success probability at least 2/3 and count oracle uses separately
from elementary gates, state preparation, and output verification. Begin
with small exact simulations and a proof valid for all N.

Retain quantum walks, linear systems, QAOA and oracle separations as
independent extensions, each with its own input/output and cost model.
Classical simulation wall time is not quantum runtime.

\subsection{Mathematical target}

Let \ensuremath{\theta}=arcsin(N\ensuremath{{}^{-1/2}}). After r Grover
iterations the success probability is
sin\ensuremath{{}^{2}}((2r+1)\ensuremath{\theta}). Choose r as a nearest
nonnegative integer to
\ensuremath{\pi}/(4\ensuremath{\theta})\ensuremath{-}1/2, evaluating
neighboring integers and handling small N explicitly. Derive the query
upper bound and pair it with a hybrid/adversary/polynomial lower bound
\ensuremath{\Omega}(\ensuremath{\surd}N) under the same promise.

An oracle separation must be written
BQP\textsuperscript{O\ensuremath{\ne}BPP}O for a specified O; it does
not establish BQP\ensuremath{\ne}BPP. A linear-system extension must
state matrix access, \ensuremath{\kappa}, precision, state preparation,
and whether the output is a quantum state or classical vector.

\subsection{Required outputs and proof obligations}

Return oracle/circuit specification, exact small-instance results, an
all-N success derivation, query/gate/resource accounting, and the
lower-bound proof with assumptions. Any QAOA extension must distinguish
a certified bound on a finite graph from a performance theorem for a
graph family.

\subsection{Validation and rejection controls}

Check the analytic success formula on n=1,\ldots,8 and all marked
positions using symmetry or exhaustive checks. Test neighboring
iteration counts to expose overshoot. Verify unitarity and
normalization. Compare with randomized classical search under identical
oracle access and error tolerance.

\subsection{Milestones and research extension}

Baseline: Grover construction and matched query bound. Strong: one
second promise problem with matched assumptions and a certified
upper/lower bracket. Research extension: a new restricted query bound or
oracle result. Reimplementation of established algorithms is
replication, not a new complexity separation.

Milestones are acceptance gates, not calendar promises. Before a run,
freeze the concrete instance/domain, tolerance, compute budget, and
permitted dependencies. A bounded search that fails to find a result
must preserve an unresolved outcome.

\subsection{Starting instruction}

Work on the baseline above. First make every input and convention
explicit, then validate the reference case. Build the checker alongside
the calculation. Label numerical evidence, finite-model certificates,
and continuum theorems separately; return the strongest justified result
and identify the exact unmet proof obligation.

\subsection{Primary references and their role}

\begin{itemize}
\tightlist
\item
  \href{https://arxiv.org/abs/quant-ph/9701001}{Bennett et al.,
  Strengths and Weaknesses of Quantum Computing} --- Oracle lower-bound
  methodology for quantum search.
\end{itemize}

References support definitions or limitations; the scope choices and
proposed milestones are editorial recommendations, not claims that these
results have already been obtained.

\subsection{Displayed formulation and claim boundary}

\href{https://github.com/rick-stevens-ai/Pure-Thought/blob/main/PRDs/25-Quantum-Algorithms-Complexity.md}{Canonical
specification}.

\[\theta=\arcsin(N^{-1/2}),\qquad
p_r=\sin^2((2r+1)\theta),\qquad
r\approx\frac{\pi}{4\theta}-\frac12.\]

Choose a nonnegative integer by checking neighboring counts and
explicitly handle small instances under the declared oracle model. The
target is a matched \(O(\sqrt N)\) algorithm and \(\Omega(\sqrt N)\)
lower bound with identical promise and error tolerance. An oracle
separation has the form \(\mathrm{BQP}^O\ne\mathrm{BPP}^O\) and does not
prove \(\mathrm{BQP}\ne\mathrm{BPP}\).

\begin{center}\rule{0.5\linewidth}{0.5pt}\end{center}

\section{26. A posteriori certification of an invariant Hamiltonian
torus}

\textbf{Revised specification --- 2026-09-11.}
\href{https://github.com/rick-stevens-ai/Pure-Thought/blob/main/reviews/26-KAM-Theory-Planetary-Stability.md}{Detailed
critique} \ensuremath{\cdot}
\href{https://github.com/rick-stevens-ai/Pure-Thought/blob/main/archive/original-PRDs/26-KAM-Theory-Planetary-Stability.md}{Original
description} \ensuremath{\cdot}
\href{https://github.com/rick-stevens-ai/Pure-Thought/blob/main/reviews/EVIDENCE-STANDARD.md}{Evidence
standard}

\subsection{Plain-language guide}

\subsubsection{The problem in everyday terms}

A weakly disturbed mechanical system can sometimes keep moving in an
orderly pattern forever, even though the motion never exactly repeats.
KAM theory explains how certain such motions survive a perturbation.
This problem starts with a proposed orderly motion and checks, with
rigorous error bounds, whether a true one must exist nearby. The goal is
a guarantee about a specified mathematical orbit family, not a
prediction that every nearby trajectory is stable.

\subsubsection{Key terms}

\begin{itemize}
\item
  \textbf{Hamiltonian} --- An energy function that generates the
  system's equations of motion.
\item
  \textbf{Integrable system} --- A system with enough conserved
  structure to describe its motion using constant actions and evolving
  angles in the setting considered here.
\item
  \textbf{Action--angle variables} --- Coordinates in which the actions
  label the motion and the angles describe its cyclic phases.
\item
  \textbf{Invariant torus} --- A torus-shaped set in phase space that
  trajectories starting on it never leave; it need not be a physical
  doughnut in ordinary space.
\item
  \textbf{Quasiperiodic motion} --- Motion combining incommensurate
  frequencies, so it remains organized without repeating exactly.
\item
  \textbf{KAM theory} --- Kolmogorov--Arnold--Moser theory, describing
  persistence of suitable quasiperiodic motions under sufficiently small
  perturbations.
\item
  \textbf{Diophantine condition} --- An arithmetic bound keeping integer
  combinations of the frequencies from becoming too close to zero.
\item
  \textbf{A posteriori theorem} --- A theorem that uses a computed
  approximate solution and its error to establish a nearby exact
  solution.
\item
  \textbf{Residual} --- The amount by which the approximate solution
  fails to satisfy the exact equation.
\item
  \textbf{Interval arithmetic} --- Computation with enclosing intervals,
  using controlled rounding to bound numerical errors.
\end{itemize}

\subsubsection{Why this matters}

Long-term stability is difficult to infer from a finite trajectory
simulation. A theorem can instead prove that a particular organized
motion exists despite perturbations. This offers a foundation for
studying stability mechanisms in celestial and other Hamiltonian
systems, provided the model and theorem assumptions are appropriate.

\subsubsection{What progress would mean}

A certificate for a nonzero perturbation would convert an approximate
torus into an existence result. Applying the method to planetary
observations would additionally require reduction of the planetary model
and a careful treatment of measured initial-condition uncertainty.

\subsection{Core question}

For an explicitly specified analytic near-integrable Hamiltonian and an
approximate torus, can every hypothesis of a quantitative KAM theorem be
verified?

\subsection{Scope and assumptions}

Begin with
H(I,\ensuremath{\theta})=\ensuremath{\omega}\ensuremath{\cdot}I+|I|\ensuremath{{}^{2}}/2+\ensuremath{\epsilon}[cos
\ensuremath{\theta}\ensuremath{{}_{1}}+cos(\ensuremath{\theta}\ensuremath{{}_{1}}\ensuremath{-}\ensuremath{\theta}\ensuremath{{}_{2}})]
in two degrees of freedom,
\ensuremath{\omega}=(1,(1+\ensuremath{\surd}5)/2), on a stated complex
neighborhood of a real action box. Choose a rational
\ensuremath{\epsilon} only after validating the unperturbed case. Seek
an embedding with rotation vector \ensuremath{\omega}.

This is a mathematical model, not a solar-system stability assertion.
Planetary extensions require an explicit canonical reduction, handling
Kepler degeneracy, collision exclusions, and input uncertainty.

\subsection{Mathematical target}

For an approximate embedding K:T\ensuremath{{}^{2}}\ensuremath{\to}phase
space define
E(\ensuremath{\theta})=X\_H(K(\ensuremath{\theta}))\ensuremath{-}DK(\ensuremath{\theta})\ensuremath{\omega}.
Select a published a posteriori KAM theorem and reproduce its hypotheses
and computable constants: analytic widths, nondegeneracy/twist inverse,
residual norm, and small-divisor bounds.

Prove
|k\ensuremath{\cdot}\ensuremath{\omega}|\ensuremath{\ge}\ensuremath{\gamma}/|k|\ensuremath{{}_{1}}\textasciicircum{}\ensuremath{\tau}
for every nonzero integer k using arithmetic properties of the exact
chosen \ensuremath{\omega}. Finite checks alone are insufficient.
Evaluate the theorem's smallness inequalities with outward-rounded
intervals and certify closeness of a true K\_* to K.

\subsection{Required outputs and proof obligations}

Export H, \ensuremath{\epsilon}, domain, exact frequency representation,
Fourier coefficients, all theorem constants, interval inequalities, and
a checker. State the conclusion: existence of a nearby quasiperiodic
invariant torus with a quantified embedding error. A failed sufficient
condition yields ``not certified,'' not ``unstable.''

\subsection{Validation and rejection controls}

Verify E=0 for
K\ensuremath{{}_{0}}(\ensuremath{\theta})=(I=0,\ensuremath{\theta}) at
\ensuremath{\epsilon}=0. Derive an all-integer Diophantine bound for the
golden-ratio frequency. Check Fourier tails and loss of analytic width.
Test a resonant frequency where the stated theorem must refuse
certification.

\subsection{Milestones and research extension}

Baseline: unperturbed certificate and complete arithmetic/theorem setup.
Strong: a nonzero-\ensuremath{\epsilon} torus certificate. Research
extension: a larger perturbation range or a reduced planetary model. One
invariant torus does not certify a neighborhood of arbitrary observed
initial conditions or a phase-space measure bound.

Milestones are acceptance gates, not calendar promises. Before a run,
freeze the concrete instance/domain, tolerance, compute budget, and
permitted dependencies. A bounded search that fails to find a result
must preserve an unresolved outcome.

\subsection{Starting instruction}

Work on the baseline above. First make every input and convention
explicit, then validate the reference case. Build the checker alongside
the calculation. Label numerical evidence, finite-model certificates,
and continuum theorems separately; return the strongest justified result
and identify the exact unmet proof obligation.

\subsection{Primary references and their role}

\begin{itemize}
\tightlist
\item
  \href{https://arxiv.org/abs/2101.07785}{Valvo and Locatelli,
  Hamiltonian Control of Magnetic Field Lines: Computer Assisted Results
  Proving the Existence of KAM Barriers} --- Example of a complete
  computer-assisted Hamiltonian KAM workflow; use a theorem matching the
  chosen formulation.
\end{itemize}

References support definitions or limitations; the scope choices and
proposed milestones are editorial recommendations, not claims that these
results have already been obtained.

\subsection{Displayed formulation and claim boundary}

\href{https://github.com/rick-stevens-ai/Pure-Thought/blob/main/PRDs/26-KAM-Theory-Planetary-Stability.md}{Canonical
specification}.

\[H(I,\theta)=\omega\cdot I+\frac{|I|^2}{2}
+\epsilon[\cos\theta_1+\cos(\theta_1-\theta_2)],
\qquad\omega=(1,(1+\sqrt5)/2),\]

\[E(\theta)=X_H(K(\theta))-DK(\theta)\omega,
\qquad |k\cdot\omega|\ge\frac\gamma{|k|_1^\tau}
\quad(k\in\mathbb Z^2\setminus\{0\}).\]

Prove the all-integer Diophantine bound and every analytic-width, twist,
tail and smallness condition of a specified quantitative theorem. Its
conclusion is a nearby invariant torus, not stability of arbitrary
planetary initial conditions. Failure of a sufficient criterion is not
instability.

\begin{center}\rule{0.5\linewidth}{0.5pt}\end{center}

\section{27. Central configurations in a bounded, symmetry-reduced
N-body problem}

\textbf{Revised specification --- 2026-09-11.}
\href{https://github.com/rick-stevens-ai/Pure-Thought/blob/main/reviews/27-N-Body-Central-Configurations.md}{Detailed
critique} \ensuremath{\cdot}
\href{https://github.com/rick-stevens-ai/Pure-Thought/blob/main/archive/original-PRDs/27-N-Body-Central-Configurations.md}{Original
description} \ensuremath{\cdot}
\href{https://github.com/rick-stevens-ai/Pure-Thought/blob/main/reviews/EVIDENCE-STANDARD.md}{Evidence
standard}

\subsection{Plain-language guide}

\subsubsection{The problem in everyday terms}

Most collections of gravitating bodies change their overall shape as
they move. Certain special arrangements can rotate or expand and
contract while keeping the same shape. These are central configurations.
This problem asks for every such arrangement in a carefully limited
case, with a proof that the list is complete and a separate analysis of
what small disturbances do to the associated rotating motion.

\subsubsection{Key terms}

\begin{itemize}
\item
  \textbf{N-body problem} --- The problem of describing the motion of N
  mutually gravitating bodies.
\item
  \textbf{Central configuration} --- An arrangement in which each body's
  gravitational acceleration is proportional to its position relative to
  the center of mass, with one common proportionality factor.
\item
  \textbf{Homographic motion} --- Motion preserving the arrangement's
  shape while its size and orientation may change.
\item
  \textbf{Relative equilibrium} --- A motion that is stationary in an
  appropriate rotating reference frame.
\item
  \textbf{Center of mass} --- The mass-weighted average position, used
  here to remove overall translation.
\item
  \textbf{Normalization} --- A fixed choice of scale, such as setting
  the moment of inertia to one, so scaled copies are not counted
  separately.
\item
  \textbf{Symmetry quotient} --- The convention for treating
  configurations related by rotations, reflections or specified
  relabelings as equivalent.
\item
  \textbf{Collision-free} --- Having strictly positive separation
  between every pair of bodies.
\item
  \textbf{Coriolis term} --- A velocity-dependent term arising in a
  rotating frame, essential to the stability calculation.
\item
  \textbf{Spectral stability} --- A statement about eigenvalues of the
  linearized dynamics; it is weaker than a full nonlinear stability
  proof.
\end{itemize}

\subsubsection{Why this matters}

Central configurations are organizing structures in gravitational
dynamics and provide exact reference cases amid otherwise complicated
motion. A complete small catalog can test numerical methods and reveal
how mass choices and symmetry affect possible motions. The problem also
links gravity to exact algebraic root finding.

\subsubsection{What progress would mean}

Progress would establish a complete list under a stated mass choice and
symmetry convention. It would not classify every periodic orbit: a
choreography such as the figure-eight changes shape and belongs to a
different question.

\subsection{Core question}

For fixed positive masses and a declared symmetry quotient, can all
collision-free central configurations in the chosen family be certified
and their relative-equilibrium linearization analyzed?

\subsection{Scope and assumptions}

Start with three labeled positive rational masses in the plane. Fix G=1,
center of mass zero, moment of inertia
I=\ensuremath{\Sigma}m\_i|r\_i|\ensuremath{{}^{2}}=1, and quotient by
O(2), so mirror images count together. Permute only equal masses when
explicitly requested. Extend to a specified four-body symmetry family
after validating the complete three-body case.

Central configurations generate homographic motions; general periodic
choreographies are a separate problem.

\subsection{Mathematical target}

Use positive
U=\ensuremath{\Sigma}\_\{i<j\}m\_im\_j/|r\_i\ensuremath{-}r\_j| and

\[\nabla_{r_i}U=-\lambda m_ir_i,\qquad \lambda =U/I>0.\]

Introduce distance variables d\_ij\textgreater0 with
d\_ij\ensuremath{{}^{2}}=|r\_i\ensuremath{-}r\_j|\ensuremath{{}^{2}},
and retain positivity/collision exclusions when polynomializing. Use
symmetry charts covering every configuration claimed. Completeness
requires certified global domain coverage or an exact
elimination/root-count argument, not just isolated numerical solutions.

\subsection{Required outputs and proof obligations}

Return normalized configurations, distance enclosures, residual
identities, existence/uniqueness proofs, equivalence representatives,
and completeness scope. For planar relative equilibria analyze the
reduced rotating-frame first-order system including Coriolis terms.
Distinguish spectral stability from nonlinear stability.

\subsection{Validation and rejection controls}

For three labeled masses under the stated O(2) quotient, recover three
collinear orderings (one per middle body) and one equilateral shape.
Check \ensuremath{\lambda}=U at I=1. Test collision artifacts introduced
by denominator clearing. For stability, reproduce a known three-body
criterion using identical mass and reduction conventions.

\subsection{Milestones and research extension}

Baseline: complete certified three-body catalog. Strong: complete
enumeration in one four-body symmetry family or mass instance, with all
charts covered. Research extension: a five-body restricted family; no
claim to solve the general finiteness problem without an appropriately
global proof.

Milestones are acceptance gates, not calendar promises. Before a run,
freeze the concrete instance/domain, tolerance, compute budget, and
permitted dependencies. A bounded search that fails to find a result
must preserve an unresolved outcome.

\subsection{Starting instruction}

Work on the baseline above. First make every input and convention
explicit, then validate the reference case. Build the checker alongside
the calculation. Label numerical evidence, finite-model certificates,
and continuum theorems separately; return the strongest justified result
and identify the exact unmet proof obligation.

\subsection{Primary references and their role}

\begin{itemize}
\tightlist
\item
  \href{https://arxiv.org/abs/2103.08906}{Yu and Zhu, On the finiteness
  of four-body central configurations} --- Context for finiteness versus
  explicit certified enumeration.
\end{itemize}

References support definitions or limitations; the scope choices and
proposed milestones are editorial recommendations, not claims that these
results have already been obtained.

\subsection{Displayed formulation and claim boundary}

\href{https://github.com/rick-stevens-ai/Pure-Thought/blob/main/PRDs/27-N-Body-Central-Configurations.md}{Canonical
specification}.

\[U=\sum_{i<j}\frac{m_im_j}{|r_i-r_j|},\quad
\sum_i m_i r_i=0,\quad I=\sum_i m_i|r_i|^2=1,\]

\[\nabla_{r_i}U=-\lambda m_i r_i,\qquad\lambda=U/I>0,
\qquad d_{ij}^2=|r_i-r_j|^2,\quad d_{ij}>0.\]

Preserve collision exclusions when clearing denominators. For three
labeled positive masses modulo \(O(2)\), recover three collinear
orderings and one equilateral shape. Complete coverage requires more
than isolated root boxes. Relative-equilibrium stability uses the
reduced rotating-frame system including Coriolis terms; a figure-eight
choreography is not a central configuration.

\begin{center}\rule{0.5\linewidth}{0.5pt}\end{center}

\section{28. Explicit finite-time action confinement from a normal form}

\textbf{Revised specification --- 2026-09-11.}
\href{https://github.com/rick-stevens-ai/Pure-Thought/blob/main/reviews/28-Nekhoroshev-Stability-Theory.md}{Detailed
critique} \ensuremath{\cdot}
\href{https://github.com/rick-stevens-ai/Pure-Thought/blob/main/archive/original-PRDs/28-Nekhoroshev-Stability-Theory.md}{Original
description} \ensuremath{\cdot}
\href{https://github.com/rick-stevens-ai/Pure-Thought/blob/main/reviews/EVIDENCE-STANDARD.md}{Evidence
standard}

\subsection{Plain-language guide}

\subsubsection{The problem in everyday terms}

A system need not move on a perfectly orderly orbit to remain near its
initial state for a very long time. Small disturbances can produce only
a slow drift in certain quantities. This problem asks for an explicit
guarantee: how far can those quantities change, and for how long is that
bound valid? It focuses on a finite time horizon rather than proving
that a special motion persists forever.

\subsubsection{Key terms}

\begin{itemize}
\item
  \textbf{Near-integrable Hamiltonian} --- An energy function consisting
  of a well-understood integrable part plus a small perturbation.
\item
  \textbf{Action} --- A coordinate that is constant in the integrable
  reference system and may drift when a perturbation is added.
\item
  \textbf{Confinement bound} --- A guaranteed maximum change of the
  specified actions over a stated time interval.
\item
  \textbf{Normal form} --- A transformed expression of the equations
  separating manageable terms from a remainder.
\item
  \textbf{Canonical transformation} --- A change of variables preserving
  the Hamiltonian structure of the dynamics.
\item
  \textbf{Remainder} --- Terms not absorbed into the simplified normal
  form; bounding their size helps bound drift.
\item
  \textbf{Resonance} --- An integer relation among frequencies that can
  change how perturbations affect the motion.
\item
  \textbf{Nekhoroshev theorem} --- A theorem giving very long
  action-confinement times for small perturbations under specified
  regularity and nondegeneracy assumptions.
\item
  \textbf{Quasi-convexity} --- A particular curvature condition on the
  integrable energy function used by some such theorems.
\item
  \textbf{Dimensional time} --- A time expressed in declared physical or
  model units, rather than an unexplained numerical factor.
\end{itemize}

\subsubsection{Why this matters}

Finite-time guarantees answer a question that simulations alone cannot
settle beyond their run length. They also distinguish a valid stability
argument from an impressive-looking exponential formula with missing
constants or units. Understanding slow drift is relevant to the
mathematical study of planetary and other nearly integrable systems.

\subsubsection{What progress would mean}

A useful first result could be a modest but fully justified confinement
time. Improving that time with controlled remainders would be real
progress; asserting billions of years would require an appropriate
physical model and every conversion and bound needed to support it.

\subsection{Core question}

For a specified analytic Hamiltonian and action domain, what rigorously
computable action excursion and confinement time follow from a certified
normal form or quantitative Nekhoroshev theorem?

\subsection{Scope and assumptions}

Use H=h(I)+\ensuremath{\epsilon}f(I,\ensuremath{\theta}),
h=\ensuremath{\omega}\ensuremath{\cdot}I+|I|\ensuremath{{}^{2}}/2, in
two or three degrees of freedom on a compact action box with an explicit
complex extension. Supply \ensuremath{\omega}, a finite Fourier
polynomial f, and rational \ensuremath{\epsilon}. Bound the gradient
away from zero where the selected theorem requires it.

The initial target is a valid dimensional time bound, even if modest.
Planetary and post-Newtonian models are extensions requiring their own
reductions and hypotheses.

\subsection{Mathematical target}

For a quasi-convex theorem verify its actual condition, for example
v\ensuremath{{}^{T}}D\ensuremath{{}^{2}}h(I)v\ensuremath{\ge}m|v|\ensuremath{{}^{2}}
for v\ensuremath{\perp}\ensuremath{\nabla}h(I), plus the theorem's
gradient, domain and analyticity conditions. A positive Hessian
determinant alone is insufficient.

Alternatively construct a canonical \ensuremath{\Phi} with
H\ensuremath{\circ}\ensuremath{\Phi}=Z(J)+R(J,\ensuremath{\theta}),
bound |I\ensuremath{-}J|\ensuremath{\le}\ensuremath{\eta} and
||\ensuremath{\partial}\_\ensuremath{\theta}R||\ensuremath{\le}r on a
validated domain. Then
|I(t)\ensuremath{-}I(0)|\ensuremath{\le}2\ensuremath{\eta}+r|t| while
the transformed orbit remains in that domain. If resonant Z depends on
angles, include its action drift or restrict the claim to
conserved/protected components.

\subsection{Required outputs and proof obligations}

Export normal-form order, transformations, inverse/domain bounds,
Fourier tails, r, \ensuremath{\eta}, and a domain-containment proof for
the claimed interval. A Nekhoroshev result must include all constants
and its smallness threshold in T=C exp(c
\ensuremath{\epsilon}\textasciicircum{}\ensuremath{-}a) with stated
units; do not choose a from a desired time target.

\subsection{Validation and rejection controls}

Verify integrable \ensuremath{\epsilon}=0 behavior and a low-order
symbolic normal form. Check symplecticity with validated remainder
control. Confirm the action bound stays inside the certified domain.
Independently recompute
exp(10\textasciicircum{}0.3)\ensuremath{\approx}7.35: it is a
dimensionless factor, not 10\ensuremath{{}^{13}} years.

\subsection{Milestones and research extension}

Baseline: one explicit finite-time confinement certificate. Strong:
optimized truncation with a stronger bound or full quantitative
quasi-convex theorem application. Research extension: exponential
scaling across a parameter family; do not infer actual solar-system
stability from \ensuremath{\epsilon} alone.

Milestones are acceptance gates, not calendar promises. Before a run,
freeze the concrete instance/domain, tolerance, compute budget, and
permitted dependencies. A bounded search that fails to find a result
must preserve an unresolved outcome.

\subsection{Starting instruction}

Work on the baseline above. First make every input and convention
explicit, then validate the reference case. Build the checker alongside
the calculation. Label numerical evidence, finite-model certificates,
and continuum theorems separately; return the strongest justified result
and identify the exact unmet proof obligation.

\subsection{Primary references and their role}

\begin{itemize}
\tightlist
\item
  \href{https://arxiv.org/abs/1701.06026}{Zhang and Zhang, Improved
  stability for analytic quasi-convex nearly integrable systems} ---
  Theorem-specific quasi-convex exponential stability results.
\end{itemize}

References support definitions or limitations; the scope choices and
proposed milestones are editorial recommendations, not claims that these
results have already been obtained.

\subsection{Displayed formulation and claim boundary}

\href{https://github.com/rick-stevens-ai/Pure-Thought/blob/main/PRDs/28-Nekhoroshev-Stability-Theory.md}{Canonical
specification}.

\[H\circ\Phi=Z(J)+R(J,\theta),\qquad
|I-J|\le\eta,\qquad\|\partial_\theta R\|\le r,\]

\[|I(t)-I(0)|\le2\eta+r|t|.\]

This estimate assumes the angle-independent normal form shown and a
proved domain-containment interval. Include resonant drift if \(Z\)
depends on angles. A quantitative Nekhoroshev theorem instead requires
all its constants, hypotheses and threshold in
\(T=C\exp(c\epsilon^{-a})\), with units supplied by \(C\). The number
\(\exp(10^{0.3})\approx7.35413\) has no unit of time.

\begin{center}\rule{0.5\linewidth}{0.5pt}\end{center}

\section{29. Autocatalytic sets: maximality, irreducibility and
dynamical viability}

\textbf{Revised specification --- 2026-09-11.}
\href{https://github.com/rick-stevens-ai/Pure-Thought/blob/main/reviews/29-Chemical-Reaction-Networks-Origins.md}{Detailed
critique} \ensuremath{\cdot}
\href{https://github.com/rick-stevens-ai/Pure-Thought/blob/main/archive/original-PRDs/29-Chemical-Reaction-Networks-Origins.md}{Original
description} \ensuremath{\cdot}
\href{https://github.com/rick-stevens-ai/Pure-Thought/blob/main/reviews/EVIDENCE-STANDARD.md}{Evidence
standard}

\subsection{Plain-language guide}

\subsubsection{The problem in everyday terms}

Could a collection of reactions collectively help produce the substances
that make those same reactions happen? Autocatalytic-set theory turns
part of that question into a precise network test. Starting from an
externally supplied food set, we ask which reactions can be supported by
products of the network and which smaller subsets already have that
property. This describes a structural prerequisite for self-sustaining
chemistry, not a complete definition of life.

\subsubsection{Key terms}

\begin{itemize}
\item
  \textbf{Catalyst} --- A substance participating in a reaction
  mechanism that increases its rate while being regenerated overall; it
  does not change the equilibrium free-energy difference.
\item
  \textbf{Autocatalysis} --- Production that directly or collectively
  helps enable more production of the relevant system.
\item
  \textbf{Food set} --- The species assumed to be available from outside
  the network.
\item
  \textbf{Closure} --- All species producible from the food set using
  the selected reaction rules; standard RAF closure temporarily ignores
  catalysis.
\item
  \textbf{RAF} --- A reflexively autocatalytic and food-generated
  reaction set: its reactants are reachable from food and each reaction
  has a catalyst available within the closure.
\item
  \textbf{Maximal RAF} --- The largest RAF contained in the supplied
  network, when a RAF exists.
\item
  \textbf{Irreducible RAF} --- A RAF containing no smaller nonempty RAF
  subset.
\item
  \textbf{Minimum-cardinality RAF} --- A RAF with the fewest reactions
  among all RAFs in the specified network; this is stronger than
  irreducibility.
\item
  \textbf{Chemostat} --- An idealized externally maintained
  concentration or chemical reservoir.
\item
  \textbf{Dynamical viability} --- The additional question of whether
  specified reaction rates and conditions support sustained activity or
  growth.
\end{itemize}

\subsubsection{Why this matters}

Collectively supported chemistry is relevant to thinking about how
organized reaction systems can arise and persist. Exact minimality
results help distinguish essential reactions from dispensable ones in a
model. Just as importantly, separating structural closure from energy
supply and kinetics prevents a network diagram from being mistaken for
proof of a viable replicating system.

\subsubsection{What progress would mean}

Success would identify supported reaction sets and prove the stated kind
of minimality. Demonstrating sustained growth, heredity or evolution
would require extra chemical and dynamical assumptions beyond the RAF
certificate.

\subsection{Core question}

For a finite catalytic reaction system, can its maximal RAF and smallest
RAF subsets be certified, while keeping structural autocatalysis
distinct from sustained growth?

\subsection{Scope and assumptions}

Input species S, reactions R with integer stoichiometry, food set F, and
catalysis relation C\ensuremath{\subset}S\ensuremath{\times}R. Begin
with at most 12 reactions for exhaustive cross-checks. Catalysis is
supplied model data, not inferred from a stoichiometric cycle.

Define three separate targets: unique maximal RAF, an inclusion-minimal
(irreducible) RAF, and a minimum-cardinality RAF. Thermodynamic and
kinetic viability require extra inputs and are separate extensions.

\subsection{Mathematical target}

For R\ensuremath{\prime}\ensuremath{\subset}R compute
cl\_\{R\ensuremath{\prime}\}(F) by repeatedly adding products whose
reactants are already present, ignoring catalysis in this closure step.
R\ensuremath{\prime} is a nonempty RAF when every reactant lies in this
closure and every reaction has at least one catalyst in it.

An irreducible RAF has no proper nonempty RAF subset. A
minimum-cardinality RAF minimizes |R\ensuremath{\prime}| over all RAF
subsets. Ordinary RAF closure permits catalysts to arise later; a
catalysis-respecting startup ordering is a stronger condition and must
be tested separately.

\subsection{Required outputs and proof obligations}

Export closure rounds, catalyst witnesses, maxRAF elimination trace, and
minimality/cardinality proof with explicit scope. For thermodynamic
extension specify chemical potentials, chemostats, and driving. For
kinetic extension provide rates, dilution, initial conditions and a
positive sustained-state/growth criterion; RAF existence alone is
insufficient.

\subsection{Validation and rejection controls}

Cross-check all subsets for tiny networks. Include food-catalyzed,
self-catalyzed, catalyst-free, and mutually dependent examples. Verify a
maximum RAF can contain several irreducible RAFs of different sizes.
Deletion tests for irreducibility must use a complete RAF detector on
each remaining network, not merely test whether the entire remainder is
itself RAF.

\subsection{Milestones and research extension}

Baseline: checked maxRAF/RAF witnesses. Strong: complete
minimum-cardinality results for a bounded network family. Research
extension: characterize which RAFs additionally support a declared
driven mass-action regime. Biological self-replication and evolution
require their own operational definitions.

Milestones are acceptance gates, not calendar promises. Before a run,
freeze the concrete instance/domain, tolerance, compute budget, and
permitted dependencies. A bounded search that fails to find a result
must preserve an unresolved outcome.

\subsection{Starting instruction}

Work on the baseline above. First make every input and convention
explicit, then validate the reference case. Build the checker alongside
the calculation. Label numerical evidence, finite-model certificates,
and continuum theorems separately; return the strongest justified result
and identify the exact unmet proof obligation.

\subsection{Primary references and their role}

\begin{itemize}
\tightlist
\item
  \href{https://pmc.ncbi.nlm.nih.gov/articles/PMC4034171/}{Hordijk and
  Steel, Autocatalytic sets in a partitioned biochemical network} ---
  RAF closure and the distinction between maximal and irreducible RAFs.
\end{itemize}

References support definitions or limitations; the scope choices and
proposed milestones are editorial recommendations, not claims that these
results have already been obtained.

\subsection{Displayed formulation and claim boundary}

\href{https://github.com/rick-stevens-ai/Pure-Thought/blob/main/PRDs/29-Chemical-Reaction-Networks-Origins.md}{Canonical
specification}.

\[R'\subseteq R,\quad R'\ne\varnothing,\quad
\operatorname{reactants}(r)\subseteq\operatorname{cl}_{R'}(F),\]

\[\forall r\in R'\ \exists s\in\operatorname{cl}_{R'}(F):\ (s,r)\in C.\]

Compute closure by reactant availability without imposing catalysis in
the closure step. These conditions define a RAF. Inclusion-maximal,
inclusion-minimal (irreducible), and minimum-cardinality RAFs are
separate targets. Catalysis-respecting startup and sustained growth
require additional conditions.

\begin{center}\rule{0.5\linewidth}{0.5pt}\end{center}

\section{30. An exact finite genotype--phenotype map and its mutation
graph}

\textbf{Revised specification --- 2026-09-11.}
\href{https://github.com/rick-stevens-ai/Pure-Thought/blob/main/reviews/30-Genotype-Phenotype-Mapping.md}{Detailed
critique} \ensuremath{\cdot}
\href{https://github.com/rick-stevens-ai/Pure-Thought/blob/main/archive/original-PRDs/30-Genotype-Phenotype-Mapping.md}{Original
description} \ensuremath{\cdot}
\href{https://github.com/rick-stevens-ai/Pure-Thought/blob/main/reviews/EVIDENCE-STANDARD.md}{Evidence
standard}

\subsection{Plain-language guide}

\subsubsection{The problem in everyday terms}

Different sequences of genetic letters can produce the same modeled
structure. That creates paths through sequence space along which
mutations leave the chosen outcome unchanged. This problem constructs a
small RNA-inspired map exactly and asks how those paths connect, how
often one mutation changes the outcome, and which new structures are
nearby. The aim is to study a controlled model of robustness and
variation.

\subsubsection{Key terms}

\begin{itemize}
\item
  \textbf{Genotype} --- The input sequence, here a string over the four
  RNA letters A, C, G and U.
\item
  \textbf{Phenotype} --- The selected output of the model, here a
  secondary-structure pattern rather than an organism's full observable
  traits.
\item
  \textbf{Secondary structure} --- A pattern of paired positions along
  an RNA sequence, with the allowed pairings and restrictions stated
  explicitly.
\item
  \textbf{Nussinov-style model} --- A combinatorial folding model
  maximizing the number of allowed pairs, not a full free-energy
  calculation.
\item
  \textbf{Tie-breaking rule} --- A deterministic rule choosing one
  output when several structures achieve the same score.
\item
  \textbf{Hamming graph} --- A graph whose vertices are sequences and
  whose edges join sequences differing at exactly one position.
\item
  \textbf{Neutral component} --- A connected group of sequences with the
  same modeled phenotype.
\item
  \textbf{Mutational robustness} --- The fraction of a genotype's
  one-letter changes that preserve its phenotype under the chosen map.
\item
  \textbf{Evolvability} --- Here, a specified count of different
  phenotypes accessible by mutation, not a complete measure of
  biological adaptability.
\item
  \textbf{Fitness} --- A separately defined measure of reproductive
  success or a model score; it is not determined automatically by the
  folding map.
\end{itemize}

\subsubsection{Why this matters}

The structure of a genotype--phenotype map helps explain how variation
can accumulate without immediately changing a selected outcome, and how
new outcomes become accessible. Exact small maps provide benchmarks for
sampling methods and expose how conclusions depend on folding rules and
tie choices.

\subsubsection{What progress would mean}

A complete short-sequence graph would establish precise robustness and
connectivity facts for the model. Those facts could motivate biological
hypotheses, but would not by themselves predict RNA function, population
evolution or a universal transition at some sequence length.

\subsection{Core question}

For a fully specified combinatorial RNA folding rule on short sequences,
what are the exact neutral components, mutation robustness, and
accessible phenotype counts?

\subsection{Scope and assumptions}

Use alphabet \{A,C,G,U\}, sequence length L=4,\ldots,8 initially,
allowed pairs AU/UA, GC/CG and GU/UG, minimum hairpin length three
unpaired bases, no pseudoknots, and one pair per base. Maximize pair
count and break ties by lexicographically smallest dot-bracket string
under a declared character order.

This defines a mathematical GP map, not a thermodynamic RNA predictor. A
nearest-neighbor free-energy extension must name its parameter set;
empirical parameters are then acknowledged model inputs.

\subsection{Mathematical target}

Let \ensuremath{\Phi} map each sequence to its deterministic selected
structure. On the Hamming graph H(L,4), exact-one-mutation neighbors
satisfy d\_H(s,t)=1 and each sequence has 3L neighbors. Define

\[r(s)=\frac{|\{t:d_H(s,t)=1,\ \Phi (t)=\Phi (s)\}|}{3L}.\]

For each phenotype analyze the induced neutral graph and every connected
component; a phenotype preimage need not be connected. Define
evolvability as the number of distinct other phenotypes reachable by one
mutation, stating whether it is measured per genotype, component or
whole preimage.

\subsection{Required outputs and proof obligations}

Export folding rules, tie-breaking, sequence-to-structure table, neutral
components, exact counts and rational robustness. A search restricted to
a component must not report the full phenotype frequency. Fitness is a
separately supplied function; evolutionary dynamics require population
size, mutation kernel, selection rule and initial state.

\subsection{Validation and rejection controls}

Cross-check dynamic programming with exhaustive legal-structure
enumeration on short sequences, especially ties and prohibited pairs.
Verify
\ensuremath{\Sigma}\_\ensuremath{\phi}|\ensuremath{\Phi}\ensuremath{{}^{-1}}(\ensuremath{\phi})|=4\textasciicircum{}L
and every Hamming vertex degree is 3L. Count neutral edges independently
and recover average robustness from twice the edge count. Treat any
larger-L sample as statistical evidence.

\subsection{Milestones and research extension}

Baseline: complete exact maps through L=8. Strong: increase L within
measured resources and characterize component-level
robustness/evolvability. Research extension: a proved asymptotic
property or a specified noisy GP-channel capacity problem. Do not
hard-code a universal percolation threshold or empirical target
correlation.

Milestones are acceptance gates, not calendar promises. Before a run,
freeze the concrete instance/domain, tolerance, compute budget, and
permitted dependencies. A bounded search that fails to find a result
must preserve an unresolved outcome.

\subsection{Starting instruction}

Work on the baseline above. First make every input and convention
explicit, then validate the reference case. Build the checker alongside
the calculation. Label numerical evidence, finite-model certificates,
and continuum theorems separately; return the strongest justified result
and identify the exact unmet proof obligation.

\subsection{Primary references and their role}

\begin{itemize}
\tightlist
\item
  \href{https://www.tbi.univie.ac.at/RNA/RNAfold}{ViennaRNA, RNAfold
  documentation} --- RNAfold's energy model and parameter choices differ
  from a pair-count model.
\end{itemize}

References support definitions or limitations; the scope choices and
proposed milestones are editorial recommendations, not claims that these
results have already been obtained.

\subsection{Displayed formulation and claim boundary}

\href{https://github.com/rick-stevens-ai/Pure-Thought/blob/main/PRDs/30-Genotype-Phenotype-Mapping.md}{Canonical
specification}.

\[\Phi:\{A,C,G,U\}^L\longrightarrow\mathcal S_L,\qquad
r(s)=\frac{|\{t:d_H(s,t)=1,\ \Phi(t)=\Phi(s)\}|}{3L},\]

\[\sum_{\varphi}|\Phi^{-1}(\varphi)|=4^L,\qquad
\deg_{H(L,4)}s=3L.\]

Specify pair-count maximization, allowed pairs, hairpin length,
pseudoknot exclusion and deterministic tie-breaking. Enumerate every
connected component of each phenotype preimage. This combinatorial map
is not a thermodynamic predictor. At \(L=20\),
\(4^L=1{,}099{,}511{,}627{,}776\); sampling does not establish a
complete catalog.

\begin{center}\rule{0.5\linewidth}{0.5pt}\end{center}

\section{S01. Crystal-field and spin--orbit anisotropy in a finite ion
model}

\textbf{Source:} synthesis challenge 11;
\href{https://github.com/rick-stevens-ai/Pure-Thought/blob/main/NOTES/30-Most-Compelling-Pure-Thought-AI-Challenges.md}{original
synthesis}. This topic is distinct from, or adds a distinct objective
to, the canonical PRD catalog.
\href{https://github.com/rick-stevens-ai/Pure-Thought/blob/main/reviews/SYNTHESIS-CROSSWALK.md}{Crosswalk}

\subsection{Plain-language guide}

\subsubsection{The problem in everyday terms}

An ion inside a crystal can respond differently to a magnetic field
pointing in different directions. Its neighboring atoms shape the
electron's available states, while the electron's spin interacts with
its orbital motion. This problem asks whether a small, explicitly
specified quantum model can establish those directional differences with
reliable numerical bounds.

\subsubsection{Key terms}

\begin{itemize}
\item
  \textbf{Ion} --- An atom or group of atoms carrying net electric
  charge.
\item
  \textbf{Crystal field} --- The influence of surrounding charges on an
  ion's electronic energy levels.
\item
  \textbf{Spin} --- An electron's intrinsic quantum angular momentum.
\item
  \textbf{Orbital angular momentum} --- Quantum angular momentum
  associated with spatial electronic motion.
\item
  \textbf{Spin--orbit coupling} --- An interaction linking spin and
  orbital angular momentum.
\item
  \textbf{Hamiltonian} --- The operator specifying the model's energies
  and quantum evolution.
\item
  \textbf{Kramers doublet} --- A pair of equal-energy states protected
  by time-reversal symmetry in this half-integer-spin setting at zero
  magnetic field.
\item
  \textbf{g tensor} --- A matrix describing the directional magnetic
  response of the selected effective doublet.
\item
  \textbf{Excitation gap} --- The energy separation from the chosen
  low-energy states to the next states.
\end{itemize}

\subsubsection{Why this matters}

Directional magnetic response is a basic ingredient in understanding
magnetic materials and controlling spin states. A finite ion model lets
us separate the roles of local geometry and spin--orbit coupling before
tackling interactions among many ions.

\subsubsection{What progress would mean}

Progress would mean certified response ranges and excitation gaps over
the specified parameter region, with clear limits on when the effective
doublet description applies. This would establish a local mechanism;
predicting a bulk magnet would require additional interactions and
material information.

\subsection{Revised core question}

For a fixed d-shell occupancy and point-group representation, which
anisotropy values are possible in a bounded crystal-field/spin--orbit
Hamiltonian family?

\subsection{Precise scope, target and outputs}

Fix one ion model (start with d\ensuremath{{}^{1}}), a chosen point
group such as C\ensuremath{{}_{4}}v, and a finite orbital-spin basis.
Supply exact matrices for L,S and the group action. Define
H(B)=H\_CF+\ensuremath{\lambda}L\ensuremath{\cdot}S+\ensuremath{\mu}\_B
B\ensuremath{\cdot}(L+2S), with crystal-field coefficients and
\ensuremath{\lambda} in a compact normalized box. Specify the
temperature or isolated low-energy multiplet before defining
susceptibility or the effective g tensor. For a Kramers doublet, extract
a g tensor from the projected magnetic moment and use its singular
values, which remove doublet-basis ambiguity. Require a certified
excitation gap above that doublet. Optimize a bounded ratio or
difference with a denominator bounded away from zero. Export matrices,
symmetry identities, interval eigenvalue bounds, and global objective
brackets.

\subsection{Detailed critique of the synthesis framing}

The synthesis's ``all d\^{}n configurations across every point group''
has no interaction model, coefficient bounds, or objective
normalization. Symmetry permits tensor components but does not determine
their magnitudes; arbitrary crystal-field or spin--orbit scaling can
make an optimization meaningless. ``Rare-earth-like'' is a motivation,
not an observable. Many-electron d\^{}n ions also require Coulomb/Hund
parameters and a specified approximation. The revision retains an
anisotropy atlas but begins with a finite, normalized spectral problem
and separates a single-ion response from bulk magnetocrystalline
anisotropy.

\subsection{Validation and milestones}

Baseline: exact symmetry-allowed Hamiltonian basis and one doublet g
tensor. Strong: certified parameter-box extrema with a doublet-gap
condition. Extend to d\ensuremath{{}^{2}} only after specifying
interactions. Check the isotropic limit, group covariance, and basis
invariance of g singular values.

Use the shared
\href{https://github.com/rick-stevens-ai/Pure-Thought/blob/main/reviews/EVIDENCE-STANDARD.md}{evidence
standard}. Unsuccessful search is unresolved unless a complete exclusion
proof is supplied. These are proposed baseline specifications, not
completed scientific results.

\subsection{Displayed formulation and claim boundary}

\href{https://github.com/rick-stevens-ai/Pure-Thought/blob/main/supplemental/S01.md}{Supplemental
specification}.

\[H(B)=H_{\mathrm{CF}}+\lambda L\cdot S+\mu_B B\cdot(L+2S),\]

\[P H(B)P=E_0I+\frac{\mu_B}{2}\sum_{a,b}B_a g_{ab}\sigma_b+O(|B|^2).\]

Use the isolated Kramers-doublet projector \(P\) and exact orbital-spin
conventions. The singular values of \(g\) remove doublet-basis
ambiguity. Certify the gap above the doublet and bound any denominator
in an anisotropy ratio on the compact parameter box. This finite ion
model does not establish bulk magnetism.

\begin{center}\rule{0.5\linewidth}{0.5pt}\end{center}

\section{S02. Bandgap optimization in a bounded wave-operator family}

\textbf{Source:} synthesis challenge 12;
\href{https://github.com/rick-stevens-ai/Pure-Thought/blob/main/NOTES/30-Most-Compelling-Pure-Thought-AI-Challenges.md}{original
synthesis}. This topic is distinct from, or adds a distinct objective
to, the canonical PRD catalog.
\href{https://github.com/rick-stevens-ai/Pure-Thought/blob/main/reviews/SYNTHESIS-CROSSWALK.md}{Crosswalk}

\subsection{Plain-language guide}

\subsubsection{The problem in everyday terms}

Some repeating structures prevent waves in certain frequency ranges from
traveling through their interior. The challenge is to choose a structure
that makes this forbidden range as large as possible within stated
material and geometric limits. A gap must survive every propagation
direction, including directions missed by a convenient plotting path.

\subsubsection{Key terms}

\begin{itemize}
\item
  \textbf{Wave operator} --- The mathematical rule governing the waves
  in the chosen physical model.
\item
  \textbf{Periodic medium} --- A material pattern repeated through
  space.
\item
  \textbf{Unit cell} --- One repeating piece of that pattern.
\item
  \textbf{Bloch wave} --- A wave adapted to a periodic medium, labeled
  by a wave vector.
\item
  \textbf{Brillouin zone} --- The set of inequivalent wave vectors
  needed to describe a periodic medium.
\item
  \textbf{Band} --- A family of allowed wave frequencies or
  corresponding operator eigenvalues.
\item
  \textbf{Complete bandgap} --- A frequency interval containing no bulk
  wave modes for any wave vector in the specified model.
\item
  \textbf{Contrast} --- The permitted ratio or variation between
  material coefficients.
\item
  \textbf{Global optimum} --- The best value across the entire declared
  design family, rather than just nearby designs.
\end{itemize}

\subsubsection{Why this matters}

Controlling which waves propagate supports ideas for sound isolation,
vibration management and optical filtering. The same mathematical
challenge arises across these applications: a promising numerical design
needs reliable evidence that its gap exists everywhere and that any
optimality claim covers the stated design space.

\subsubsection{What progress would mean}

A successful result would certify a nonzero gap, bound the best
achievable objective, or establish an optimum within the bounded family.
A scalar-wave benchmark would test the method; applying it to
electromagnetic devices requires the appropriate Maxwell model and
physical constraints.

\subsection{Revised core question}

Within a fixed finite family of periodic media, what lower bound on a
complete bandgap and what upper bound on the best possible gap can be
certified?

\subsection{Precise scope, target and outputs}

Choose either a scalar acoustic operator or a full Maxwell operator,
never an unspecified mixture. Initially choose a periodic scalar
elliptic operator
\ensuremath{-}\ensuremath{\nabla}\ensuremath{\cdot}a(r)\ensuremath{\nabla}
with 0<a\_min\ensuremath{\le}a(r)\ensuremath{\le}a\_max, a fixed unit
cell and a finite geometry parameter box. For adjacent Bloch eigenvalues
define the complete gap g=min\_k
\ensuremath{\lambda}\_\{n+1\}(k)\ensuremath{-}max\_k
\ensuremath{\lambda}\_n(k), with a declared normalized objective. A
grid-local separation is not a complete gap. Use validated eigenvalue
bounds plus k-space variation estimates to cover the Brillouin zone.
Branch-and-bound over parameters yields a candidate lower bound and an
optimization upper bound, keeping the inequality direction explicit.

\subsection{Detailed critique of the synthesis framing}

The old description mixes photonic and phononic operators, index
contrast and generic gap optimization without fixing constitutive laws,
geometry regularity, or constraints. ``Lowest possible contrast'' and
``largest possible gap'' are two objectives needing an ordering or
Pareto definition. A converged plane-wave plot does not establish a
complete gap, and a locally optimized geometry does not prove global
optimality. This is distinct from PRD 11's topological invariant: a wide
ordinary bandgap can be topologically trivial.

\subsection{Validation and milestones}

Baseline: one certified complete gap for a supplied geometry. Strong:
bound the optimum to a declared tolerance over the entire compact
parameter box, or report unresolved boxes. Validate a homogeneous medium
and a simple one-dimensional layered case before multidimensional
optimization.

Use the shared
\href{https://github.com/rick-stevens-ai/Pure-Thought/blob/main/reviews/EVIDENCE-STANDARD.md}{evidence
standard}. Unsuccessful search is unresolved unless a complete exclusion
proof is supplied. These are proposed baseline specifications, not
completed scientific results.

\subsection{Displayed formulation and claim boundary}

\href{https://github.com/rick-stevens-ai/Pure-Thought/blob/main/supplemental/S02.md}{Supplemental
specification}.

\[\lambda_- =\max_k\lambda_n(k),\qquad
\lambda_+=\min_k\lambda_{n+1}(k),\qquad
g=\lambda_+-\lambda_-. \]

A possible normalized objective is \(G=2g/(\lambda_++\lambda_-)\) when
the denominator is positive; frequency-gap objectives must instead use
\(\omega_j=\sqrt{\lambda_j}\) consistently.

Use a fixed self-adjoint wave model, bounded coefficients and geometry
class, with a positive denominator. A complete gap needs
\(\lambda_+>\lambda_-\) across the entire Brillouin zone and all
required polarizations. A certified candidate lower bound does not
establish global optimality; that needs coverage of the full design
class.

\begin{center}\rule{0.5\linewidth}{0.5pt}\end{center}

\section{S03. Effective conductivity bounds for two-phase composites}

\textbf{Source:} synthesis challenge 13;
\href{https://github.com/rick-stevens-ai/Pure-Thought/blob/main/NOTES/30-Most-Compelling-Pure-Thought-AI-Challenges.md}{original
synthesis}. This topic is distinct from, or adds a distinct objective
to, the canonical PRD catalog.
\href{https://github.com/rick-stevens-ai/Pure-Thought/blob/main/reviews/SYNTHESIS-CROSSWALK.md}{Crosswalk}

\subsection{Plain-language guide}

\subsubsection{The problem in everyday terms}

Mix two materials that conduct heat or electricity differently. The
overall conductivity depends on how the materials are arranged, not just
how much of each is present. This problem asks for reliable upper and
lower bounds on the conductivity of a specified repeating arrangement,
and a careful distinction between its performance and limits that apply
to every arrangement.

\subsubsection{Key terms}

\begin{itemize}
\item
  \textbf{Conductivity} --- A coefficient relating a driving gradient to
  the resulting current or heat flux in the selected linear model.
\item
  \textbf{Two-phase composite} --- A material made from two constituent
  material types.
\item
  \textbf{Volume fraction} --- The fraction of space occupied by one
  constituent.
\item
  \textbf{Effective conductivity} --- The large-scale response obtained
  after averaging the microscopic structure.
\item
  \textbf{Conductivity tensor} --- A matrix allowing that response to
  depend on direction.
\item
  \textbf{Potential} --- A scalar field, such as electric potential,
  whose gradient drives transport.
\item
  \textbf{Flux} --- The local flow of the transported quantity per unit
  area.
\item
  \textbf{Divergence-free flux} --- A steady flow with no internal
  accumulation or sources.
\item
  \textbf{Variational bound} --- An inequality obtained by optimizing or
  testing an energy principle.
\end{itemize}

\subsubsection{Why this matters}

Composite design seeks useful bulk properties from available
ingredients. Bounds reveal whether a proposed arrangement performs well,
whether a computation is accurate, and how much improvement remains
possible. Simple layered structures also provide exact reference cases
for checking more complicated methods.

\subsubsection{What progress would mean}

Progress would produce a certified interval for a chosen directional
response and clarify which bounds are geometry-specific or universal.
The result would concern the stated positive, linear conductivities;
nonlinear effects, interfaces and manufacturing constraints would need
separate treatment.

\subsection{Revised core question}

For fixed positive phase conductivities and volume fraction, can
rigorous effective-conductivity bounds and matching constructions be
obtained in a specified microstructure class?

\subsection{Precise scope, target and outputs}

Choose scalar conductivities
a\ensuremath{{}_{1}},a\ensuremath{{}_{2}}>0, periodic unit cell Y, phase
fraction f, and a unit direction e. Define e\ensuremath{{}^{T}}A\_eff
e=min\_\{periodic
u\}\ensuremath{\langle}a|e+\ensuremath{\nabla}u|\ensuremath{{}^{2}}\ensuremath{\rangle}.
Use an admissible potential trial field for an upper bound and a
divergence-free flux variational principle for a lower bound.
Distinguish bounds valid for every arrangement from optimization over a
finite geometry family. Fix isotropy assumptions before comparing
isotropic bounds. Cross-property elasticity/thermal results are separate
extensions with both constitutive laws specified.

\subsection{Detailed critique of the synthesis framing}

The original ``universal bounds'' problem spans conductivity,
permittivity and elasticity, including tensor and frequency effects with
different variational principles. ``Beyond Hashin--Shtrikman'' is not a
measurable novelty requirement until the symmetry class and available
information are fixed. A finite set of candidate microstructures cannot
characterize the unrestricted G-closure. The revision makes one scalar
static problem precise and makes the bound direction and construction
obligation explicit.

\subsection{Validation and milestones}

Baseline: verify arithmetic and harmonic directional bounds with layered
composites. Strong: matching bounds for one declared geometry class or a
certified gap to the unrestricted bound. Check volume fraction and
ellipticity exactly; report discretization and geometry errors
separately.

Use the shared
\href{https://github.com/rick-stevens-ai/Pure-Thought/blob/main/reviews/EVIDENCE-STANDARD.md}{evidence
standard}. Unsuccessful search is unresolved unless a complete exclusion
proof is supplied. These are proposed baseline specifications, not
completed scientific results.

\subsection{Displayed formulation and claim boundary}

\href{https://github.com/rick-stevens-ai/Pure-Thought/blob/main/supplemental/S03.md}{Supplemental
specification}.

\[e\cdot A^*e=\min_{\phi\ \mathrm{periodic}}
\langle a(x)|e+\nabla\phi|^2\rangle,\qquad a(x)\in\{a_1,a_2\},\]

\[\left(\frac f{a_1}+\frac{1-f}{a_2}\right)^{-1}|e|^2
\le e\cdot A^*e\le(fa_1+(1-f)a_2)|e|^2.\]

Assume positive scalar phases, fixed volume fraction \(f\), and a
normalized cell average. Potential trials give upper bounds and the
corresponding flux principle gives lower bounds. A finite microstructure
list does not characterize an unrestricted G-closure.

\begin{center}\rule{0.5\linewidth}{0.5pt}\end{center}

\section{S04. Finite-volume certification of a disordered topological
index}

\textbf{Source:} synthesis challenge 15;
\href{https://github.com/rick-stevens-ai/Pure-Thought/blob/main/NOTES/30-Most-Compelling-Pure-Thought-AI-Challenges.md}{original
synthesis}. This topic is distinct from, or adds a distinct objective
to, the canonical PRD catalog.
\href{https://github.com/rick-stevens-ai/Pure-Thought/blob/main/reviews/SYNTHESIS-CROSSWALK.md}{Crosswalk}

\subsection{Plain-language guide}

\subsubsection{The problem in everyday terms}

A perfectly repeating lattice is convenient for calculating topological
properties, but disorder breaks that repetition. This problem asks when
a calculation on a finite, disordered patch can prove a topological
property of the corresponding infinite system. The key issue is whether
the patch is large enough and the model controlled enough for its answer
to be trustworthy.

\subsubsection{Key terms}

\begin{itemize}
\item
  \textbf{Hamiltonian} --- The operator describing the lattice system's
  energies and couplings.
\item
  \textbf{Disorder} --- Spatial variation in parameters that need not
  repeat periodically.
\item
  \textbf{Locality} --- A restriction on how far or how strongly sites
  can interact.
\item
  \textbf{Fermi energy} --- The reference energy separating occupied and
  unoccupied states in the specified zero-temperature filling model.
\item
  \textbf{Spectral gap} --- An open energy interval with no spectrum of
  the bulk Hamiltonian.
\item
  \textbf{Topological index} --- A discrete quantity that distinguishes
  phases and stays fixed under allowed gap-preserving changes.
\item
  \textbf{Finite volume} --- A bounded patch of a larger or infinite
  lattice.
\item
  \textbf{Spectral localizer} --- A matrix combining Hamiltonian and
  position information to compute an index under stated conditions.
\item
  \textbf{Mobility gap} --- A regime characterized by localization
  rather than necessarily an empty spectral interval; it requires
  different hypotheses.
\end{itemize}

\subsubsection{Why this matters}

Topological behavior is interesting partly because it can survive
certain imperfections. Establishing that robustness requires methods
that work without a perfect crystal and distinguish a finite-sample
observation from a bulk conclusion. A rigorous bridge between the two
makes numerical evidence more informative.

\subsubsection{What progress would mean}

Progress would give checkable size, gap and locality conditions under
which a finite calculation equals the bulk index. It would not
automatically establish robustness to every disorder strength, a
mobility-gap theorem, or the transport response of an experimental
device.

\subsection{Revised core question}

For a local disordered lattice Hamiltonian with a specified bulk
spectral-gap hypothesis, when does a finite-volume index provably equal
its infinite-volume counterpart?

\subsection{Precise scope, target and outputs}

Fix a finite-range two-dimensional tight-binding Hamiltonian on
Z\ensuremath{{}^{2}}, bounded disorder, Fermi energy E\_F and a bulk
spectral gap. Initially use bounded perturbations small enough to
preserve a known clean gap. Choose one index construction, such as a
spectral localizer, and one theorem relating finite and infinite volume.
Define position operators, truncation radius and boundary convention;
compute the finite Hermitian signature or index with interval spectral
bounds. Verify every locality, commutator, radius and scale inequality
in the selected theorem. Mobility-gap cases require a different theorem
and are an extension.

\subsection{Detailed critique of the synthesis framing}

The synthesis lists Bott indices, noncommutative Chern numbers and
localizers as if they shared one universal finite-size error bound. They
have different hypotheses and finite-volume behavior. An integer from a
finite matrix is not automatically the bulk invariant, and one disorder
realization does not establish robustness for all disorder. The revised
task certifies the finite-to-infinite transfer under a stated
gap/locality theorem; it is distinct from the periodic model
classification in PRD 09.

\subsection{Validation and milestones}

Baseline: recover a clean model's known index with a validated
finite-to-bulk theorem. Strong: certify a full bounded disorder class
using gap preservation. Include trivial controls and a gap-closing case
that the checker declines to certify.

Use the shared
\href{https://github.com/rick-stevens-ai/Pure-Thought/blob/main/reviews/EVIDENCE-STANDARD.md}{evidence
standard}. Unsuccessful search is unresolved unless a complete exclusion
proof is supplied. These are proposed baseline specifications, not
completed scientific results.

\subsection{Displayed formulation and claim boundary}

\href{https://github.com/rick-stevens-ai/Pure-Thought/blob/main/supplemental/S04.md}{Supplemental
specification}.

\[\operatorname{dist}(E_F,\operatorname{spec}H_0)\ge g>0,
\qquad\|V\|<g\quad\Longrightarrow\quad
\operatorname{dist}(E_F,\operatorname{spec}(H_0+V))\ge g-\|V\|.\]

This bounded self-adjoint perturbation estimate supplies a spectral-gap
hypothesis. Equality of a chosen finite-volume index with the bulk index
still requires a specific localizer or other index theorem with all
radius, locality, commutator and scale conditions checked. Mobility-gap
claims need a different theorem.

\begin{center}\rule{0.5\linewidth}{0.5pt}\end{center}

\section{S05. Finite-basis density functionals and exact convex
constraints}

\textbf{Source:} synthesis challenge 17;
\href{https://github.com/rick-stevens-ai/Pure-Thought/blob/main/NOTES/30-Most-Compelling-Pure-Thought-AI-Challenges.md}{original
synthesis}. This topic is distinct from, or adds a distinct objective
to, the canonical PRD catalog.
\href{https://github.com/rick-stevens-ai/Pure-Thought/blob/main/reviews/SYNTHESIS-CROSSWALK.md}{Crosswalk}

\subsection{Plain-language guide}

\subsubsection{The problem in everyday terms}

A many-electron state contains vastly more information than the average
number of electrons at each site. Can we still describe the minimum
energy using only those average occupancies? This problem studies that
question in a finite lattice, where the exact relationship can be
defined and bounded without pretending that a convenient approximate
formula is exact.

\subsubsection{Key terms}

\begin{itemize}
\item
  \textbf{Density} --- Here, the expected electron occupancy at each
  lattice site.
\item
  \textbf{Density functional} --- A rule assigning an energy
  contribution to a full density profile.
\item
  \textbf{Finite basis} --- A fixed finite collection of states or
  orbitals used to define the model.
\item
  \textbf{Ensemble} --- A statistical mixture of allowed quantum states.
\item
  \textbf{Constrained search} --- Minimization over all states with a
  prescribed density.
\item
  \textbf{Kinetic energy} --- The energy associated with electron motion
  or hopping in the model.
\item
  \textbf{Interaction energy} --- The energy arising from interactions
  between electrons.
\item
  \textbf{Convexity} --- The property that mixing two density profiles
  cannot raise the functional above the corresponding mixture of
  functional values.
\item
  \textbf{Duality} --- An equivalent optimization viewpoint using
  external potentials to constrain energy values.
\item
  \textbf{Exchange--correlation energy} --- A remainder defined only
  after specifying the noninteracting kinetic and Hartree reference
  contributions.
\end{itemize}

\subsubsection{Why this matters}

Density-based methods reduce a difficult many-electron description to a
much smaller set of variables. Exact finite examples reveal which
constraints an approximation should respect and provide benchmarks for
evaluating proposed energy formulas.

\subsubsection{What progress would mean}

A useful result would bracket the functional on the declared density
domain and verify its convex and dual structure. It would give a
controlled model benchmark, with a separate argument needed before
extending conclusions to continuum chemistry or real materials.

\subsection{Revised core question}

For a fixed finite lattice Hamiltonian, can the ensemble
constrained-search density functional be bracketed and its
convexity/duality verified?

\subsection{Precise scope, target and outputs}

Fix a finite N-particle Hilbert space, kinetic T, interaction W, and
commuting site-density operators \(\widehat n_i\). Define
\(F(n)=\min_{\rho\succeq0,\operatorname{Tr}\rho=1,\operatorname{Tr}(\rho\widehat n_i)=n_i}\operatorname{Tr}\rho(T+W)\)
on ensemble-representable densities. Derive the dual
\(F(n)=\sup_v[E(v)-v\cdot n]\) with
\(E(v)=\min_{\rho\succeq0,\operatorname{Tr}\rho=1}\operatorname{Tr}\rho(T+W+\sum_i v_i\widehat n_i)\),
under the finite-dimensional convex formulation. Use exact states for
primal upper bounds and potentials with certified E(v) lower bounds for
dual lower bounds. If defining an exchange-correlation remainder,
specify T\_s and Hartree conventions separately.

\subsection{Detailed critique of the synthesis framing}

``The cone of exchange-correlation functionals satisfying all exact
constraints'' is not justified: normalization and affine conditions need
not define a cone, and convexity of the ensemble universal functional
does not imply convexity of E\_xc after subtracting other functionals.
Constraint satisfaction alone does not identify the exact physical
functional. The broad continuum task also leaves admissible density
spaces and spin scaling unspecified. The revision first tests exact
convex duality in a finite model, preserving nonempirical functional
development as an extension.

\subsection{Validation and milestones}

Baseline: exact functional values and dual brackets on a two-site model.
Strong: a certified piecewise characterization or a new valid inequality
for a larger finite model. Verify particle-number bounds, convexity
inequalities, and agreement with direct constrained search.

Use the shared
\href{https://github.com/rick-stevens-ai/Pure-Thought/blob/main/reviews/EVIDENCE-STANDARD.md}{evidence
standard}. Unsuccessful search is unresolved unless a complete exclusion
proof is supplied. These are proposed baseline specifications, not
completed scientific results.

\subsection{Displayed formulation and claim boundary}

\href{https://github.com/rick-stevens-ai/Pure-Thought/blob/main/supplemental/S05.md}{Supplemental
specification}.

\[F(n)=\min_{\substack{\rho\succeq0,\ \operatorname{Tr}\rho=1\\
\operatorname{Tr}(\rho\widehat n_i)=n_i}}
\operatorname{Tr}\rho(T+W),\]

\[E(v)=\min_{\rho\succeq0,\ \operatorname{Tr}\rho=1}
\operatorname{Tr}\rho\!\left(T+W+\sum_i v_i\widehat n_i\right),
\qquad F(n)=\sup_v[E(v)-v\cdot n].\]

Work in the fixed finite \(N\)-particle Hilbert space and on
ensemble-representable densities. The supremum need not be attained by a
finite potential at boundary densities. Feasible density matrices give
upper bounds; certified energy lower bounds at chosen potentials give
lower bounds. Convexity of \(F\) does not imply convexity of an
exchange-correlation remainder.

\begin{center}\rule{0.5\linewidth}{0.5pt}\end{center}

\section{S06. Inverse pair-potential design over an explicit competitor
class}

\textbf{Source:} synthesis challenge 20;
\href{https://github.com/rick-stevens-ai/Pure-Thought/blob/main/NOTES/30-Most-Compelling-Pure-Thought-AI-Challenges.md}{original
synthesis}. This topic is distinct from, or adds a distinct objective
to, the canonical PRD catalog.
\href{https://github.com/rick-stevens-ai/Pure-Thought/blob/main/reviews/SYNTHESIS-CROSSWALK.md}{Crosswalk}

\subsection{Plain-language guide}

\subsubsection{The problem in everyday terms}

Instead of asking which arrangement particles prefer under known forces,
turn the question around: which interaction would favor an arrangement
we want? Here, candidate pair interactions are tuned so that a target
repeating structure has lower energy than an explicitly listed set of
alternatives across a range of densities.

\subsubsection{Key terms}

\begin{itemize}
\item
  \textbf{Pair potential} --- A function giving the interaction energy
  of two particles as a function of their separation.
\item
  \textbf{Inverse design} --- Choosing model parameters to produce a
  desired outcome.
\item
  \textbf{Target lattice} --- The repeating particle arrangement the
  design aims to favor.
\item
  \textbf{Competitor class} --- The precise collection of alternative
  arrangements against which the target is compared.
\item
  \textbf{Number density} --- The number of particles per unit volume or
  area.
\item
  \textbf{Energy per particle} --- Total interaction energy normalized
  by the number of particles in the periodic limit.
\item
  \textbf{Lattice sum} --- A sum of interactions over particles in a
  repeating structure.
\item
  \textbf{Tail bound} --- A guaranteed bound on contributions omitted
  when truncating an infinite sum.
\item
  \textbf{Ground state} --- A lowest-energy configuration within a
  stated admissible class.
\end{itemize}

\subsubsection{Why this matters}

Interaction design is a way to explore self-assembly: simple local rules
may favor organized structures. Reliable comparisons expose whether
apparent success comes from a useful mechanism or from overlooking a
competing arrangement or a long-range contribution.

\subsubsection{What progress would mean}

Progress would prove the target's energy advantage, or identify a
defeating competitor, within the declared family and density interval.
Beating that family alone does not establish the ground state among
every possible configuration or show that particles will assemble into
it at finite temperature.

\subsection{Revised core question}

Can a parameterized pair potential make a target periodic structure an
energy minimizer among a precisely specified set of competitors over a
density interval?

\subsection{Precise scope, target and outputs}

Fix dimension, density interval, and
v\_\ensuremath{\theta}(r)=\ensuremath{\Sigma}\_j\ensuremath{\theta}\_j\ensuremath{\phi}\_j(r)
with bounded coefficients and convergent lattice sums. Define
E\_\ensuremath{\theta}(\ensuremath{\Lambda})=1/2
\ensuremath{\Sigma}\_\{x\ensuremath{\in}\ensuremath{\Lambda}\ensuremath{\setminus}\{0\}\}v\_\ensuremath{\theta}(|x|).
Begin with a finite list of normalized lattice competitors and use
interval sums plus tail bounds to certify
E\_target\ensuremath{\le}E\_competitor\ensuremath{-}\ensuremath{\delta}
uniformly on a density interval. Continuous competitors require a
compact parameterization and global coverage. State whether particle
displacements or variable cell size are allowed. General configuration
ground states require additional universal lower bounds beyond lattice
comparisons.

\subsection{Detailed critique of the synthesis framing}

The original jumps from optimizing lattice energies to provable crystals
over both densities and temperatures. Beating selected lattices does not
exclude aperiodic competitors or infinitesimal deformations, and a
zero-temperature energy gap is not a finite-temperature free-energy
theorem. A strictly uniform energy gap against all arbitrarily close
deformations may be impossible under continuity. The rewrite gives a
restricted inverse-design result whose proof scope is honest;
unrestricted crystallization and thermal stability remain separate
research extensions.

\subsection{Validation and milestones}

Baseline: exact or interval-verified ordering for a finite competitor
list. Strong: uniform density/parameter-region certificate. Validate a
direct long-cutoff sum against the certified tail bound; label the
result ``best among these competitors'' unless a universal proof is
supplied.

Use the shared
\href{https://github.com/rick-stevens-ai/Pure-Thought/blob/main/reviews/EVIDENCE-STANDARD.md}{evidence
standard}. Unsuccessful search is unresolved unless a complete exclusion
proof is supplied. These are proposed baseline specifications, not
completed scientific results.

\subsection{Displayed formulation and claim boundary}

\href{https://github.com/rick-stevens-ai/Pure-Thought/blob/main/supplemental/S06.md}{Supplemental
specification}.

\[v_\theta(r)=\sum_j\theta_j\phi_j(r),\qquad
E_\theta(\Lambda)=\frac12\sum_{x\in\Lambda\setminus\{0\}}v_\theta(|x|),\]

\[E_\theta(\Lambda_{\mathrm{target}})\le
E_\theta(\Lambda_{\mathrm{competitor}})-\delta.\]

Require convergent sums and certified tails, uniformly on the declared
density interval and competitor set. The conclusion is best among those
competitors. Unrestricted configurations, infinitesimal deformations and
finite-temperature free energies require separate arguments.

\begin{center}\rule{0.5\linewidth}{0.5pt}\end{center}

\section{S07. Finite-size criteria for parent-Hamiltonian spectral gaps}

\textbf{Source:} synthesis challenge 23;
\href{https://github.com/rick-stevens-ai/Pure-Thought/blob/main/NOTES/30-Most-Compelling-Pure-Thought-AI-Challenges.md}{original
synthesis}. This topic is distinct from, or adds a distinct objective
to, the canonical PRD catalog.
\href{https://github.com/rick-stevens-ai/Pure-Thought/blob/main/reviews/SYNTHESIS-CROSSWALK.md}{Crosswalk}

\subsection{Plain-language guide}

\subsubsection{The problem in everyday terms}

A small quantum chain can have a clear energy separation above its
lowest state while that separation shrinks away as the chain grows. This
problem asks when an exact calculation on a short segment can rule out
that collapse for an entire specified family of longer chains. A theorem
connecting local and global behavior is the essential step.

\subsubsection{Key terms}

\begin{itemize}
\item
  \textbf{Quantum chain} --- A one-dimensional sequence of interacting
  quantum degrees of freedom.
\item
  \textbf{Matrix product state} --- A structured representation of a
  many-body state using products of small matrices.
\item
  \textbf{Parent Hamiltonian} --- A Hamiltonian constructed so that a
  specified state belongs to its ground-state space.
\item
  \textbf{Ground-state space} --- The collection of states attaining the
  minimum energy.
\item
  \textbf{Frustration-free} --- Having ground states that simultaneously
  minimize every local interaction term.
\item
  \textbf{Spectral gap} --- The energy separation above the entire
  ground-state space.
\item
  \textbf{Uniform gap} --- A positive lower bound on that gap that does
  not vanish as chain length increases.
\item
  \textbf{Thermodynamic limit} --- The limit of arbitrarily large
  systems.
\item
  \textbf{Finite-size criterion} --- A theorem translating a suitable
  local calculation into a conclusion for larger systems.
\end{itemize}

\subsubsection{Why this matters}

The gap helps distinguish low-energy behavior and enters rigorous
statements about correlations and stability under suitable conditions.
Proving a uniform gap turns a finite numerical observation into a
statement about a many-body family, which is often the scientifically
relevant object.

\subsubsection{What progress would mean}

Progress would verify all assumptions of a chosen finite-size theorem
and obtain a positive size-independent lower bound. Failing its
threshold would leave the question unresolved rather than prove
gaplessness. The result would apply to the specified family, not solve
the general spectral-gap problem.

\subsection{Revised core question}

For a specified frustration-free one-dimensional parent-Hamiltonian
family, can a validated local spectral calculation imply a uniform
many-body gap?

\subsection{Precise scope, target and outputs}

Fix exact MPS tensors, on-site dimension, interaction range, boundary
convention, and local positive projectors h\_i. Verify frustration
freeness and ground-space assumptions. Select a precise Knabe-type or
martingale theorem appropriate to that model. Compute interval
enclosures for the required finite-segment nonzero gaps and verify the
theorem's threshold to derive a uniform lower bound for the stated
sizes/boundaries. Degenerate ground spaces require identifying the first
strictly positive eigenvalue rather than assuming a unique ground state.

\subsection{Detailed critique of the synthesis framing}

The synthesis's jump from MPS to arbitrary PEPS makes a tractable
finite-size criterion look like a general spectral-gap decision method.
Such general algorithms are obstructed by undecidability results. A
positive finite-chain gap alone says nothing uniform as size grows; one
must meet a theorem's hypotheses and threshold. The revision targets
certificate automation for an explicit restricted family and preserves
unresolved outcomes when the chosen sufficient criterion fails.

\subsection{Validation and milestones}

Baseline: certify a standard parent model with a valid finite-size
theorem. Strong: a uniform parameter interval with an explicit gap lower
bound. Test a deliberately gapless/threshold-failing control without
claiming that criterion failure proves gaplessness.

Use the shared
\href{https://github.com/rick-stevens-ai/Pure-Thought/blob/main/reviews/EVIDENCE-STANDARD.md}{evidence
standard}. Unsuccessful search is unresolved unless a complete exclusion
proof is supplied. These are proposed baseline specifications, not
completed scientific results.

\href{https://arxiv.org/abs/1502.04573}{Cubitt, Pérez-García and Wolf,
Undecidability of the Spectral Gap} motivates the restriction to a
specified certifiable family.

\subsection{Displayed formulation and claim boundary}

\href{https://github.com/rick-stevens-ai/Pure-Thought/blob/main/supplemental/S07.md}{Supplemental
specification}.

\[H_N=\sum_i h_i,\qquad h_i\succeq0,\qquad
\Delta_N=\min(\operatorname{spec}H_N\setminus\{0\}),\]

\[\text{target: }\inf_{N\ \mathrm{in\ the\ stated\ family}}\Delta_N\ge\delta>0.\]

Verify frustration freeness and ground-space assumptions. A specified
finite-size gap theorem must supply the implication from validated local
gaps to the displayed uniform bound. Positive gaps on finitely many
chains are insufficient; failure of a sufficient criterion is not a
gaplessness proof.

\begin{center}\rule{0.5\linewidth}{0.5pt}\end{center}

\section{S08. Gapped boundaries from exact anyon data}

\textbf{Source:} synthesis challenge 24;
\href{https://github.com/rick-stevens-ai/Pure-Thought/blob/main/NOTES/30-Most-Compelling-Pure-Thought-AI-Challenges.md}{original
synthesis}. This topic is distinct from, or adds a distinct objective
to, the canonical PRD catalog.
\href{https://github.com/rick-stevens-ai/Pure-Thought/blob/main/reviews/SYNTHESIS-CROSSWALK.md}{Crosswalk}

\subsection{Plain-language guide}

\subsubsection{The problem in everyday terms}

In some two-dimensional quantum phases, collective excitations behave
neither like ordinary bosons nor fermions. Their possible combinations
and exchanges restrict how the phase can end at an edge. This problem
asks which proposed ways of absorbing these excitations at a boundary
are mathematically consistent and can describe an edge with a nonzero
energy gap.

\subsubsection{Key terms}

\begin{itemize}
\item
  \textbf{Topological order} --- A type of quantum organization
  characterized by nonlocal structure and, here, anyon excitation data.
\item
  \textbf{Anyon} --- A two-dimensional excitation with exchange or
  braiding behavior beyond ordinary bosonic and fermionic possibilities.
\item
  \textbf{Fusion} --- The possible outcomes when excitations are
  combined.
\item
  \textbf{Braiding} --- The transformation associated with moving
  excitations around one another.
\item
  \textbf{Topological twist} --- The phase associated with a full
  rotation of an anyon.
\item
  \textbf{F and R data} --- Numbers or matrices specifying how fusion
  regrouping and exchanges act in chosen conventions.
\item
  \textbf{Condensate} --- A collection of excitations treated as
  absorbable into the boundary vacuum, together with consistent
  combination rules.
\item
  \textbf{Lagrangian algebra} --- The mathematical structure expressing
  a sufficiently complete compatible condensate for the boundary
  considered here.
\item
  \textbf{Gapped boundary} --- An edge whose excitations retain a
  positive energy threshold in the appropriate large-system limit.
\end{itemize}

\subsubsection{Why this matters}

Boundaries determine how a topological phase interfaces with its
surroundings and which excitations can disappear or remain
distinguishable there. Understanding these options supports the theory
of topological phases and the organization of quantum information in
such systems.

\subsubsection{What progress would mean}

Progress would verify a boundary candidate using complete consistency
data, or reject it for a precise algebraic reason. Matching a list of
fusion outcomes alone is insufficient. Turning the algebraic result into
a particular microscopic boundary Hamiltonian is an additional
construction.

\subsection{Revised core question}

For one fully specified two-dimensional topological order, can candidate
condensates and their gapped boundaries be verified from exact
categorical data?

\subsection{Precise scope, target and outputs}

Start with the toric-code modular category, explicit simple objects,
fusion, twists, F and R symbols and conventions. Enumerate candidate
condensable algebra objects within a declared multiplicity bound. Check
multiplication/unit, associativity, commutativity and separability
conditions, and the Lagrangian dimension condition for a boundary to
vacuum. For the abelian toric-code benchmark recover electric and
magnetic condensates. If constructing a lattice boundary, supply its
Hamiltonian and check commuting-projector and gap claims separately.

\subsection{Detailed critique of the synthesis framing}

Fusion rules and modular S,T tables alone do not in general specify or
prove all categorical coherence data. The synthesis also mixes 2D anyons
with 3D topological orders, which require different structures. A listed
condensate is not automatically a microscopic boundary realization. The
rewrite retains an exact categorical task, bounds the enumeration and
separates algebraic classification from lattice realization and
anomaly/inflow interpretation.

\subsection{Validation and milestones}

Baseline: verify the toric-code coherence data and two standard vacuum
boundaries. Strong: a second fully specified category or a completeness
proof within a bounded class. Invalid F/R symbols and a nonbosonic
condensate are rejection controls.

Use the shared
\href{https://github.com/rick-stevens-ai/Pure-Thought/blob/main/reviews/EVIDENCE-STANDARD.md}{evidence
standard}. Unsuccessful search is unresolved unless a complete exclusion
proof is supplied. These are proposed baseline specifications, not
completed scientific results.

\href{https://arxiv.org/abs/2504.19512}{Li, Yang and Dong, Gapped
Boundaries of Kitaev's Quantum Double Models} provides a primary example
of linking Lagrangian algebras to explicit lattice boundary terms.

\subsection{Displayed formulation and claim boundary}

\href{https://github.com/rick-stevens-ai/Pure-Thought/blob/main/supplemental/S08.md}{Supplemental
specification}.

\[A=\bigoplus_a n_a a,\qquad n_a\in\mathbb Z_{\ge0},\qquad
(\dim A)^2=\operatorname{Dim}\mathcal C.\]

For a vacuum boundary this Lagrangian dimension condition supplements,
rather than replaces, the unit, multiplication, associativity,
commutativity and separability checks using exact categorical coherence
data. In the toric-code category the reference condensates are
\(1\oplus e\) and \(1\oplus m\). A microscopic boundary Hamiltonian
needs its own gap proof.

\begin{center}\rule{0.5\linewidth}{0.5pt}\end{center}

\section{S09. Stoquasticity under a bounded class of basis changes}

\textbf{Source:} synthesis challenge 25;
\href{https://github.com/rick-stevens-ai/Pure-Thought/blob/main/NOTES/30-Most-Compelling-Pure-Thought-AI-Challenges.md}{original
synthesis}. This topic is distinct from, or adds a distinct objective
to, the canonical PRD catalog.
\href{https://github.com/rick-stevens-ai/Pure-Thought/blob/main/reviews/SYNTHESIS-CROSSWALK.md}{Crosswalk}

\subsection{Plain-language guide}

\subsubsection{The problem in everyday terms}

Quantum calculations can involve large positive and negative
contributions that almost cancel, making some sampling methods
inefficient. Sometimes a change in the mathematical basis removes
troublesome signs. This problem asks whether that is possible within a
finite, explicitly allowed set of local basis changes, so that both
success and failure can be checked completely.

\subsubsection{Key terms}

\begin{itemize}
\item
  \textbf{Basis} --- The reference states used to write a quantum
  operator as a matrix.
\item
  \textbf{Hamiltonian} --- The matrix describing the system's energy and
  couplings.
\item
  \textbf{Off-diagonal entry} --- A matrix element linking two different
  basis states.
\item
  \textbf{Stoquastic} --- Having real, nonpositive off-diagonal
  Hamiltonian entries in the specified basis.
\item
  \textbf{Sign problem} --- A sampling difficulty caused by signed or
  complex weights and cancellation in a chosen representation.
\item
  \textbf{Pauli operator} --- One of the standard single-qubit matrices
  used to express spin interactions.
\item
  \textbf{Clifford transformation} --- A basis transformation that maps
  Pauli operators to Pauli operators under conjugation.
\item
  \textbf{Product transformation} --- A change of basis applied
  independently to the individual qubits.
\item
  \textbf{Satisfiability certificate} --- A checkable assignment meeting
  encoded constraints, or an appropriate proof that none exists.
\end{itemize}

\subsubsection{Why this matters}

Knowing when signs can be removed helps separate difficulties caused by
representation from those that survive a proposed remedy. A bounded
search also creates small, exact benchmarks for algorithms intended to
find useful basis changes.

\subsubsection{What progress would mean}

A successful result would supply and verify a transformation, or exclude
every transformation in the allowed finite class. The complete
Hamiltonian must be checked because contributions can cancel.
Stoquasticity alone does not guarantee fast sampling, and failure for
product Clifford changes does not exclude all possible bases.

\subsection{Revised core question}

For a finite rational Pauli Hamiltonian, is there an allowed product
Clifford transformation making all off-diagonal entries real and
nonpositive?

\subsection{Precise scope, target and outputs}

Fix n qubits, rational coefficients, and allowed
U=\ensuremath{\otimes}\_iU\_i with U\_i from the finite single-qubit
Clifford group modulo phase. Define stoquasticity in the computational
basis for the full transformed H, not automatically term by term.
Exhaustively enumerate or use an exactly equivalent constraint encoding.
Export U and an exact matrix/sign check if successful; otherwise export
complete enumeration or a checked UNSAT proof. Continuous local
unitaries require a semialgebraic formulation and are a distinct
extension.

\subsection{Detailed critique of the synthesis framing}

The original omits which local basis changes are allowed, so
``unavoidable sign problem'' has no definite quantifier. A SAT search
over finitely many transformations cannot exclude all continuous local
rotations. Requiring every term to be stoquastic can also be stronger
than requiring their sum to be stoquastic because entries can cancel.
Finally, stoquasticity does not guarantee efficient sampling or mixing.
The rewrite makes a finite, falsifiable basis-cure problem and limits
every no-go conclusion to its transformation class.

\subsection{Validation and milestones}

Baseline: exhaustive small-n decisions with explicit transformations or
coverage. Strong: verified SAT encoding checked against brute force.
Extend to a continuous rotation class only with exact equations,
inequalities, and global proof machinery.

Use the shared
\href{https://github.com/rick-stevens-ai/Pure-Thought/blob/main/reviews/EVIDENCE-STANDARD.md}{evidence
standard}. Unsuccessful search is unresolved unless a complete exclusion
proof is supplied. These are proposed baseline specifications, not
completed scientific results.

\href{https://pmc.ncbi.nlm.nih.gov/articles/PMC6450938/}{Marvian, Lidar
and Hen, On the computational complexity of curing non-stoquastic
Hamiltonians} supports the need to specify the allowed basis class.

\subsection{Displayed formulation and claim boundary}

\href{https://github.com/rick-stevens-ai/Pure-Thought/blob/main/supplemental/S09.md}{Supplemental
specification}.

\[U=\bigotimes_{i=1}^nU_i,\quad U_i\in\mathrm{Clifford}_1/\mathrm{phase},
\qquad H'=UHU^\dagger,\]

\[\operatorname{Im}H'_{xy}=0,\qquad\operatorname{Re}H'_{xy}\le0\quad(x\ne y).\]

Apply these conditions to the full transformed matrix, allowing
cancellations between terms. Exhaustive coverage or an equivalent
checked UNSAT proof excludes only this finite transformation class.
Stoquasticity alone does not guarantee efficient mixing or sampling.

\begin{center}\rule{0.5\linewidth}{0.5pt}\end{center}

\section{S10. Certified periodic orbits of a specified three-body model}

\textbf{Source:} synthesis challenge 27;
\href{https://github.com/rick-stevens-ai/Pure-Thought/blob/main/NOTES/30-Most-Compelling-Pure-Thought-AI-Challenges.md}{original
synthesis}. This topic is distinct from, or adds a distinct objective
to, the canonical PRD catalog.
\href{https://github.com/rick-stevens-ai/Pure-Thought/blob/main/reviews/SYNTHESIS-CROSSWALK.md}{Crosswalk}

\subsection{Plain-language guide}

\subsubsection{The problem in everyday terms}

A simulated three-body trajectory may nearly return to where it started,
but a small numerical error can disguise whether an exactly repeating
orbit exists. This problem converts a promising numerical trajectory
into a proof that a nearby periodic solution exists and stays away from
collisions, then examines how small disturbances behave over one cycle.

\subsubsection{Key terms}

\begin{itemize}
\item
  \textbf{Three-body problem} --- The motion of three mutually
  gravitating bodies in the specified Newtonian model.
\item
  \textbf{Restricted model} --- A version in which one body has
  negligible mass and does not affect the other two.
\item
  \textbf{Periodic orbit} --- A trajectory that returns to the same
  state after a fixed positive time.
\item
  \textbf{Shooting method} --- Adjusting an initial state until its
  evolved state satisfies the return conditions.
\item
  \textbf{Phase condition} --- A rule fixing the arbitrary starting
  point along a periodic trajectory.
\item
  \textbf{Validated integration} --- Numerical evolution with guaranteed
  enclosures for errors.
\item
  \textbf{Collision exclusion} --- A proved positive lower bound on
  relevant body separations along the orbit.
\item
  \textbf{Floquet multiplier} --- An eigenvalue describing linear
  perturbation growth or rotation after one orbit.
\item
  \textbf{Neutral direction} --- A direction with constrained or trivial
  linear behavior arising from symmetries or conservation laws.
\end{itemize}

\subsubsection{Why this matters}

Periodic orbits are organizing reference motions in gravitational
dynamics. Certifying their existence gives subsequent studies of nearby
motion a dependable starting point and separates exact solutions from
visual patterns in a simulation.

\subsubsection{What progress would mean}

Progress would enclose an orbit and its period, exclude collisions, and
classify the reduced linear stability when the bounds permit. Linear
stability would not alone prove nonlinear or indefinite stability, and a
periodic-orbit proof would not establish connecting trajectories between
different orbits.

\subsection{Revised core question}

Can a numerically proposed periodic orbit be enclosed by a validated
shooting proof and classified by its reduced Floquet spectrum?

\subsection{Precise scope, target and outputs}

Choose one nondimensional Newtonian three-body model with positive
rational masses, or a planar circular restricted model with fixed mass
ratio. Declare collision exclusions and conserved quantities. Define a
Poincaré section or phase condition to remove time translation, and
reduce continuous symmetries. Solve the resulting shooting residual
using a validated flow integrator and interval Newton/Krawczyk bounds
for initial conditions and period. Integrate variational equations with
enclosures for the monodromy matrix; remove the expected neutral
directions before interpreting multipliers.

\subsection{Detailed critique of the synthesis framing}

The original bundles orbit existence, complete family discovery, linear
stability, invariant manifolds and heteroclinic connections into one
certificate. A small shooting residual is not existence, and
stable-looking Floquet multipliers are not nonlinear stability.
Near-return trajectories do not establish a connection between invariant
manifolds. The revised first target is a single orbit with a phase-fixed
existence proof; continuation, manifold enclosure and connection proofs
are explicit later stages. This also preserves the figure-eight idea
incorrectly folded into PRD 27's central configurations.

\subsection{Validation and milestones}

Baseline: certify one supplied orbit and its period. Strong: validated
continuation over a parameter interval and a reduced Floquet
classification. A manifold connection requires additional
covering/intersection proof, not a plot of approaching trajectories.

Use the shared
\href{https://github.com/rick-stevens-ai/Pure-Thought/blob/main/reviews/EVIDENCE-STANDARD.md}{evidence
standard}. Unsuccessful search is unresolved unless a complete exclusion
proof is supplied. These are proposed baseline specifications, not
completed scientific results.

\subsection{Displayed formulation and claim boundary}

\href{https://github.com/rick-stevens-ai/Pure-Thought/blob/main/supplemental/S10.md}{Supplemental
specification}.

\[F(z_0,T)=\varphi_T(z_0)-z_0=0,\qquad
M=D\varphi_T(z_0).\]

The shooting residual must first be augmented and reduced using a phase
condition, conserved quantities and continuous symmetries so the
validated root problem has the correct dimension. Interval flow and root
bounds establish existence; the reduced Floquet multipliers classify
linear behavior. Neither a near return nor stable-looking multipliers
prove nonlinear stability or a manifold connection.

\begin{center}\rule{0.5\linewidth}{0.5pt}\end{center}

\section{S11. Minimal reaction implementations of a specified
computational behavior}

\textbf{Source:} synthesis challenge 29;
\href{https://github.com/rick-stevens-ai/Pure-Thought/blob/main/NOTES/30-Most-Compelling-Pure-Thought-AI-Challenges.md}{original
synthesis}. This topic is distinct from, or adds a distinct objective
to, the canonical PRD catalog.
\href{https://github.com/rick-stevens-ai/Pure-Thought/blob/main/reviews/SYNTHESIS-CROSSWALK.md}{Crosswalk}

\subsection{Plain-language guide}

\subsubsection{The problem in everyday terms}

A reaction network can encode inputs as molecular amounts and produce an
output through its reactions. But calling a network a computer is
ambiguous unless we specify what counts as the input, the answer and an
acceptable error. This problem asks for the smallest network that
implements one precisely defined behavior within a bounded set of
permitted reactions.

\subsubsection{Key terms}

\begin{itemize}
\item
  \textbf{Reaction network} --- A set of chemical species and rules for
  converting reactants into products.
\item
  \textbf{Reaction grammar} --- The restrictions defining which species
  and reaction forms are allowed.
\item
  \textbf{Encoding} --- The mapping from abstract inputs to initial
  molecular states or concentrations.
\item
  \textbf{Boolean function} --- A rule mapping a finite collection of
  binary inputs to a binary output.
\item
  \textbf{Deterministic semantics} --- A model in which concentrations
  follow fixed evolution equations.
\item
  \textbf{Stochastic semantics} --- A model in which individual reaction
  events occur randomly.
\item
  \textbf{Readout} --- The rule and time for translating the chemical
  state into an output.
\item
  \textbf{Error bound} --- A limit on deviation from the required
  output, or on the probability of an incorrect answer.
\item
  \textbf{Cost model} --- The explicit quantity being minimized, such as
  reaction or species count.
\item
  \textbf{Universality} --- The ability to implement arbitrary
  computations under a stated encoding and resource scheme.
\end{itemize}

\subsubsection{Why this matters}

Chemical computation explores how local reactions can perform
information processing, including ideas relevant to molecular sensing
and control. Precise semantics make proposed constructions comparable
and prevent hidden resources or flexible readout choices from creating
misleading claims of simplicity.

\subsubsection{What progress would mean}

Progress would provide a verified implementation and, if claimed, prove
minimality within the bounded grammar and cost model. Implementing a
finite truth table would establish that task only; computational
universality requires a construction for arbitrary programs and an
explicit account of resources.

\subsection{Revised core question}

Within a bounded reaction grammar, what is the smallest network
implementing a precisely defined input/output behavior with a proved
error bound?

\subsection{Precise scope, target and outputs}

Choose deterministic mass-action or stochastic reaction semantics,
species/reaction bounds, stoichiometric arity, input encoding, output
decoding and cost (species count first, then reaction count, for
example). Begin with a finite Boolean input/output function and an
explicit settling-time/error requirement. Enumerate networks modulo
species relabeling; prove correctness over the complete allowed input
set and exclude every lower-cost network within the grammar. For a
universal-computation extension, supply a uniform compiler and a
simulation theorem covering arbitrary program lengths, resource scaling
and precision.

\subsection{Detailed critique of the synthesis framing}

The synthesis joins minimal autocatalysis, self-replication and
universal computation, but none implies the others. ``Smallest
chemistry'' depends on kinetics, encoding, concentration precision and
which resources count as free. A fixed finite truth table does not prove
Turing universality, while an analog model with unbounded precision may
hide computation in its inputs. The rewrite separates a bounded minimal
implementation result from a universal simulation theorem and from the
RAF structure in PRD 29.

\subsection{Validation and milestones}

Baseline: one certified reaction implementation with explicit semantics.
Strong: minimality within a finite grammar and input domain. Extension:
a compiler/simulation theorem; avoid calling the baseline universal or
biologically self-replicating.

Use the shared
\href{https://github.com/rick-stevens-ai/Pure-Thought/blob/main/reviews/EVIDENCE-STANDARD.md}{evidence
standard}. Unsuccessful search is unresolved unless a complete exclusion
proof is supplied. These are proposed baseline specifications, not
completed scientific results.

\subsection{Displayed formulation and claim boundary}

\href{https://github.com/rick-stevens-ai/Pure-Thought/blob/main/supplemental/S11.md}{Supplemental
specification}.

\[\min_{\mathcal R\in\mathfrak G}\operatorname{cost}(\mathcal R)
\quad\text{subject to}\quad
\forall x\in\mathcal X:\ \operatorname{error}(\mathcal R,x;T)\le\varepsilon.\]

Declare the finite reaction grammar \(\mathfrak G\), semantics,
encoding, decoding, time requirement, input set and cost ordering.
Correctness of one network gives an upper cost bound; excluding every
lower-cost network gives minimality within that grammar. A finite truth
table does not prove universal computation or self-replication.

\begin{center}\rule{0.5\linewidth}{0.5pt}\end{center}

\section{S12. Capacity of a specified finite noisy genotype--phenotype
channel}

\textbf{Source:} synthesis challenge 30;
\href{https://github.com/rick-stevens-ai/Pure-Thought/blob/main/NOTES/30-Most-Compelling-Pure-Thought-AI-Challenges.md}{original
synthesis}. This topic is distinct from, or adds a distinct objective
to, the canonical PRD catalog.
\href{https://github.com/rick-stevens-ai/Pure-Thought/blob/main/reviews/SYNTHESIS-CROSSWALK.md}{Crosswalk}

\subsection{Plain-language guide}

\subsubsection{The problem in everyday terms}

Suppose the same genetic sequence can produce different modeled outcomes
because of noise. How reliably can an observer infer the chosen sequence
from the outcome? This problem treats the relationship as an explicitly
given communication channel and asks for certified bounds on the
greatest information it can transmit under any stated input-cost limit.

\subsubsection{Key terms}

\begin{itemize}
\item
  \textbf{Genotype} --- An input genetic sequence or label in the finite
  model.
\item
  \textbf{Phenotype} --- An output trait or state as defined by the
  model.
\item
  \textbf{Channel law} --- The conditional probability of each output
  given each input.
\item
  \textbf{Noise} --- Variation in the output for a fixed input, as
  encoded by that channel law.
\item
  \textbf{Input distribution} --- The probabilities assigned to
  selecting the possible genotypes.
\item
  \textbf{Mutual information} --- A measure of how much observing the
  output reduces uncertainty about the input.
\item
  \textbf{Shannon capacity} --- The largest mutual information over the
  allowed input distributions for this channel.
\item
  \textbf{Bit} --- The information unit obtained when logarithms use
  base two.
\item
  \textbf{Cost constraint} --- A bound on average input cost, such as a
  prescribed sequence-dependent resource measure.
\item
  \textbf{Relative entropy} --- A measure of discrepancy between
  probability distributions, also called Kullback--Leibler divergence.
\end{itemize}

\subsubsection{Why this matters}

An explicit channel separates the number of possible outputs from their
reliability and distinguishability. Capacity bounds give a rigorous way
to compare modeled noise mechanisms or resource constraints, and to
check whether a numerical optimization has found the best input
distribution.

\subsubsection{What progress would mean}

Progress would bracket capacity with a feasible input distribution and a
valid upper certificate. For a deterministic unconstrained channel, the
number of reachable outputs supplies a simple benchmark. The answer
describes the specified probabilities and costs; it is not an intrinsic
information value of a genome or a measure of evolutionary fitness.

\subsection{Revised core question}

For a finite conditional phenotype law P(y\textbar g) and optional
genotype cost constraint, what certified lower and upper bounds can be
placed on Shannon capacity?

\subsection{Precise scope, target and outputs}

Supply finite genotype and phenotype alphabets, a stochastic matrix
W(y\textbar g), and optional costs c(g) with E[c]\ensuremath{\le}C.
Define capacity max\_p I\_p(G;Y), with logarithms base two. A chosen
input distribution gives a lower bound. For the unconstrained channel,
any output distribution q with adequate support gives upper bound max\_g
D(W(.\textbar g)\textbar\textbar q); constrained dual bounds must
include the cost multiplier. Use interval logarithms and validated
probabilities. Connect to PRD 30 by constructing W from a fully declared
folding/noise rule.

\subsection{Detailed critique of the synthesis framing}

Mutual information is a property of a joint distribution, not a map
alone. The original omits the input distribution, noise, phenotype
alphabet and optimization constraints. For a deterministic finite map
with unrestricted input distribution, capacity is simply
log\ensuremath{{}_{2}} of the number of reachable phenotypes, so a
purported hard capacity problem may collapse. Fano/Le Cam bounds answer
specific inference questions and are not automatic capacity
certificates. The rewrite specifies a nontrivial noisy channel,
separates achievability from converse bounds, and keeps biophysical
interpretation conditional on the chosen channel.

\subsection{Validation and milestones}

Baseline: binary symmetric-channel capacity and deterministic-map
capacity reproduced. Strong: a certified bracket for a finite GP channel
and a cost-constrained extension with a valid dual. Check rows sum to
one, handle zero probabilities explicitly, and report unresolved
numerical gaps.

Use the shared
\href{https://github.com/rick-stevens-ai/Pure-Thought/blob/main/reviews/EVIDENCE-STANDARD.md}{evidence
standard}. Unsuccessful search is unresolved unless a complete exclusion
proof is supplied. These are proposed baseline specifications, not
completed scientific results.

\subsection{Displayed formulation and claim boundary}

\href{https://github.com/rick-stevens-ai/Pure-Thought/blob/main/supplemental/S12.md}{Supplemental
specification}.

\[C=\max_p I_p(G;Y),\qquad
I_p(G;Y)\le C\le\max_g D(W(\cdot|g)\|q).\]

For a cost limit \(c_*\), any \(\lambda\ge0\) and adequately supported
\(q\) give

\[C(c_*)\le\lambda c_*+
\max_g\{D(W(\cdot|g)\|q)-\lambda c(g)\}.\]

Use base-two logarithms and verified stochastic rows. A feasible input
distribution supplies a lower bound, and interval logarithms certify
numerical values. For a deterministic unrestricted finite map, capacity
is \(\log_2|\operatorname{im}\Phi|\); this differs from neutral-network
connectivity.

\section{Evidence and execution standard}

This standard applies to the revised specifications. It is a proposed
contract for future research runs; this editorial revision does not
supply the scientific certificates themselves.

\subsection{What ``pure thought'' means here}

Use explicitly stated mathematical models and deductions without fitting
conclusions to experimental outcomes. Synthetic examples, exact model
inputs, analytic references, numerical search, and established software
are allowed. Physical constants, electronic-structure approximations,
constitutive laws, and empirical parameter sets must be identified when
used. A model-derived result is conditional on that model.

The original ban on legacy software conflicts with its proposed reliance
on NumPy, SciPy, CVXPY, nauty, and other libraries. Reimplementing every
optimizer and algebra system is not necessary for scientific
independence. Prefer a small independently checkable certificate and a
documented trusted base. A fresh implementation can be a project
objective, but must not be confused with a mathematical proof or novelty
claim.

\subsection{Four evidence levels}

{\def\LTcaptype{none} % do not increment counter
\begin{longtable}[]{@{}
  >{\raggedright\arraybackslash}p{(\linewidth - 4\tabcolsep) * \real{0.3333}}
  >{\raggedright\arraybackslash}p{(\linewidth - 4\tabcolsep) * \real{0.3333}}
  >{\raggedright\arraybackslash}p{(\linewidth - 4\tabcolsep) * \real{0.3333}}@{}}
\toprule\noalign{}
\begin{minipage}[b]{\linewidth}\raggedright
Label
\end{minipage} & \begin{minipage}[b]{\linewidth}\raggedright
What it establishes
\end{minipage} & \begin{minipage}[b]{\linewidth}\raggedright
What it does not establish
\end{minipage} \\
\midrule\noalign{}
\endhead
\bottomrule\noalign{}
\endlastfoot
Numerical evidence & A reproducible computation at stated precision,
sample size, and discretization & Exact truth, continuum coverage, or a
universal theorem \\
Finite-model certificate & An exact/validated statement about the
specified finite model or finite relaxation & An omitted tail,
thermodynamic limit, or physical realization \\
Conditional mathematical theorem & A proof over the stated domain with
every hypothesis and error controlled & Truth outside those assumptions
or a demonstrated physical system \\
Formal proof & A checked theorem in a named proof assistant, with
imported axioms/dependencies listed & Correct physical modeling merely
by virtue of formalization \\
\end{longtable}
}

Formalization is an optional extra deliverable unless the specific
project makes it mandatory. High precision, JSON output, solver status
and small residuals do not by themselves confer any of these proof
levels.

\subsection{Claim directions}

For minimization over a larger feasible set, the optimum is a lower
bound on the original minimum. A physical feasible point supplies an
upper bound. For maximization the directions reverse. Track feasibility
of the primal/dual witnesses separately from their objective values.

Infeasibility of an inner ansatz excludes only that ansatz.
Infeasibility of an outer relaxation can exclude the original problem if
the map into the relaxation is proved and all approximations preserve
necessity. Passing necessary tests means ``not excluded''; existence
needs a constructive sufficient argument.

A smallest observed example is not a minimality theorem. A complete
catalog must define the universe, equivalence relation, boundaries, and
proof that pruning misses nothing. A finite-size numerical trend does
not prove an asymptotic limit.

\subsection{Minimum run manifest}

Before calculation, record the problem ID and revision hash, exact
inputs and conventions, quantified domain, equivalence relation, target
observable, accepted evidence level, tolerances, and resource limits.
Record dependencies and precision/rounding mode. Where randomization is
used, preserve seeds, trial counts and raw sufficient statistics.
Resource limits are planning decisions, not evidence that a mathematical
solution cannot exist.

An output bundle should include:

\begin{itemize}
\tightlist
\item
  The statement being claimed, including all assumptions and
  quantifiers.
\item
  Construction or proof object, plus an independently runnable verifier
  and its trust assumptions.
\item
  Numerical/error ledger separating roundoff, discretization,
  truncation, optimization, sampling and modeling error where relevant.
\item
  Positive reference cases and negative controls that the verifier must
  reject.
\item
  Coverage: included domain, omitted domain, unresolved regions and
  termination reason.
\end{itemize}

These are requirements for future implementations, not a promise of
polynomial-size certificates for every problem.

\subsection{Checking specific certificate types}

\textbf{LP/SDP/SOS:} verify equality constraints, objective direction,
and cone feasibility using exact arithmetic or outward-rounded
intervals. Rationalizing a numerical solution can break PSD constraints.
Correct residuals with justified norm bounds; do not silently project
and reuse the old objective. Distinguish a proof of a polynomial
approximation from a proof for the original function.

\textbf{Interval numerics:} state function enclosures, domains and
outward rounding. A small root box needs existence and, if claimed,
uniqueness. Cover the complement for a global enumeration. Include
singularities and tails rather than discarding them numerically.

\textbf{SAT/SMT:} a solver's UNSAT label is not a certificate. Preserve
a proof supported by a checker and the proof that the encoding, symmetry
reduction and blocked solutions match the mathematical universe. DRAT
applies to an appropriate Boolean CNF proof, not automatically to every
real-arithmetic or SDP claim.

\textbf{Statistics:} specify the estimand, dependence assumptions,
confidence procedure and bias. Zero observed failures does not prove
zero failure probability. MCMC error bars based on estimated
autocorrelation are not generally rigorous mixing guarantees.

\subsection{Milestones and novelty}

Baseline means a working, independently validated result on a concrete
reference problem. Strong means a broader or sharper result with its
claimed evidence level. Research extension means a possible novel result
whose feasibility and significance need review. A negative theorem can
be valuable; an unsuccessful search can also be useful if its scope is
documented, but the two are not interchangeable.

Do not promise ``publication-quality'' in a fixed number of months.
Before claiming novelty, conduct a focused current literature comparison
against the exact theorem/model/metric. The sources in this revision
check key definitions and limitations; they are not an exhaustive
novelty audit across all 30 fields.

\end{document}
